
\documentclass[ucs,9pt, mathserif]{beamer}

\usepackage[utf8x]{inputenc}   
\usepackage[german, ngerman]{babel}
\usepackage{amsmath}
\usepackage{amsfonts}   
\usepackage{amssymb}
\usepackage{amsthm}
\usepackage{nicefrac}
\usepackage{siunitx}
\usepackage{xstring}
\usepackage{nicefrac}
\usepackage{ulem}
\usepackage{tabu}
\usepackage[eulergreek]{sansmath}
\usepackage{multicol}


\usepackage{listings}
\usepackage[ruled,vlined,german,linesnumbered, titlenotnumbered]{algorithm2e} % Algorithmenpacket

\usepackage{graphicx}
\usepackage{float}
\usepackage{wrapfig}
\usepackage{tikz}
\usetikzlibrary{positioning, arrows, trees}
\usepackage{fix-cm}
\usepackage{pgfplots}
\usepackage{multimedia}



\pgfdeclareimage[height=0.7cm]{university-logo}{logo}
\logo{\pgfuseimage{ohmlogo}}

% large version for upper right corner of title page
%\pgfdeclareimage[height=1.085cm]{big-university-logo}{FULogo_RGB}
\newcommand{\titleimage}[1]{\pgfdeclareimage[height=4cm]{title-image}{#1}}


\titlegraphic{\pgfuseimage{title-image}}
%%% end logo

% NOTE: 1cm = 0.393 in = 28.346 pt;    1 pt = 1/72 in = 0.0352 cm
\setbeamersize{text margin right=3.5mm, text margin left=3.5mm}  % text margin

% colors to be used
\definecolor{text-grey}{rgb}{0.45, 0.45, 0.45} % grey text on white background
\definecolor{bg-grey}{rgb}{0.66, 0.65, 0.60} % grey background (for white text)

% Farben
%\definecolor{grey}{grey}{0.5} % die Farbe grau wird definiert
%\definecolor{light-grey}{grey}{0.95}

\definecolor{ohmyellow}{cmyk}{0.01, 0.08, 0.92, 0.01}
\definecolor{ohmblue}{cmyk}{0.94, 0.63, 0.03, 0.05}
\definecolor{ohmred}{cmyk}{0.02, 0.94, 0.90, 0.06}
\definecolor{ohmorange}{cmyk}{0.01, 0.57, 0.91, 0.02}
\definecolor{ohmcoolgrey}{cmyk}{0.55, 0.31, 0.25, 0.02}
\definecolor{ohmgreen}{cmyk}{0.83, 0.03, 0.93, 0.06}

%%% colors
\usecolortheme{lily}
 
\setbeamercolor{frametitle}{fg=ohmblue}
\setbeamercolor{title}{fg= ohmblue}

\setbeamercolor*{normal text}{fg=black,bg=white}
\setbeamercolor*{alerted text}{fg=red}
\setbeamercolor*{example text}{fg=green}
\setbeamercolor*{structure}{fg=ohmblue}

\setbeamercolor*{block title}{fg=white,bg=black!50}
\setbeamercolor*{block title alerted}{fg=white,bg=black!50}
\setbeamercolor*{block title example}{fg=white,bg=black!50}

\setbeamercolor*{block body}{bg=black!10}
\setbeamercolor*{block body alerted}{bg=black!10}
\setbeamercolor*{block body example}{bg=black!10}

\setbeamercolor{bibliography entry author}{fg=ohmblue}
\setbeamercolor{bibliography entry journal}{fg=text-grey}

\setbeamercolor{item}{fg=ohmblue}
\setbeamercolor{navigation symbols}{fg=text-grey,bg=bg-grey}
%%% end colors



\setbeamersize{sidebar width left=0cm, sidebar width right=0mm}
\setbeamertemplate{sidebar right}{}
\setbeamertemplate{sidebar left}{}
%    XOR
% \useoutertheme{sidebar}

\setbeamertemplate{frametitle}

\setbeamertemplate{frametitle}{%
    \vskip-30pt \color{black}\large\bfseries%
    \begin{minipage}[b][23pt]{90.5mm}%
    \flushleft\insertframetitle%
    \end{minipage}%
}


\setbeamertemplate{title page}{
    \hfill%\begin{minipage}[t]{1cm} 
        \includegraphics[scale=0.5]{src/logo.png}
    %\end{minipage} 
    % necessary line
   %\begin{minipage}[b]{4.3cm} 
   %     \small{\textsf{\textbf{Technische Hochschule Nürnberg} \\
   %     Georg-Simon-Ohm\\[0.16cm]
   %     Fakultät Informatik }}
   % \end{minipage} \\
   % \vskip0pt%\hskip3pt
%    \begin{minipage}{15.6cm}
%        \hspace{-1mm}\inserttitlegraphic
%    \end{minipage}

    % set the title and the author
    \vskip14pt
    %\parbox[top][1.35cm][c]{11cm}{\huge\inserttitle \\ \small \insertsubtitle}
    \vskip11pt
    \begin{center}{\Huge Seminar Künstliche Intelligenz}
    
        \noindent\rule{\textwidth}{0.5pt} \\[0.5em]
        {\large\textcolor{ohmblue}{\inserttitle} } 
        \noindent\rule{\textwidth}{0.5pt} \\[1cm]
        \textcolor{text-grey}{Autor} \\
        {\small \insertauthor} \\[3mm]
    \end{center}
    \vfill
    \parbox[top][1.35cm][c]{11cm}{\small \insertdate} 
}
%%% end title page



%%% headline
\setbeamertemplate{headline}
{
    \vskip3pt\hspace{3.0mm}
    \begin{minipage}[b]{8cm}
        \sffamily \tiny \textcolor{text-grey}{\inserttitle} \\[0.5cm]
    \end{minipage} \hfill 
    % necessary line
    %\begin{minipage}[t]{0.75cm} 
        \includegraphics[scale=0.25]{src/logo.png}\hspace*{1em}
    %\end{minipage} 
    % necessary line
   %\begin{minipage}[b]{2.8cm} 
   %     \tiny{\textsf{\textbf{Technische Hochschule Nürnberg} \\
   %     Georg-Simon-Ohm\\[0.16cm]
   %     Fakultät Informatik }}
   % \end{minipage} \hspace{3.0mm}
    %\vskip-5pt\hfill\insertlogo\hspace{1.5mm} % logo on the right
    \vskip3pt\color{black}\rule{\textwidth}{0.4pt} % horizontal line
}
%%% end headline

%%% footline
%\newcommand{\footlinetext}{\insertshortinstitute, \insertshorttitle, \insertshortdate}
\setbeamertemplate{footline}{
    %\vskip5pt\color{black}\rule{\textwidth}{0.4pt}\\ % horizontal line

    %\hspace*{\fill} \makebox(0, 6)[c]{\insertpagenumber\,/\,\insertpresentationendpage} 
    \hspace*{\fill}%
    \llap{\usebeamertemplate*{navigation symbols} \hspace{3.5mm}}

    \vskip3pt
}
%%% end footline


\lstdefinestyle{code} {
	language		= Python,
	basicstyle		= \ttfamily,
	identifierstyle	= \ttfamily,
	stringstyle		= \ttfamily\color{ohmred},
	commentstyle	= \ttfamily\color{ohmgreen},
	keywordstyle	= \ttfamily\color{ohmblue},
	columns			= fullflexible,
	showstringspaces= false,
	emphstyle 		= \color{ohmorange},
	frame			= single,
	morekeywords	= {REFERENCES, DATETIME},
	emph			= {Michi},
	framesep 		= 15pt,
	rulesep			= 15pt,
	framerule		= 0.25pt,
	captionpos		= b,
	%linewidth		= 0.75\textwidth,
	xleftmargin 	= 15pt,
	xrightmargin	= 15pt,
	aboveskip		= 15pt,
	belowskip		= 15pt
 }


% matlabtikz bugfix
\newlength\figureheight 
\newlength\figurewidth 


\newcommand{\nicequote}[2]
{
    \begin{center}{\Huge \bfseries \rmfamily ``}
        \begin{minipage}[t]{0.75\textwidth} \itshape
            #2 
            \begin{flushright} \vspace*{-0em}
                 -- #1
            \end{flushright}
        \end{minipage} \, {\Huge \bfseries \rmfamily ''} 
    \end{center}
}




\tikzset{
    invisible/.style={opacity=0},
    visible on/.style={alt=#1{}{invisible}},
    alt/.code args={<#1>#2#3}{%
      \alt<#1>{\pgfkeysalso{#2}}{\pgfkeysalso{#3}} % \pgfkeysalso doesn't change the path
    },
  }










\title[Softcomputing] 
{Entscheidungsbäume -- Eine Einführung}

\subtitle
{}

\author[] 
{Michael~Dorner}

\institute[Softcomputing] % (optional, but mostly needed)
{Softcomputing}


\date[] 
{Erlangen, 16. Januar 2014}

\subject{Seminar Künstliche Intelligenz}

% \insertshortauthor, \insertshortinstitute, \insertshorttitle, \insertshortdate, ...
%\renewcommand{\footlinetext}{\insertshortinstitute, \insertshorttitle, \insertshortdate}

\AtBeginSubsection[]
{
    \begin{frame}<beamer>{Übersicht}
    \begin{multicols}{2}
        \tableofcontents[currentsection,currentsubsection]
    \end{multicols}
    \end{frame}
}


\begin{document}

%\renewcommand{\thealgorithm}{}


\begin{frame}[plain]
  \titlepage
\end{frame}

\begin{frame}{Übersicht}
  \begin{multicols}{2}
        \tableofcontents
    \end{multicols}
\end{frame}


\section{Entscheidungsbäume}

\begin{frame}{Einfach}

\nicequote{Albert Einstein}{Everything should be made as simple as possible, but not simpler.}

\end{frame}





\subsection{Definition}

\begin{frame}{Definition}
\subtitle{uuuuu}
\textbf{Was ist ein Entscheidungsbaum?} \\[2em]
\begin{Definition}

\begin{itemize}[<+->]
    \item Ein Entscheidungsbaum ist ein Baum mit folgenden Eigenschaften:
    \begin{itemize}[<+->]
        \item ein innerer Knoten repräsentiert ein \textcolor{ohmred}{Attribut}
        \item eine Kante repräsentiert einen Test auf dem Attribut des Vaterknotens
        \item ein Blatt repräsentiert eine der \textcolor{ohmgreen}{Klassen}
    \end{itemize}
    \item  Konstruktion eines Entscheidungsbaums
        \begin{itemize}[<+->]
            %\item Greedy Search
            \item anhand der Trainingsmenge
            \item Top-Down
        \end{itemize}
    \item Anwendung eines Entscheidungsbaums
    \begin{itemize}[<+->]
        \item Durchlaufen des Entscheidungsbaum von der Wurzel zu einem der Blätter (somit eindeutiger Pfad)
        \item Zuordnung des Objekts zur Klasse des erreichten Blattes
    \end{itemize}
\end{itemize}

\end{Definition}

\end{frame}


\subsubsection{Beispielentscheidungsbaum}

\begin{frame}{Beispiel (1)}

\textbf{Folge ich dem Vortrag aufmerksam?}\\[2em]


\begin{center}
\begin{tikzpicture}[level distance=1.5cm,
level 1/.style={sibling distance=7cm},
level 2/.style={sibling distance=2cm}]
\tikzstyle{every node}=[rectangle,draw]    
\node (Root) [ohmred] {Vortragsthema}
child [visible on=<2->] {
    node [ohmred, visible on=<4->] {Mathematiklastig}    
    child [visible on=<5->]{ 
        node [ohmgreen] {nein} 
        edge from parent node[left,draw=none] {ja}
    }
    child [visible on=<6->]{ 
        node [ohmred, visible on=<7->] {Beginn} 
        child [visible on=<8->]{
            node [ohmgreen] {nein}
            edge from parent node[left,draw=none] {Einsteinzitat}
        }
        child [visible on=<9->]{
            node [ohmgreen] {ja}
            edge from parent node[right,draw=none] {lustiges Katzenbild}
        }
        edge from parent node[right,draw=none] {nein}
    }
    edge from parent node[left,draw=none] {interessant $\;$}
}
child [visible on=<3->] {
    node [ohmred, visible on=<10->] {Vortragender}
    child [visible on=<11->]{ 
        node [ohmgreen, visible on=<13->] {ja}         
        edge from parent node[left,draw=none] {attraktive Frau}
    }
    child [visible on=<12->]{ 
        node [ohmgreen, visible on=<14->] {nein} 
        edge from parent node[right,draw=none] {unattraktiver Mann}
    }
    edge from parent node[right,draw=none] { $\;$ uninteressant}
};
\end{tikzpicture}
\end{center}

\end{frame}


\begin{frame}{Beispiel (2)}

\textbf{Wie sieht der Entscheidungsbaum als Tabelle aus?} \\[2em] 
\centering
\begin{tabular}{| l l l l | l |} \hline
    \textbf{\textcolor{ohmred}{Vortragsthema}} & \textbf{\textcolor{ohmred}{Mathematiklastig}} & \textbf{\textcolor{ohmred}{Vortragender}} & \textbf{\textcolor{ohmred}{Beginn}} & \textbf{\textcolor{ohmgreen}{Aufmerksam?}} \\ \hline
    interessant & nein & Frau & Katzenbild & ja \\
    uninteressant & ja & Frau & Definition & ja \\
    uninteressant & ja & Mann & Definition & nein \\
    uninteressant & nein & Mann & Definition & nein \\
    interessant & nein & Mann & Katzenbild & ja \\
    interessant & ja & Mann & Definition & nein \\ \hline
\end{tabular}

\end{frame}

\subsection{Konstruktion}
\begin{frame}{Konstruktion}

\textbf{Wie konstruiere ich Entscheidungsbäume generisch?}\\[2em]



\begin{algorithm}[H]
\caption{Konstruktion von Entscheidungsbäumen}
\begin{enumerate}[<+->]
    \item Wähle ein (bestes) Attribut
    \item Erstelle Knoten für dieses Attribut
    \item Füge für jeden Attributwert eine Verzweigung zu dem neuen Knoten
    \item Partitioniere die Trainingsdaten entsprechend der Attributwerte
    \item Wiederhole \textcolor{ohmblue}{1.} -- \textcolor{ohmblue}{4.}, bis alle Daten im neuen Knoten der gleichen Klasse angehören
\end{enumerate}
\end{algorithm}

\end{frame}


\subsubsection{Beispielkonstruktion}
\begin{frame}{Beispielkonstruktion eines Entscheidungsbaums (1)} 

\textbf{Möchte ich jetzt schlafen?} \\[2em]

\pause

\begin{columns}[c]
    \begin{column}{0.5\linewidth}
        \begin{tabular}{| l l | l |}\hline
            \textbf{\textcolor{ohmred}{Müde}} & \textbf{\textcolor{ohmred}{Bett}} & \textbf{\textcolor{ohmgreen}{Schlafen?}} \\ \hline
            nein & nein & nein \\
            nein & ja & nein \\
            ja & nein & nein \\
            ja & ja & ja \\ \hline
        \end{tabular}
    \end{column} \pause
    \begin{column}{0.5\linewidth}
        \begin{tikzpicture}[level distance=1.5cm,
            level 1/.style={sibling distance=2cm},
            level 2/.style={sibling distance=1cm}]
            \tikzstyle{every node}=[rectangle,draw]    
            \node (Root) [visible on=<1->, ohmred] {Müde}
            child [visible on=<2->] {
                node [visible on=<4->, ohmred] {Bett}
                child [visible on=<5->] {
                    node [visible on=<6->, ohmgreen] {ja}
                    edge from parent node[left,draw=none] {ja}
                } 
                child [visible on=<5->] {
                    node [visible on=<7->, ohmgreen] {nein}
                    edge from parent node[right,draw=none] {nein}
               }   
                edge from parent node[left,draw=none] {ja}
            }
            child [visible on=<2->] {
                node [visible on=<3->, ohmgreen] {nein}      
                edge from parent node[right,draw=none] {nein}
            };
        \end{tikzpicture}
    \end{column}
\end{columns}



\end{frame}




\begin{frame}{Beispielkonstruktion eines Entscheidungsbaums (2)} 

\textbf{Gibt es einen weiteren/besseren Entscheidungsbaum?} \\[2em]

\pause

Ja, offensichtlich, denn \\[1em]

\begin{columns}[b]
    \column[]{0.45\textwidth} \centering
\begin{tikzpicture}[level distance=1.5cm,
level 1/.style={sibling distance=2cm},
level 2/.style={sibling distance=1cm}]
\tikzstyle{every node}=[rectangle,draw]    
\node (Root) [ohmred] {Bett}
child {
    node [ohmred] {Müde}  
    child {
        node [ohmgreen] {ja}
        edge from parent node[left,draw=none] {ja}
    } 
    child {
        node [ohmgreen] {nein}
        edge from parent node[right,draw=none] {nein}
   }   
    edge from parent node[left,draw=none] {ja}
}
child {
    node [ohmgreen] {nein}      
    edge from parent node[right,draw=none] {nein}
};
\end{tikzpicture}

    \column[]{0.45\textwidth} \centering
\begin{tikzpicture}[level distance=1.5cm,
level 1/.style={sibling distance=2cm},
level 2/.style={sibling distance=1cm}]
\tikzstyle{every node}=[rectangle,draw]    
\node (Root) [ohmred] {Müde}
child {
    node [ohmred] {Bett}  
    child {
        node [ohmgreen] {ja}
        edge from parent node[left,draw=none] {ja}
    } 
    child {
        node [ohmgreen] {nein}
        edge from parent node[right,draw=none] {nein}
   }   
    edge from parent node[left,draw=none] {ja}
}
child {
    node [ohmgreen] {nein}      
    edge from parent node[right,draw=none] {nein}
};
\end{tikzpicture}
\end{columns}

\vfill

\pause

\begin{center}
Welches ist das beste Attribut?
\end{center}

\end{frame}





\section{ID3 Algorithmus}

\subsection{Exkurs in die Physik}


\begin{frame}{Exkurs in die Physik}

\textbf{Warum haben Atomkraftwerke Kühltürme?}\\[2em]
\centering
\includegraphics[height=0.8\textheight]{src/akw.jpg}

\end{frame}



\begin{frame}{Blick zur Tafel}

\centering

\includegraphics[width=\linewidth]{src/blackboard}

\end{frame}



\subsection{Informationsgehalt und -gewinn}

\begin{frame}{Informationsgehalt \& Informationsgewinn}
\begin{Definition}

Den mittleren Informationsgehalt $H(P)$ einer Wahrscheinlichkeitsverteilung $P$ über einer endlichen Menge $S$ bezeichnet man als Entropie von $P$:
\begin{equation*}
    H(S) = - \sum_{i=1}^{} P(C_i) \log_2 P(C_i)
\end{equation*}
mit $P(C_i)$ als Auftrittswahrscheinlichkeit der Klasse $C_i$ in $S$.

\end{Definition}

\pause

\begin{Definition}
Um den Informationsgewinn (\textit{Information Gain}) von Attribut $A$ zu quantifizieren, bilden wir die Differenz der ursprünglichen Information und der Restinformation:
\begin{equation*}
G(S, A) = H(S) - \sum_{i \in \text{Werte}(A)} \frac{|S_i|}{|S|}H(S_i)
\end{equation*}
mit \begin{itemize}
 \item $\text{Werte}(A)$ als alle Ausprägungen von $A$ und
 \item $S_i$ als Teilmenge von $S$, wobei $A$ den Wert $i$ annimmt.
 \end{itemize}

\end{Definition}


\end{frame}

\subsubsection{Beispiel}


\begin{frame}{Beispiel (1)}

\textbf{Welchen Informationsgehalt (Entropie) haben folgende Nachrichten?}\\[1em]

\begin{itemize}[<+->]
    \item \texttt{'aaaa'}: 
        \begin{itemize}
            \item Wahrscheinlichkeitsverteilung: 
            \[P(\texttt{'a'}) = \frac{4}{4}\]        
            \item Entropie:
    \[H(\texttt{'aaaa'}) = - \left( \frac{4}{4} \log_2 \frac{4}{4} \right) = 0 \]
        \end{itemize}   
    \item \texttt{'aabc'}: 
        \begin{itemize}[<+->]
            \item Wahrscheinlichkeitsverteilungen: 
            \[P(\texttt{'a'}) = \frac{2}{4}, \qquad P(\texttt{'b'}) = P(\texttt{'c'}) = \frac{1}{4}\]        
            \item Entropie:
    \[H(\texttt{'aabc'}) = - \left( \frac{2}{4} \log_2 \frac{2}{4} + \frac{1}{4} \log_2 \frac{1}{4} + \frac{1}{4} \log_2 \frac{1}{4} \right) = 1.5 \]        
        \end{itemize}        
\end{itemize}

\end{frame}



\begin{frame}{Beispiele (2)}

\begin{itemize}[<+->]
   \item \texttt{'abcd'}: 
        \begin{itemize}[<+->]
            \item Wahrscheinlichkeitsverteilungen: 
            \[P(\texttt{'a'}) = P(\texttt{'b'}) = P(\texttt{'c'}) = P(\texttt{'d'}) = \frac{1}{4}\]        
            \item Entropie:
    \[H(\texttt{'abcd'}) = - \left( \frac{1}{4} \log_2 \frac{1}{4} + \frac{1}{4} \log_2 \frac{1}{4} + \frac{1}{4} \log_2 \frac{1}{4} + \frac{1}{4} \log_2 \frac{1}{4} \right) = 2 \]
        \end{itemize} 
    \item \texttt{'aaaabcdefg'}: 
        \begin{itemize}[<+->]
            \item Wahrscheinlichkeitsverteilungen: 
            \[P(\texttt{'a'}) = \frac{4}{10}, \quad P(\texttt{'b'}) = P(\texttt{'c'}) = P(\texttt{'d'}) = P(\texttt{'e'}) = P(\texttt{'f'}) = P(\texttt{'g'}) = \frac{1}{10}\]        
            \item Entropie:
    \[H(\texttt{'aaaabcdefg'}) = - \left( \frac{4}{10} \log_2 \frac{4}{10} + 6 \cdot \left( \frac{1}{10} \log_2 \frac{1}{10} \right) \right) = 2{,}52193 \]
        \end{itemize}
\end{itemize} 
\end{frame}



\subsection{Pseudocode}
\begin{frame}{Pseudocode}
\begin{algorithm}[H]
\caption{ID3 Algorithmus}
\begin{enumerate}
    \item Wähle \sout{\textcolor{gray}{ein (bestes) Attribut}} das Attribut mit größtem Informationgewinn
    \item Erstelle Knoten für dieses Attribut
    \item Füge für jeden Attributwert eine Verzweigung zu dem neuen Knoten
    \item Partitioniere die Trainingsdaten entsprechend der Attributwerte
    \item Wiederhole \textcolor{ohmblue}{1.} -- \textcolor{ohmblue}{4.}, bis alle Daten im neuen Knoten der gleichen Klasse angehören
\end{enumerate}
\end{algorithm}
\end{frame}

\subsubsection{Beispiel}



\begin{frame}{Beispiel}
\textbf{Informationsgehalt von Krankheit}\\[1em]

\begin{tabular}{|l l l l l l | l |}\hline
\textbf{Patient} & \textbf{\textcolor{ohmred}{Fieber}} & \textbf{\textcolor{ohmred}{Husten}} & \textbf{\textcolor{ohmred}{Röntgen}} & \textbf{\textcolor{ohmred}{BSG}} & \textbf{\textcolor{ohmred}{Abhören}} &\textbf{\textcolor{ohmgreen}{Krankheit}}\\ \hline
1 & hoch & stark & Flocken & normal & blubbernd & Lungenentzündung \\
2 & mittel & stark & Flocken & normal & blubbernd & Lungenentzündung \\
3 & niedrig & leicht & Punkt & normal & fiepend & Lungenentzündung \\
4 & hoch & mittel & Flocken & normal & blubbernd & Lungenentzündung \\
5 & mittel & leicht & Flocken & normal & blubbernd & Lungenentzündung \\
6 & ohne & leicht & Streifen & normal & normal & Tuberkulose \\
7 & hoch & stark & Loch & schnell & fiepend & Tuberkulose \\
8 & niedrig & leicht & Streifen & normal & normal & Tuberkulose \\
9 & ohne & leicht & Punkt & schnell & fiepend & Tuberkulose \\
10 & niedrig & mittel & Flocken & schnell & normal & Tuberkulose \\ \hline 
\end{tabular}

\vfill

\pause 

\[H(S) = - \sum_{i=1}^{} P(C_i) \log_2 P(C_i) = - \nicefrac{5}{10} \log_2 \nicefrac{5}{10} - \nicefrac{5}{10} \log_2 \nicefrac{5}{10} = 1\]

\end{frame}



\begin{frame}{Beispiel}

\textbf{Informationsgehalt für Fieber}\\[1em]

\begin{tabular}{|l l l l l l | l |}\hline
\textbf{Patient} & \textbf{\textcolor{ohmred}{Fieber}} & \textbf{\textcolor{ohmred}{Husten}} & \textbf{\textcolor{ohmred}{Röntgen}} & \textbf{\textcolor{ohmred}{BSG}} & \textbf{\textcolor{ohmred}{Abhören}} &\textbf{\textcolor{ohmgreen}{Krankheit}}\\ \hline
1 & hoch & stark & Flocken & normal & blubbernd & Lungenentzündung \\
2 & mittel & stark & Flocken & normal & blubbernd & Lungenentzündung \\
3 & niedrig & leicht & Punkt & normal & fiepend & Lungenentzündung \\
4 & hoch & mittel & Flocken & normal & blubbernd & Lungenentzündung \\
5 & mittel & leicht & Flocken & normal & blubbernd & Lungenentzündung \\
6 & ohne & leicht & Streifen & normal & normal & Tuberkulose \\
7 & hoch & stark & Loch & schnell & fiepend & Tuberkulose \\
8 & niedrig & leicht & Streifen & normal & normal & Tuberkulose \\
9 & ohne & leicht & Punkt & schnell & fiepend & Tuberkulose \\
10 & niedrig & mittel & Flocken & schnell & normal & Tuberkulose \\ \hline 
\end{tabular}

\vfill

\begin{tabular}{| l<{\onslide<2->} | l<{\onslide<3->} | l<{\onslide<5->} | l<{\onslide<6->} | l<{\onslide<7->} | l<{\onslide} |}\hline 
\textbf{\textcolor{ohmred}{Fieber}} & \textbf{Anzahl} & \textbf{LE} & \textbf{TBC} & \textbf{Formel} & \textbf{Ergebnis} \\ \hline
hoch & 3 & $ \nicefrac{2}{3}$ & $ \nicefrac{1}{3}$ & $ \nicefrac{3}{10} \cdot \left(- \nicefrac{2}{3} \log_2 \nicefrac{2}{3} - \nicefrac{1}{3} \log_2 \nicefrac{1}{3} \right)$ & $=0{,}2755$\\ \hline
mittel & 2 & $ \nicefrac{2}{2}$ & $\nicefrac{0}{2}$ & $ \nicefrac{2}{10} \cdot \left(- \nicefrac{2}{2} \log_2 \nicefrac{2}{2} - \nicefrac{0}{2} \log_2 \nicefrac{0}{2} \right)$ & $=0$\\ \hline
niedrig & 3 & $\nicefrac{1}{3}$ & $\nicefrac{2}{3}$ & $\nicefrac{3}{10} \cdot \left(- \nicefrac{1}{3} \log_2 \nicefrac{1}{3} - \nicefrac{2}{3} \log_2 \nicefrac{2}{3} \right)$ & $=0{,}2755$\\ \hline
ohne & 2 & $\nicefrac{0}{2}$ & $\nicefrac{2}{2}$ & $\nicefrac{2}{10} \cdot \left(- \nicefrac{0}{2} \log_2 \nicefrac{0}{2} - \nicefrac{2}{2} \log_2 \nicefrac{2}{2} \right)$ & $=0$\\ \hline
\multicolumn{5}{r|}{$\boldsymbol{\sum}$} & \onslide<8->{$0{,}5510$} \\ \cline{6-6}
\end{tabular}


\end{frame}




\begin{frame}{Beispiel}

\textbf{Informationsgewinn für Fieber}

\[ \displaystyle
\begin{array}{r l l l}
    G(S, \text{Fieber}) &= H(S) &- \sum\limits_{\substack{i \in \{\text{hoch, mittel,} \\ \text{niedrig, ohne}\}}} \cfrac{|S_i|}{|S|}H(S_i) &= \pause \\[2em]
    &= 1 &-0{,}5510 &= 0{,}4490 
\end{array}
\]

\end{frame}



\begin{frame}{Beispiel}

\textbf{Informationsgehalt für Husten} \\[1em]

\begin{tabular}{|l l l l l l | l |}\hline
\textbf{Patient} & \textbf{\textcolor{ohmred}{Fieber}} & \textbf{\textcolor{ohmred}{Husten}} & \textbf{\textcolor{ohmred}{Röntgen}} & \textbf{\textcolor{ohmred}{BSG}} & \textbf{\textcolor{ohmred}{Abhören}} &\textbf{\textcolor{ohmgreen}{Krankheit}}\\ \hline
1 & hoch & stark & Flocken & normal & blubbernd & Lungenentzündung \\
2 & mittel & stark & Flocken & normal & blubbernd & Lungenentzündung \\
3 & niedrig & leicht & Punkt & normal & fiepend & Lungenentzündung \\
4 & hoch & mittel & Flocken & normal & blubbernd & Lungenentzündung \\
5 & mittel & leicht & Flocken & normal & blubbernd & Lungenentzündung \\
6 & ohne & leicht & Streifen & normal & normal & Tuberkulose \\
7 & hoch & stark & Loch & schnell & fiepend & Tuberkulose \\
8 & niedrig & leicht & Streifen & normal & normal & Tuberkulose \\
9 & ohne & leicht & Punkt & schnell & fiepend & Tuberkulose \\
10 & niedrig & mittel & Flocken & schnell & normal & Tuberkulose \\ \hline 
\end{tabular}

\vfill

\begin{tabular}{| l<{\onslide<2->} | l<{\onslide<3->} | l<{\onslide<5->} | l<{\onslide<6->} | l<{\onslide<7->} | l<{\onslide} |}\hline 
\textbf{\textcolor{ohmred}{Husten}} & \textbf{Anzahl} & \textbf{LE} & \textbf{TBC} & \textbf{Formel} & \textbf{Ergebnis} \\ \hline  
stark & 3 & $ \nicefrac{2}{3}$ & $ \nicefrac{1}{3}$ & $ \nicefrac{3}{10} \cdot \left(- \nicefrac{2}{3} \log_2 \nicefrac{2}{3} - \nicefrac{1}{3} \log_2 \nicefrac{1}{3} \right)$ & $=0{,}2755$ \\ \hline 
mittel & 2 & $ \nicefrac{1}{2}$ & $\nicefrac{1}{2}$ & $ \nicefrac{2}{10} \cdot \left(- \nicefrac{1}{2} \log_2 \nicefrac{1}{2} - \nicefrac{1}{2} \log_2 \nicefrac{1}{2} \right)$ & $=0{,}2$  \\ \hline 
leicht & 5 & $\nicefrac{2}{5}$ & $\nicefrac{4}{5}$ & $\nicefrac{3}{10} \cdot \left(- \nicefrac{2}{5} \log_2 \nicefrac{2}{5} - \nicefrac{3}{5} \log_2 \nicefrac{3}{5} \right)$ & $=0{,}4855$  \\ \hline  
\multicolumn{5}{r|}{$\boldsymbol{\sum}$} & \onslide<8->{$0{,}9610$} \\ \cline{6-6}
\end{tabular}

\end{frame}


\begin{frame}{Beispiel}

\textbf{Informationsgewinn für Husten}

\[
\begin{array}{r l l l}
    G(S, \text{Husten}) &= H(S) &- \sum\limits_{\substack{i \in \{\text{stark, } \\ \text{mittel,leicht}\}}} \cfrac{|S_i|}{|S|}H(S_i) &= \pause \\[1em]
    &= 1 &-0{,}9610 &= 0{,}0390 
\end{array}
\]

\end{frame}





\begin{frame}{Beispiel}

\textbf{Informationsgehalt für Abhören} \\[1em]

\begin{tabular}{|l l l l l l | l |}\hline
\textbf{Patient} & \textbf{\textcolor{ohmred}{Fieber}} & \textbf{\textcolor{ohmred}{Husten}} & \textbf{\textcolor{ohmred}{Röntgen}} & \textbf{\textcolor{ohmred}{BSG}} & \textbf{\textcolor{ohmred}{Abhören}} &\textbf{\textcolor{ohmgreen}{Krankheit}}\\ \hline
1 & hoch & stark & Flocken & normal & blubbernd & Lungenentzündung \\
2 & mittel & stark & Flocken & normal & blubbernd & Lungenentzündung \\
3 & niedrig & leicht & Punkt & normal & fiepend & Lungenentzündung \\
4 & hoch & mittel & Flocken & normal & blubbernd & Lungenentzündung \\
5 & mittel & leicht & Flocken & normal & blubbernd & Lungenentzündung \\
6 & ohne & leicht & Streifen & normal & normal & Tuberkulose \\
7 & hoch & stark & Loch & schnell & fiepend & Tuberkulose \\
8 & niedrig & leicht & Streifen & normal & normal & Tuberkulose \\
9 & ohne & leicht & Punkt & schnell & fiepend & Tuberkulose \\
10 & niedrig & mittel & Flocken & schnell & normal & Tuberkulose \\ \hline 
\end{tabular}

\vfill

\begin{tabular}{| l<{\onslide<2->} | l<{\onslide<3->} | l<{\onslide<5->} | l<{\onslide<6->} | l<{\onslide<7->} | l<{\onslide} |}\hline 
\textbf{\textcolor{ohmred}{Fieber}} & \textbf{Anzahl} & \textbf{LE} & \textbf{TBC} & \textbf{Formel} & \textbf{Ergebnis} \\ \hline 
blubbernd & 4 & $ \nicefrac{4}{4}$ & $ \nicefrac{0}{4}$ & $ \nicefrac{4}{10} \cdot \left(- \nicefrac{4}{4} \log_2 \nicefrac{4}{4} - \nicefrac{0}{4} \log_2 \nicefrac{0}{4} \right)$ & $=0$\\ \hline 
fiepend & 3 & $ \nicefrac{1}{3}$ & $\nicefrac{2}{3}$ & $ \nicefrac{3}{10} \cdot \left(- \nicefrac{1}{3} \log_2 \nicefrac{1}{3} - \nicefrac{2}{3} \log_2 \nicefrac{2}{3} \right)$ & $=0{,}2755$\\ \hline 
normal & 3 & $\nicefrac{0}{3}$ & $\nicefrac{3}{3}$ & $\nicefrac{3}{10} \cdot \left(- \nicefrac{0}{3} \log_2 \nicefrac{0}{3} - \nicefrac{3}{3} \log_2 \nicefrac{3}{3} \right)$ & $=0$\\ \hline 
\multicolumn{5}{r|}{$\boldsymbol{\sum}$} & \onslide<8->{$0{,}2755$} \\ \cline{6-6}
\end{tabular}


\end{frame}


\begin{frame}{Beispiel}

Informationsgewinn für Abhören

\[
\begin{array}{r l l l}
    G(S, \text{Abhören}) &= H(S) &- \sum\limits_{\substack{i \in \{\text{blubbernd} \\ \text{fiepend, normal}\}}} \cfrac{|S_i|}{|S|}H(S_i) &= \pause \\[1em]
    &= 1 &-0{,}2755 &= 0{,}7245 
\end{array}
\]

\end{frame}



\begin{frame}{Beispiel}

\textbf{Wahl des ersten Knoten} \\[1em]

\begin{center}
\begin{tabular}{|l l l l l l|} \hline
\textbf{\textcolor{ohmred}{Fieber}} & \textbf{\textcolor{ohmred}{Husten}} & \textbf{\textcolor{ohmred}{Röntgen}} & \textbf{\textcolor{ohmred}{BSG}} & \textbf{\textcolor{ohmred}{Abhören}} &\textbf{\textcolor{ohmgreen}{Krankheit}}\\ \hline
0{,}4490 & 0{,}0390 & 0{,}4390 & 0{,}3958 & 0{,}7245 & 1 \\ \hline
\end{tabular}

\vfill

\pause

\begin{tikzpicture}[level distance=1.5cm,
level 1/.style={sibling distance=4cm},
level 2/.style={sibling distance=2cm}]
\tikzstyle{every node}=[rectangle,draw]    
\node (Root) [ohmred] {Abhören}
child {
    node [ohmgreen] {Lungenentzündung}    
    edge from parent node[left,draw=none] {blubbernd $\;$}
}
child {
    node [ohmred] {?}
    edge from parent node[right,draw=none] { $\;$ fiepend}
}
child {
    node [ohmgreen] {TBC}
    edge from parent node[right,draw=none] { $\;$ normal}
};
\end{tikzpicture}
\end{center}


\end{frame}




\begin{frame}{Beispiel}

\textbf{Informationsgehalt für Fieber} \\[1em]

\begin{tabu}{|l l l l l l | l |}\hline
\textbf{Patient} & \textbf{\textcolor{ohmred}{Fieber}} & \textbf{\textcolor{ohmred}{Husten}} & \textbf{\textcolor{ohmred}{Röntgen}} & \textbf{\textcolor{ohmred}{BSG}} & \textbf{\textcolor{ohmred}{Abhören}} &\textbf{\textcolor{ohmgreen}{Krankheit}}\\ \hline
\rowfont{\color{gray}}1 & hoch & stark & Flocken & normal & blubbernd & Lungenentzündung \\
\rowfont{\color{gray}}2 & mittel & stark & Flocken & normal & blubbernd & Lungenentzündung \\
3 & niedrig & leicht & Punkt & normal & fiepend & Lungenentzündung \\
\rowfont{\color{gray}}4 & hoch & mittel & Flocken & normal & blubbernd & Lungenentzündung \\
\rowfont{\color{gray}}5 & mittel & leicht & Flocken & normal & blubbernd & Lungenentzündung \\
\rowfont{\color{gray}}6 & ohne & leicht & Streifen & normal & normal & Tuberkulose \\
7 & hoch & stark & Loch & schnell & fiepend & Tuberkulose \\
\rowfont{\color{gray}}8 & niedrig & leicht & Streifen & normal & normal & Tuberkulose \\
9 & ohne & leicht & Punkt & schnell & fiepend & Tuberkulose \\
\rowfont{\color{gray}}10 & niedrig & mittel & Flocken & schnell & normal & Tuberkulose \\ \hline 
\end{tabu}

\vfill

\begin{tabular}{| l<{\onslide<2->} | l<{\onslide<3->} | l<{\onslide<5->} | l<{\onslide<6->} | l<{\onslide<7->} | l<{\onslide} |}\hline 
\textbf{\textcolor{ohmred}{Fieber}} & \textbf{Anzahl} & \textbf{LE} & \textbf{TBC} & \textbf{Formel} & \textbf{Ergebnis} \\ \hline 
niedrig & 1 & $ \nicefrac{1}{1}$ & $ \nicefrac{0}{1}$ & $ \nicefrac{1}{3} \cdot \left(- \nicefrac{1}{1} \log_2 \nicefrac{1}{1} - \nicefrac{0}{1} \log_2 \nicefrac{0}{1} \right)$ & $=0$\\ \hline 
hoch & 1 & $ \nicefrac{0}{1}$ & $\nicefrac{1}{1}$ & $ \nicefrac{1}{3} \cdot \left(- \nicefrac{0}{1} \log_2 \nicefrac{0}{1} - \nicefrac{1}{1} \log_2 \nicefrac{1}{1} \right)$ & $=0$\\ \hline 
ohne & 1 & $\nicefrac{0}{1}$ & $\nicefrac{1}{1}$ & $\nicefrac{1}{3} \cdot \left(- \nicefrac{0}{1} \log_2 \nicefrac{0}{1} - \nicefrac{1}{11} \log_2 \nicefrac{1}{1} \right)$ & $=0$\\ \hline 
\multicolumn{5}{r|}{$\boldsymbol{\sum}$} & \onslide<7->{$0$} \\ \cline{6-6}
\end{tabular}


\end{frame}


\begin{frame}{Beispiel}

\textbf{Informationsgewinn für Fieber}\\[1em]

\[
\begin{array}{r l l l}
    G(S, \text{Fieber}) &= H(S) &- \sum\limits_{\substack{i \in \{\text{niedrig} \\ \text{hoch, normal}\}}} \cfrac{|S_i|}{|S|}H(S_i) &= \pause \\[1em]
    &= 0{,}9183 &-0 &= 0{,}9183
\end{array}
\]

\end{frame}


\begin{frame}{Beispiel}

\textbf{Resultierender Entscheidungsbaum} \\[2em]

\begin{center}

\begin{tikzpicture}[level distance=1.5cm,
level 1/.style={sibling distance=4cm},
level 2/.style={sibling distance=3cm}]
\tikzstyle{every node}=[rectangle,draw]    
\node (Root) [ohmred] {Abhören}
child {
    node [ohmgreen] {Lungenentzündung}    
    edge from parent node[left,draw=none] {blubbernd $\quad$}
}
child {
    node [ohmred] {Fieber}
    child {
        node [ohmgreen] {Lungenentzündung}    
        edge from parent node[left,draw=none] {niedrig $\;$}
    }
    child {
        node [ohmgreen] {TBC}    
        edge from parent node[left,draw=none] {hoch }
    }
    child {
        node [ohmgreen] {TBC}    
        edge from parent node[right,draw=none] {$\;$ ohne}
    }
    edge from parent node[right,draw=none] {fiepend}
}
child {
    node [ohmgreen] {TBC}
    edge from parent node[right,draw=none] { $\quad$ normal}
};
\end{tikzpicture}
\end{center}


\end{frame}





\section{Anwendung -- Bewegungsdaten}
\subsection{Einführung}
\begin{frame}{Einführung}
\textbf{Allgemeines}\\[1em]
\begin{itemize}[<+->]
    \item Daten aus meiner Bachelorarbeit
    \item Nicht vollständig
    \item Uniformisierte Daten (vier Klassen à 876 Samples)
    \item GPS Geschwindigkeit und Beschleunigungslevel
\end{itemize}

\vfill \pause

\textbf{Berechnung}\\[1em]
Betrag der Beschleunigung 
\begin{equation}
    m_i = \sqrt{x_i^2 + y_i^2 + z_i^2}, \qquad i = 1, \dots, n 
\end{equation}
mit $x_i, y_i, z_i$ als Beschleunigungswerte entlang der entsprechenden Achse für jeden Datenpunkt $i$ der gesamtem Samples $n$.

\end{frame}


\begin{frame}{Filterung}
    \centering 
    \setlength\figureheight{0.4\textwidth} 
    \setlength\figurewidth{0.8\textwidth} 
    % This file was created by matlab2tikz v0.4.3.
% Copyright (c) 2008--2013, Nico Schlömer <nico.schloemer@gmail.com>
% All rights reserved.
% 
% The latest updates can be retrieved from
%   http://www.mathworks.com/matlabcentral/fileexchange/22022-matlab2tikz
% where you can also make suggestions and rate matlab2tikz.
% 
\begin{tikzpicture}[font=\sffamily\sansmath]

\begin{axis}[%
width=\figurewidth,
height=\figureheight,
scale only axis,
xmin=0,
xmax=1000,
xtick={0,100,200,300,400,500,600,700,800,900,1000},
xticklabels={0,1,2,3,4,5,6,7,8,9,10},
xlabel={Zeit in s},
ymin=0,
ymax=2,
ylabel={Beschleunigungslevel},
legend style={draw=black,fill=white,legend cell align=left}
]
\addplot [
color=ohmorange,
solid
]
table[row sep=crcr]{
1 0.826946236625417\\
2 0.89345639509675\\
3 0.836697025637272\\
4 0.794763592560089\\
5 0.690463613708784\\
6 0.80504532358636\\
7 1.00175845628395\\
8 0.874739427751146\\
9 0.936489126429378\\
10 0.70770365811582\\
11 1.12145981765175\\
12 1.96329758254411\\
13 1.54249963413657\\
14 0.187570380783917\\
15 0.322573670603304\\
16 0.390554372703014\\
17 1.5103800538576\\
18 1.6458512846547\\
19 1.42938880339412\\
20 1.06533372635507\\
21 0.789764119631556\\
22 0.48730917098695\\
23 0.222803330563807\\
24 0.292805962523945\\
25 0.431114450369552\\
26 0.372627595984717\\
27 0.314140741599902\\
28 0.255653887215067\\
29 0.307019932996571\\
30 0.438606706491643\\
31 0.347023311983369\\
32 0.293337260900488\\
33 0.157018524725836\\
34 0.155631279172974\\
35 0.210927427488056\\
36 0.164328269285296\\
37 0.173452077847861\\
38 0.212735041833654\\
39 0.247005582746526\\
40 0.259406136931726\\
41 0.386154789228354\\
42 0.374967024543669\\
43 0.327469509720848\\
44 0.309957737336816\\
45 0.284405110384139\\
46 0.262473715324631\\
47 0.248227423535248\\
48 0.148617140378895\\
49 0.106963346733287\\
50 0.129051557624627\\
51 0.271581041210228\\
52 0.284938445929823\\
53 0.266505670630534\\
54 0.2329280532872\\
55 0.213617794089397\\
56 0.34020632577597\\
57 0.342286837128385\\
58 0.3443673484808\\
59 0.346447859833214\\
60 0.348528371185628\\
61 0.350608882538042\\
62 0.352689393890457\\
63 0.354769905242871\\
64 0.356850416595285\\
65 0.358930927947699\\
66 0.361011439300114\\
67 0.363091950652528\\
68 0.365172462004942\\
69 0.367252973357357\\
70 0.369333484709771\\
71 0.371413996062185\\
72 0.373494507414599\\
73 0.375575018767014\\
74 0.377655530119428\\
75 0.379736041471842\\
76 0.381816552824257\\
77 0.387388059574251\\
78 0.448043612271019\\
79 0.4777536746224\\
80 0.413866972860873\\
81 0.391117347005782\\
82 0.450063884915062\\
83 0.612647496666893\\
84 0.883840045010305\\
85 0.788200599576318\\
86 0.564632698094378\\
87 0.602837284378904\\
88 0.687038226076507\\
89 0.604945678796504\\
90 0.302804598317828\\
91 0.496438930332615\\
92 0.532215295323869\\
93 0.583256833554523\\
94 0.521382431937511\\
95 0.41689997746759\\
96 0.539272596283708\\
97 0.656477860763717\\
98 0.675024459242232\\
99 0.528593213910749\\
100 0.467096789580911\\
101 0.538861402325346\\
102 0.624693386267935\\
103 0.634141656862085\\
104 0.625959075292732\\
105 0.54620695823444\\
106 0.45999544940353\\
107 0.41574256944031\\
108 0.418481352893212\\
109 0.396338506578263\\
110 0.325463169432679\\
111 0.295334424075992\\
112 0.306164252188163\\
113 0.316994080300331\\
114 0.598224684513932\\
115 0.971384486033635\\
116 1.16449056196194\\
117 1.23654349351869\\
118 1.08757719995928\\
119 0.640670126211815\\
120 0.394412802756728\\
121 0.283236242606362\\
122 0.28752366495582\\
123 0.428504875455565\\
124 0.564631369042774\\
125 0.712556874426314\\
126 1.16246888329986\\
127 1.33597734720925\\
128 1.16697489447683\\
129 1.10244877806555\\
130 1.05483055550019\\
131 0.736526939648251\\
132 0.646112873347704\\
133 0.611507981943704\\
134 0.606708882655337\\
135 0.603946339309033\\
136 0.601183795962729\\
137 0.598421252616424\\
138 0.595658709270121\\
139 0.592896165923816\\
140 0.590133622577513\\
141 0.587371079231208\\
142 0.584608535884904\\
143 0.5818459925386\\
144 0.579083449192295\\
145 0.576320905845992\\
146 0.573558362499687\\
147 0.570795819153384\\
148 0.568033275807079\\
149 0.565270732460776\\
150 0.562508189114471\\
151 0.559745645768166\\
152 0.556983102421863\\
153 0.554220559075558\\
154 0.551458015729255\\
155 0.643172495300215\\
156 0.823707081206004\\
157 0.833788864771492\\
158 0.730363322784545\\
159 0.391428862434961\\
160 0.18450880253225\\
161 0.209297609815698\\
162 0.119410101762262\\
163 0.189821555245313\\
164 0.384286245304908\\
165 0.53921603347069\\
166 0.409557849801379\\
167 0.16769287122947\\
168 0.286710611469116\\
169 0.257590837216216\\
170 0.160981545445401\\
171 0.0887418560586131\\
172 0.0679166537664057\\
173 0.0847971174909366\\
174 0.0811804248968975\\
175 0.177061316763884\\
176 0.169713628752862\\
177 0.0734990040787747\\
178 0.13137446813284\\
179 0.116382542348094\\
180 0.215071975583583\\
181 0.310133629931576\\
182 0.336365108005066\\
183 0.379014016937426\\
184 0.383257956172106\\
185 0.368280047182998\\
186 0.361269456790221\\
187 0.352749539512899\\
188 0.295668461884676\\
189 0.268006135316513\\
190 0.292589251100715\\
191 0.321897091230589\\
192 0.376284411834378\\
193 0.358090298543325\\
194 0.314736872870454\\
195 0.296303807531941\\
196 0.462456097933091\\
197 0.45237526063644\\
198 0.442635886918503\\
199 0.541407533287959\\
200 0.644656835823282\\
201 0.647080422142428\\
202 0.63884495194015\\
203 0.630609481737869\\
204 0.622374011535592\\
205 0.614138541333311\\
206 0.605903071131033\\
207 0.597667600928752\\
208 0.589432130726471\\
209 0.581196660524193\\
210 0.572961190321913\\
211 0.564725720119635\\
212 0.556490249917354\\
213 0.548254779715076\\
214 0.540019309512795\\
215 0.531783839310518\\
216 0.523548369108237\\
217 0.515312898905956\\
218 0.507077428703678\\
219 0.498841958501397\\
220 0.49060648829912\\
221 0.482371018096839\\
222 0.474135547894561\\
223 0.46590007769228\\
224 0.457664607489999\\
225 0.449429137287721\\
226 0.259765603232901\\
227 0.276192216258961\\
228 0.294908203741661\\
229 0.317395415571093\\
230 0.304854266121399\\
231 0.343969809526213\\
232 0.442841095338768\\
233 0.382438942426169\\
234 0.377904172726308\\
235 0.334220619746717\\
236 0.360466344883407\\
237 0.535697218215304\\
238 0.571385405774414\\
239 0.569120451889954\\
240 0.570624227477278\\
241 0.541562180495263\\
242 0.525863736995735\\
243 0.601180818782732\\
244 0.578429557536713\\
245 0.721668023603801\\
246 0.564395330046985\\
247 0.476909642178322\\
248 0.451520156833312\\
249 0.364976424214218\\
250 0.383162602403364\\
251 0.544125832495188\\
252 0.776803183913567\\
253 0.934429395865476\\
254 0.943993822243762\\
255 0.593819712622071\\
256 0.257252852758264\\
257 0.182140389883617\\
258 0.257770651125644\\
259 0.595787676393185\\
260 1.00635745965074\\
261 1.313042990733\\
262 1.45371685616321\\
263 1.27304343056384\\
264 0.944084992942152\\
265 1.07978426809779\\
266 1.06506622218476\\
267 0.859227466832736\\
268 0.843208675816808\\
269 0.465019363676038\\
270 1.05279249674395\\
271 1.2721228693386\\
272 0.704508059619473\\
273 0.340722439288971\\
274 0.238737787034483\\
275 0.0580316741996017\\
276 0.0674424920644725\\
277 0.07685330992934\\
278 0.0862641277942109\\
279 0.0956749456590784\\
280 0.105085763523949\\
281 0.114496581388817\\
282 0.123907399253688\\
283 0.133318217118559\\
284 0.142729034983426\\
285 0.152139852848297\\
286 0.161550670713164\\
287 0.170961488578035\\
288 0.180372306442903\\
289 0.189783124307774\\
290 0.199193942172641\\
291 0.208604760037512\\
292 0.218015577902383\\
293 0.22742639576725\\
294 0.236837213632121\\
295 0.246248031496989\\
296 0.25565884936186\\
297 0.265069667226727\\
298 0.288801561380046\\
299 0.185235217071515\\
300 0.214296090380335\\
301 0.253082733664612\\
302 0.303280711184584\\
303 0.398894944198058\\
304 0.443830770010657\\
305 0.434095567537413\\
306 0.407319301742352\\
307 0.366011147834084\\
308 0.325740928441419\\
309 0.285878359190837\\
310 0.260538833527603\\
311 0.127212660616881\\
312 0.253532272706978\\
313 0.208270814592016\\
314 0.118283183976558\\
315 0.134764658722759\\
316 0.121829729607972\\
317 0.075995615897289\\
318 0.0502022518904005\\
319 0.0533704631480991\\
320 0.104698024644272\\
321 0.113668893600329\\
322 0.0870765118942553\\
323 0.100153516672669\\
324 0.114842397529064\\
325 0.157335041235668\\
326 0.18421682860352\\
327 0.203689621573069\\
328 0.234657887676317\\
329 0.254791229479862\\
330 0.274473711689021\\
331 0.294949444707673\\
332 0.297102393260791\\
333 0.297171939119152\\
334 0.28902154190809\\
335 0.314943642598546\\
336 0.31937380918399\\
337 0.399654605710203\\
338 0.429747934801652\\
339 0.42953310942894\\
340 0.463495148608492\\
341 0.265646902126968\\
342 0.178161883832239\\
343 0.195703560284496\\
344 0.35936432159843\\
345 0.259744722171167\\
346 0.440961047190813\\
347 0.506296214683772\\
348 0.464412270661576\\
349 0.229781261588603\\
350 0.219892410257484\\
351 0.219961295998078\\
352 0.220030181738671\\
353 0.220099067479265\\
354 0.220167953219858\\
355 0.220236838960452\\
356 0.220305724701045\\
357 0.220374610441638\\
358 0.220443496182232\\
359 0.220512381922825\\
360 0.220581267663419\\
361 0.220650153404012\\
362 0.220719039144606\\
363 0.220787924885199\\
364 0.220856810625792\\
365 0.220925696366386\\
366 0.220994582106979\\
367 0.221063467847573\\
368 0.221132353588166\\
369 0.22120123932876\\
370 0.221270125069353\\
371 0.221339010809947\\
372 0.22140789655054\\
373 0.221476782291134\\
374 0.221545668031727\\
375 0.22161455377232\\
376 0.233652831383778\\
377 0.251646965663161\\
378 0.360421897150029\\
379 0.396885723337714\\
380 0.446085874136308\\
381 0.422999482418163\\
382 0.397042589456323\\
383 0.452794924448968\\
384 0.380714258115937\\
385 0.401105789355871\\
386 0.452163278954503\\
387 0.491451321883851\\
388 0.507012873093353\\
389 0.511874445006997\\
390 0.496584481271502\\
391 0.498571506135273\\
392 0.549836272538536\\
393 0.556670482423889\\
394 0.619761651858077\\
395 0.711127341580251\\
396 0.664267870097678\\
397 0.566249322294779\\
398 0.519909236212152\\
399 0.417645648721629\\
400 0.352711591444409\\
401 0.28960214652324\\
402 0.300061281658726\\
403 0.318113719902467\\
404 0.357621801005327\\
405 0.443017303957852\\
406 0.64806397147621\\
407 0.786674546596615\\
408 0.888925757511283\\
409 0.964858792676511\\
410 1.04079182784174\\
411 0.994719473728443\\
412 0.91945784823023\\
413 0.70280157290871\\
414 0.418985084431052\\
415 0.231683256837635\\
416 0.377859973231387\\
417 0.481584398898112\\
418 0.49484481696439\\
419 0.902417440644516\\
420 0.985879631241601\\
421 1.16158840906024\\
422 1.16174152577028\\
423 1.02456052777298\\
424 1.07801760907767\\
425 0.896230598266811\\
426 0.688089383505781\\
427 0.749854727894118\\
428 0.869503815321064\\
429 0.857040306673\\
430 0.787705230723463\\
431 1.06069198457332\\
432 1.18349694580518\\
433 1.17550829522121\\
434 1.16751964463726\\
435 1.1595309940533\\
436 1.15154234346933\\
437 1.14355369288538\\
438 1.13556504230142\\
439 1.12757639171746\\
440 1.1195877411335\\
441 1.11159909054954\\
442 1.10361043996558\\
443 1.09562178938162\\
444 1.08763313879766\\
445 1.0796444882137\\
446 1.07165583762974\\
447 1.06366718704578\\
448 1.05567853646182\\
449 1.04768988587786\\
450 1.0397012352939\\
451 1.03171258470994\\
452 1.02372393412598\\
453 1.01573528354202\\
454 1.00774663295806\\
455 0.859618679731686\\
456 0.700373919741125\\
457 0.470277595394889\\
458 0.0443466099477412\\
459 0.137809879201257\\
460 0.204030227438765\\
461 0.231683564850098\\
462 0.240547401269679\\
463 0.338400595254361\\
464 0.403637683416094\\
465 0.422375671680238\\
466 0.396179356910801\\
467 0.385284806574502\\
468 0.406706223511096\\
469 0.411790936144918\\
470 0.193811540955921\\
471 0.306892924056071\\
472 0.37525440377661\\
473 0.283923183485231\\
474 0.229298893239304\\
475 0.106279913639411\\
476 0.0964045691722652\\
477 0.0781920938424905\\
478 0.145835635108699\\
479 0.23778965255037\\
480 0.202330143141266\\
481 0.170345442909976\\
482 0.16861514966421\\
483 0.198116649685219\\
484 0.18238507145799\\
485 0.225032853095413\\
486 0.31640465663228\\
487 0.366341509182874\\
488 0.325528114538302\\
489 0.305031064503463\\
490 0.260896190933973\\
491 0.239076990307085\\
492 0.195764593593542\\
493 0.200102042912673\\
494 0.213247754762397\\
495 0.347814487819526\\
496 0.263488466491903\\
497 0.187998938970359\\
498 0.265995810866794\\
499 0.507835356524386\\
500 0.464521252434553\\
501 0.499407627065111\\
502 0.559105769884303\\
503 0.585609290076664\\
504 0.411178007508424\\
505 0.266689716093913\\
506 0.305157620053317\\
507 0.300551645648011\\
508 0.295945671242702\\
509 0.309413209094835\\
510 0.309568637187637\\
511 0.30972406528044\\
512 0.309879493373242\\
513 0.310034921466044\\
514 0.310190349558846\\
515 0.310345777651648\\
516 0.310501205744451\\
517 0.310656633837253\\
518 0.310812061930055\\
519 0.310967490022857\\
520 0.311122918115659\\
521 0.311278346208461\\
522 0.311433774301264\\
523 0.311589202394066\\
524 0.311744630486868\\
525 0.31190005857967\\
526 0.312055486672472\\
527 0.312210914765274\\
528 0.312366342858077\\
529 0.312521770950879\\
530 0.312677199043681\\
531 0.312832627136483\\
532 0.312988055229285\\
533 0.313143483322087\\
534 0.351943232790952\\
535 0.50316444244725\\
536 0.546644719573752\\
537 0.497500047955092\\
538 0.474083639175967\\
539 0.469055686630424\\
540 0.458365673445249\\
541 0.540800361686568\\
542 0.640008409233217\\
543 0.674559163287379\\
544 0.6084594866556\\
545 0.560164091403123\\
546 0.456440497227447\\
547 0.501618115763769\\
548 0.504534315950149\\
549 0.661261789214055\\
550 0.596804179543835\\
551 0.551542024383108\\
552 0.415890142964278\\
553 0.275368454428828\\
554 0.260592424256854\\
555 0.411327657328889\\
556 0.870620230959382\\
557 1.12215195046847\\
558 0.490467924836911\\
559 0.344312894999991\\
560 0.308123126904508\\
561 0.37221476130794\\
562 0.653458638923895\\
563 1.15857567405259\\
564 1.24882962091985\\
565 1.10026943004727\\
566 1.02387137230677\\
567 1.16029979909791\\
568 1.11023379084625\\
569 0.820714647055862\\
570 0.728855382506047\\
571 0.831574967626845\\
572 1.00801066077217\\
573 1.2448819872749\\
574 0.387612291621228\\
575 0.268118769091398\\
576 0.302885038080599\\
577 0.306836676047635\\
578 0.0931743976371783\\
579 0.180381012731648\\
580 0.310040940191804\\
581 0.307937139312865\\
582 0.307593317082174\\
583 0.307249494851482\\
584 0.306905672620791\\
585 0.3065618503901\\
586 0.306218028159408\\
587 0.305874205928717\\
588 0.305530383698026\\
589 0.305186561467335\\
590 0.304842739236643\\
591 0.304498917005952\\
592 0.304155094775261\\
593 0.303811272544569\\
594 0.303467450313878\\
595 0.303123628083187\\
596 0.302779805852496\\
597 0.302435983621804\\
598 0.302092161391113\\
599 0.301748339160422\\
600 0.30140451692973\\
601 0.301060694699039\\
602 0.256639619665468\\
603 0.249060025511405\\
604 0.235455981085893\\
605 0.12253051583902\\
606 0.157920421490981\\
607 0.186482554276392\\
608 0.193682086115454\\
609 0.200881617954515\\
610 0.204565368227019\\
611 0.207880505627594\\
612 0.242777180570531\\
613 0.193220983246427\\
614 0.189435874545521\\
615 0.0821276374929219\\
616 0.0838803750188513\\
617 0.135932829531935\\
618 0.180979203846406\\
619 0.130677003548264\\
620 0.177928784554347\\
621 0.231167198841538\\
622 0.176829044053726\\
623 0.155368120956646\\
624 0.269954759989914\\
625 0.316728062041472\\
626 0.326605266544081\\
627 0.327381605722724\\
628 0.313846764846853\\
629 0.255354050608453\\
630 0.274581106790452\\
631 0.286467513078632\\
632 0.352399844716199\\
633 0.404498463833684\\
634 0.444349739239712\\
635 0.430320808956354\\
636 0.389888401874792\\
637 0.203178801752923\\
638 0.275373704751226\\
639 0.371987790333855\\
640 0.289922020465667\\
641 0.326752749602076\\
642 0.531589465035964\\
643 0.501275915276576\\
644 0.226739336994872\\
645 0.229202794654555\\
646 0.247605062779919\\
647 0.250801594362938\\
648 0.347793497125595\\
649 0.197276137706483\\
650 0.22956683647629\\
651 0.47196960292594\\
652 0.366715160424131\\
653 0.337786595722564\\
654 0.340453604673937\\
655 0.343120613625308\\
656 0.345787622576679\\
657 0.348454631528052\\
658 0.351121640479424\\
659 0.353788649430795\\
660 0.356455658382166\\
661 0.359122667333539\\
662 0.36178967628491\\
663 0.364456685236281\\
664 0.367123694187652\\
665 0.369790703139025\\
666 0.372457712090397\\
667 0.375124721041768\\
668 0.377791729993141\\
669 0.380458738944512\\
670 0.383125747895883\\
671 0.385792756847254\\
672 0.388459765798627\\
673 0.391126774749998\\
674 0.39379378370137\\
675 0.396460792652743\\
676 0.399229515559883\\
677 0.419460195679859\\
678 0.440547351452646\\
679 0.419830128441014\\
680 0.47842529485995\\
681 0.52434583373015\\
682 0.542035247490701\\
683 0.58059188360195\\
684 0.549993406843702\\
685 0.503353679502633\\
686 0.520277029154098\\
687 0.402752430326425\\
688 0.337750179956546\\
689 0.249547949581286\\
690 0.203191097118529\\
691 0.425478267881327\\
692 0.581114845269121\\
693 0.885852562141737\\
694 0.777661346349539\\
695 0.406228058723441\\
696 0.301384349337235\\
697 0.194853728675045\\
698 0.16320257340712\\
699 0.402801406747106\\
700 0.749974699615523\\
701 0.925797164487001\\
702 0.911406390395424\\
703 1.03856509182799\\
704 0.566170337445383\\
705 0.79387180412277\\
706 0.787414954473041\\
707 0.793975488760531\\
708 0.782830757913221\\
709 0.760907859803832\\
710 0.738984961694443\\
711 0.717062063585038\\
712 0.6964815483678\\
713 0.628370264540688\\
714 0.747882302148477\\
715 0.986832981006019\\
716 1.0565213052818\\
717 1.14631769693449\\
718 0.985627741699727\\
719 0.364887685615936\\
720 0.512219840636466\\
721 0.488942548644851\\
722 0.376750643024103\\
723 0.334347597351422\\
724 0.440458916864602\\
725 0.554316134897994\\
726 0.506515109767618\\
727 0.393615447212912\\
728 0.362295124740135\\
729 0.3584971903007\\
730 0.354699255861268\\
731 0.350901321421836\\
732 0.347103386982401\\
733 0.343305452542969\\
734 0.339507518103537\\
735 0.335709583664104\\
736 0.331911649224669\\
737 0.328113714785237\\
738 0.324315780345805\\
739 0.320517845906373\\
740 0.316719911466938\\
741 0.312921977027505\\
742 0.309124042588073\\
743 0.305326108148638\\
744 0.301528173709206\\
745 0.297730239269774\\
746 0.293932304830342\\
747 0.290134370390907\\
748 0.286336435951474\\
749 0.282538501512042\\
750 0.264842870436365\\
751 0.211005672399023\\
752 0.142008437511614\\
753 0.142897657637007\\
754 0.199797989740186\\
755 0.237159644381204\\
756 0.207589734203596\\
757 0.102793223620854\\
758 0.122673871058344\\
759 0.18836103408237\\
760 0.199568557959804\\
761 0.186488342833612\\
762 0.161926645402794\\
763 0.0870293338922727\\
764 0.0572599796149279\\
765 0.152816338535635\\
766 0.175518051085788\\
767 0.182814349307231\\
768 0.110240832466058\\
769 0.0963908512591913\\
770 0.14449887527371\\
771 0.235141108525634\\
772 0.253777010738637\\
773 0.260308561266511\\
774 0.288631965076596\\
775 0.283895901461631\\
776 0.286595986909989\\
777 0.292811340132054\\
778 0.298358791273993\\
779 0.316179410475228\\
780 0.297176429052106\\
781 0.31864334014309\\
782 0.329675611388073\\
783 0.262904640245445\\
784 0.198552278880533\\
785 0.110416158965931\\
786 0.102087384451635\\
787 0.237689195741231\\
788 0.368428839802749\\
789 0.271801505499627\\
790 0.294782618744043\\
791 0.563011174972378\\
792 0.507737766306192\\
793 0.306314042197305\\
794 0.275096887725373\\
795 0.176186963092207\\
796 0.272726538136066\\
797 0.423815392899555\\
798 0.435265339353562\\
799 0.439251209336582\\
800 0.443237079319604\\
801 0.447222949302624\\
802 0.451208819285644\\
803 0.455194689268664\\
804 0.459180559251687\\
805 0.463166429234707\\
806 0.467152299217726\\
807 0.471138169200749\\
808 0.475124039183769\\
809 0.479109909166789\\
810 0.483095779149809\\
811 0.487081649132832\\
812 0.491067519115851\\
813 0.495053389098871\\
814 0.499039259081891\\
815 0.503025129064914\\
816 0.507010999047934\\
817 0.510996869030954\\
818 0.514982739013976\\
819 0.518968608996996\\
820 0.522954478980016\\
821 0.526940348963036\\
822 0.530926218946059\\
823 0.461559164607039\\
824 0.479007801857828\\
825 0.378197257090277\\
826 0.330327726223291\\
827 0.353640989610037\\
828 0.308230840282157\\
829 0.330388313330553\\
830 0.375867178921023\\
831 0.406047643653093\\
832 0.416097748384272\\
833 0.450502741600686\\
834 0.468870816380491\\
835 0.485882139616237\\
836 0.525856641702649\\
837 0.556045886472842\\
838 0.535603360417665\\
839 0.629786071217595\\
840 0.544055871577955\\
841 0.439503273012405\\
842 0.368646434586001\\
843 0.305441921427465\\
844 0.241401539990402\\
845 0.155648202820925\\
846 0.126851843550177\\
847 0.142152966184794\\
848 0.154732957377287\\
849 0.27222937244473\\
850 0.494400478793157\\
851 0.712646597383477\\
852 0.754449901698927\\
853 0.640716056250783\\
854 0.573054140965587\\
855 0.381829836072261\\
856 0.274882747381405\\
857 0.173595858505963\\
858 0.159931178887658\\
859 0.226245243229196\\
860 0.379950318791074\\
861 0.471750016499906\\
862 0.627035254908594\\
863 0.771777392120453\\
864 0.821734259942306\\
865 0.810636994536167\\
866 0.673294392497529\\
867 0.571456694829784\\
868 0.582240367234136\\
869 0.471329424060504\\
870 0.530512250211488\\
871 0.588910153715277\\
872 0.638256358595504\\
873 0.88245388150277\\
874 0.931537140769452\\
875 0.918373511256605\\
876 0.756282577975687\\
877 0.284944586356999\\
878 0.111947652018117\\
879 0.133267178868349\\
880 0.125917261678761\\
881 0.130743699771138\\
882 0.137729822475451\\
883 0.144715945179759\\
884 0.151702067884066\\
885 0.158688190588374\\
886 0.165674313292687\\
887 0.172660435996994\\
888 0.179646558701302\\
889 0.18663268140561\\
890 0.193618804109922\\
891 0.20060492681423\\
892 0.207591049518538\\
893 0.21457717222285\\
894 0.221563294927158\\
895 0.228549417631466\\
896 0.235535540335774\\
897 0.242521663040086\\
898 0.249507785744394\\
899 0.256493908448702\\
900 0.263480031153014\\
901 0.270466153857322\\
902 0.27745227656163\\
903 0.315544293959409\\
904 0.376059881406947\\
905 0.373809550437492\\
906 0.368581174207881\\
907 0.363352797978266\\
908 0.378191757254714\\
909 0.364429312410893\\
910 0.341955483484641\\
911 0.316911625307434\\
912 0.279632168172909\\
913 0.256702784860458\\
914 0.223129375588976\\
915 0.272498706355382\\
916 0.255609953506594\\
917 0.177661930478477\\
918 0.173148626017587\\
919 0.1161377514061\\
920 0.16168851045378\\
921 0.114724782329119\\
922 0.0471741244231732\\
923 0.0891698947843067\\
924 0.0950510983641213\\
925 0.108424462428846\\
926 0.0938238644832193\\
927 0.0471814320280151\\
928 0.094508484732654\\
929 0.0913125238380743\\
930 0.0951413539386766\\
931 0.0726189012403012\\
932 0.064387683879592\\
933 0.171344982468373\\
934 0.222603630078509\\
935 0.215347763757412\\
936 0.247076381959494\\
937 0.258360522143196\\
938 0.256274110323274\\
939 0.246835210211295\\
940 0.201123225624816\\
941 0.221051503894657\\
942 0.235378752621458\\
943 0.234014810011584\\
944 0.391927709387444\\
945 0.415517802870123\\
946 0.378232267417721\\
947 0.36039174194791\\
948 0.21794916075657\\
949 0.213816349633388\\
950 0.45788506228382\\
951 0.484083864621108\\
952 0.393011894382596\\
953 0.62862600812714\\
954 0.816444280223296\\
955 0.552052293727547\\
956 0.382014939185702\\
957 0.278936717517055\\
958 0.219393317453587\\
959 0.534019955580901\\
960 0.52796409208806\\
961 0.346713665216273\\
962 0.320772110404265\\
963 0.321803824467164\\
964 0.322835538530063\\
965 0.323867252592963\\
966 0.324898966655862\\
967 0.325930680718761\\
968 0.326962394781661\\
969 0.32799410884456\\
970 0.329025822907459\\
971 0.330057536970358\\
972 0.331089251033257\\
973 0.332120965096156\\
974 0.333152679159055\\
975 0.334184393221955\\
976 0.335216107284854\\
977 0.336247821347753\\
978 0.337279535410652\\
979 0.338311249473552\\
980 0.339342963536451\\
981 0.34037467759935\\
982 0.341406391662249\\
983 0.33952851380053\\
984 0.346434564068906\\
985 0.363924863660704\\
986 0.404060385266917\\
987 0.437575787514899\\
988 0.418367481163777\\
989 0.513599179808108\\
990 0.54444241511737\\
991 0.523899000282808\\
992 0.520325670415668\\
993 0.562063176792688\\
994 0.575299585322052\\
995 0.693658608443137\\
996 0.624613597786029\\
997 0.595237925288579\\
998 0.491832773315122\\
999 0.427044839012323\\
1000 0.339832328676591\\
1001 0.342458113754092\\
};
\addlegendentry{ungefiltert};

\addplot [
color=ohmgreen,
solid
]
table[row sep=crcr]{
1 0.350018661833295\\
2 0.350186853798652\\
3 0.350353138328617\\
4 0.350517877742661\\
5 0.35068152144753\\
6 0.350844606807404\\
7 0.351007759800017\\
8 0.35117169544851\\
9 0.351337218021128\\
10 0.351505220993894\\
11 0.35167668677369\\
12 0.351852686180712\\
13 0.352034377691886\\
14 0.352223006452171\\
15 0.352419903065399\\
16 0.352626482174603\\
17 0.352844240833177\\
18 0.353074756661216\\
19 0.353319685783473\\
20 0.353580760554926\\
21 0.353859787088762\\
22 0.354158642603413\\
23 0.354479272600399\\
24 0.354823687877081\\
25 0.355193961371621\\
26 0.355592224833357\\
27 0.356020665311344\\
28 0.356481521455852\\
29 0.356977079630223\\
30 0.357509669832336\\
31 0.358081661426556\\
32 0.358695458689097\\
33 0.359353496172084\\
34 0.360058233893125\\
35 0.360812152357633\\
36 0.361617747420896\\
37 0.362477524997226\\
38 0.36339399562448\\
39 0.364369668893665\\
40 0.365407047754753\\
41 0.366508622711396\\
42 0.367676865918898\\
43 0.368914225201573\\
44 0.370223118007125\\
45 0.371605925316583\\
46 0.373064985528268\\
47 0.374602588333602\\
48 0.37622096860178\\
49 0.377922300289409\\
50 0.37970869039027\\
51 0.381582172939133\\
52 0.383544703082852\\
53 0.385598151232063\\
54 0.387744297307871\\
55 0.389984825098748\\
56 0.392321316743031\\
57 0.394755247351875\\
58 0.397287979787133\\
59 0.399920759608738\\
60 0.402654710206277\\
61 0.405490828129224\\
62 0.408429978629595\\
63 0.411472891429905\\
64 0.414620156728399\\
65 0.417872221452583\\
66 0.421229385771215\\
67 0.424691799873961\\
68 0.428259461027042\\
69 0.431932210912301\\
70 0.435709733256227\\
71 0.43959155175459\\
72 0.443577028297508\\
73 0.447665361498885\\
74 0.451855585533363\\
75 0.456146569283084\\
76 0.460537015795805\\
77 0.465025462055085\\
78 0.469610279062582\\
79 0.474289672231808\\
80 0.47906168209236\\
81 0.483924185303305\\
82 0.488874895974039\\
83 0.493911367290045\\
84 0.499030993440344\\
85 0.50423101184377\\
86 0.509508505672886\\
87 0.514860406675851\\
88 0.520283498296284\\
89 0.525774419088852\\
90 0.531329666425727\\
91 0.536945600486907\\
92 0.542618448525237\\
93 0.54834430939472\\
94 0.554119158329481\\
95 0.55993885196092\\
96 0.565799133561012\\
97 0.57169563849948\\
98 0.577623899902085\\
99 0.583579354497583\\
100 0.5895573486418\\
101 0.595553144507595\\
102 0.60156192642867\\
103 0.607578807384071\\
104 0.613598835609942\\
105 0.619617001325814\\
106 0.62562824356329\\
107 0.631627457084832\\
108 0.637609499379207\\
109 0.643569197718746\\
110 0.649501356262361\\
111 0.655400763187739\\
112 0.661262197835849\\
113 0.667080437850777\\
114 0.672850266297962\\
115 0.678566478744928\\
116 0.684223890291298\\
117 0.689817342539786\\
118 0.695341710505871\\
119 0.70079190946875\\
120 0.70616290176733\\
121 0.711449703541717\\
122 0.716647391414031\\
123 0.721751109095457\\
124 0.726756073902138\\
125 0.731657583161782\\
126 0.736451020495287\\
127 0.741131861962144\\
128 0.745695682064091\\
129 0.750138159607358\\
130 0.754455083427757\\
131 0.758642357983257\\
132 0.762696008816096\\
133 0.766612187882199\\
134 0.77038717874109\\
135 0.774017401595684\\
136 0.777499418169174\\
137 0.780829936405844\\
138 0.784005814983137\\
139 0.787024067623156\\
140 0.789881867192693\\
141 0.792576549581809\\
142 0.795105617351839\\
143 0.797466743144606\\
144 0.799657772845461\\
145 0.801676728493585\\
146 0.803521810933831\\
147 0.805191402205134\\
148 0.806684067661313\\
149 0.807998557820832\\
150 0.809133809942794\\
151 0.810088949327173\\
152 0.810863290337931\\
153 0.811456337148363\\
154 0.81186778420863\\
155 0.812097516436079\\
156 0.812145609129696\\
157 0.812012327611352\\
158 0.811698126598417\\
159 0.811203649314533\\
160 0.810529726346113\\
161 0.809677374250407\\
162 0.808647793917096\\
163 0.807442368681405\\
164 0.806062662184659\\
165 0.804510415978527\\
166 0.802787546871636\\
167 0.800896144020867\\
168 0.798838465772699\\
169 0.796616936260985\\
170 0.794234141766849\\
171 0.791692826845667\\
172 0.788995890225849\\
173 0.786146380484023\\
174 0.78314749150092\\
175 0.780002557702406\\
176 0.776715049090921\\
177 0.773288566074073\\
178 0.769726834098755\\
179 0.766033698100291\\
180 0.762213116776774\\
181 0.758269156699295\\
182 0.754205986269831\\
183 0.750027869540239\\
184 0.745739159907623\\
185 0.74134429370285\\
186 0.736847783689672\\
187 0.732254212491988\\
188 0.727568225966212\\
189 0.722794526534897\\
190 0.717937866496669\\
191 0.713003041326334\\
192 0.707994882978067\\
193 0.702918253204093\\
194 0.697778036901329\\
195 0.692579135498542\\
196 0.68732646039633\\
197 0.682024926471633\\
198 0.676679445658241\\
199 0.671294920615041\\
200 0.665876238494231\\
201 0.660428264822044\\
202 0.654955837504951\\
203 0.649463760974649\\
204 0.643956800485056\\
205 0.638439676573654\\
206 0.632917059698123\\
207 0.627393565057681\\
208 0.621873747607109\\
209 0.616362097270068\\
210 0.610863034357065\\
211 0.605380905192255\\
212 0.599919977952202\\
213 0.59448443871872\\
214 0.589078387747015\\
215 0.583705835949492\\
216 0.578370701594867\\
217 0.573076807221481\\
218 0.567827876763131\\
219 0.562627532885127\\
220 0.557479294527778\\
221 0.552386574654065\\
222 0.547352678197828\\
223 0.542380800208432\\
224 0.537474024187582\\
225 0.532635320613642\\
226 0.527867545648565\\
227 0.52317344002206\\
228 0.518555628086686\\
229 0.514016617036383\\
230 0.509558796280074\\
231 0.505184436962075\\
232 0.50089569162179\\
233 0.496694593986167\\
234 0.492583058889623\\
235 0.488562882317142\\
236 0.484635741566884\\
237 0.480803195528643\\
238 0.47706668507452\\
239 0.473427533558965\\
240 0.469886947426865\\
241 0.466446016929809\\
242 0.463105716951353\\
243 0.459866907941833\\
244 0.456730336962653\\
245 0.453696638839533\\
246 0.450766337424109\\
247 0.447939846963534\\
248 0.445217473577556\\
249 0.442599416841443\\
250 0.440085771471209\\
251 0.437676529105701\\
252 0.435371580179126\\
253 0.433170715878107\\
254 0.431073630179858\\
255 0.429079921971499\\
256 0.427189097253166\\
257 0.425400571426913\\
258 0.423713671669336\\
259 0.422127639380165\\
260 0.420641632695643\\
261 0.419254729056455\\
262 0.417965927825372\\
263 0.416774152957332\\
264 0.415678255731147\\
265 0.414677017554634\\
266 0.413769152852957\\
267 0.412953312045497\\
268 0.412228084612263\\
269 0.411592002247594\\
270 0.411043542095816\\
271 0.410581130060483\\
272 0.410203144177656\\
273 0.40990791804498\\
274 0.409693744299403\\
275 0.409558878133735\\
276 0.409501540836292\\
277 0.409519923332404\\
278 0.409612189703799\\
279 0.409776480662027\\
280 0.410010916953809\\
281 0.410313602678849\\
282 0.410682628503282\\
283 0.411116074754514\\
284 0.411612014385569\\
285 0.412168515799194\\
286 0.412783645524048\\
287 0.41345547073708\\
288 0.414182061627945\\
289 0.414961493602816\\
290 0.415791849326391\\
291 0.416671220602147\\
292 0.417597710092076\\
293 0.418569432878154\\
294 0.419584517868735\\
295 0.420641109053887\\
296 0.421737366614408\\
297 0.422871467889906\\
298 0.424041608211873\\
299 0.42524600160815\\
300 0.42648288138547\\
301 0.427750500596738\\
302 0.429047132399388\\
303 0.430371070311113\\
304 0.431720628369765\\
305 0.433094141205398\\
306 0.43448996403372\\
307 0.43590647258108\\
308 0.437342062951084\\
309 0.438795151442118\\
310 0.440264174323666\\
311 0.441747587577827\\
312 0.443243866610941\\
313 0.444751505938863\\
314 0.446269018848456\\
315 0.447794937037376\\
316 0.449327810233967\\
317 0.450866205798672\\
318 0.452408708307863\\
319 0.453953919120645\\
320 0.455500455928951\\
321 0.457046952291123\\
322 0.458592057149526\\
323 0.460134434333312\\
324 0.461672762048197\\
325 0.46320573235549\\
326 0.464732050642983\\
327 0.466250435090772\\
328 0.467759616135774\\
329 0.469258335939345\\
330 0.470745347863023\\
331 0.47221941595778\\
332 0.473679314472433\\
333 0.475123827386971\\
334 0.47655174797643\\
335 0.477961878410677\\
336 0.479353029394971\\
337 0.480724019855721\\
338 0.48207367667553\\
339 0.483400834481638\\
340 0.484704335492131\\
341 0.485983029424412\\
342 0.487235773470082\\
343 0.488461432338997\\
344 0.489658878373006\\
345 0.490826991727704\\
346 0.4919646606196\\
347 0.493070781636584\\
348 0.49414426011075\\
349 0.495184010553931\\
350 0.496188957156912\\
351 0.497158034352496\\
352 0.498090187440396\\
353 0.498984373269755\\
354 0.499839560973865\\
355 0.500654732751469\\
356 0.501428884689204\\
357 0.502161027620001\\
358 0.502850188012486\\
359 0.50349540888668\\
360 0.504095750751485\\
361 0.504650292559698\\
362 0.505158132676441\\
363 0.50561838985716\\
364 0.506030204231466\\
365 0.506392738289315\\
366 0.506705177866194\\
367 0.506966733124115\\
368 0.507176639525425\\
369 0.507334158796553\\
370 0.507438579878992\\
371 0.507489219864946\\
372 0.507485424915207\\
373 0.507426571156971\\
374 0.507312065559426\\
375 0.507141346785063\\
376 0.506913886014794\\
377 0.506629187745105\\
378 0.506286790555677\\
379 0.505886267846355\\
380 0.505427228543192\\
381 0.504909317774471\\
382 0.504332217518815\\
383 0.503695647228079\\
384 0.502999364427728\\
385 0.502243165296898\\
386 0.501426885229732\\
387 0.500550399379042\\
388 0.499613623183036\\
389 0.498616512875984\\
390 0.49755906598401\\
391 0.496441321807437\\
392 0.495263361890881\\
393 0.494025310481966\\
394 0.492727334979179\\
395 0.491369646369464\\
396 0.489952499656454\\
397 0.488476194280795\\
398 0.486941074534275\\
399 0.485347529968836\\
400 0.48369599579989\\
401 0.481986953301143\\
402 0.480220930185906\\
403 0.478398500968115\\
404 0.476520287295235\\
405 0.474586958245187\\
406 0.472599230580287\\
407 0.470557868953003\\
408 0.46846368606103\\
409 0.466317542752601\\
410 0.464120348086366\\
411 0.461873059352769\\
412 0.45957668206528\\
413 0.457232269929875\\
414 0.45484092479966\\
415 0.452403796618155\\
416 0.449922083349866\\
417 0.447397030891479\\
418 0.444829932953502\\
419 0.442222130901693\\
420 0.439575013549958\\
421 0.436890016900576\\
422 0.43416862383247\\
423 0.4314123637428\\
424 0.42862281214992\\
425 0.425801590266186\\
426 0.422950364547447\\
427 0.420070846223233\\
428 0.417164790808209\\
429 0.414233997592346\\
430 0.411280309105772\\
431 0.408305610554372\\
432 0.405311829223197\\
433 0.402300933846192\\
434 0.399274933942799\\
435 0.396235879124274\\
436 0.39318585837376\\
437 0.390126999303978\\
438 0.3870614673955\\
439 0.383991465217463\\
440 0.38091923163161\\
441 0.37784704097968\\
442 0.374777202253434\\
443 0.371712058245929\\
444 0.368653984682128\\
445 0.365605389326461\\
446 0.362568711064592\\
447 0.359546418956325\\
448 0.356541011256393\\
449 0.353555014399677\\
450 0.350590981947344\\
451 0.347651493490321\\
452 0.344739153506542\\
453 0.341856590168482\\
454 0.339006454097541\\
455 0.336191417062024\\
456 0.333414170615302\\
457 0.330677424670133\\
458 0.327983906003231\\
459 0.32533635668105\\
460 0.322737532393703\\
461 0.320190200680631\\
462 0.317697139030897\\
463 0.315261132843135\\
464 0.312884973233871\\
465 0.310571454686877\\
466 0.308323372539958\\
467 0.306143520308971\\
468 0.304034686851503\\
469 0.301999653374332\\
470 0.300041190290007\\
471 0.298162053928632\\
472 0.29636498311125\\
473 0.294652695591282\\
474 0.293027884370897\\
475 0.291493213900059\\
476 0.290051316166471\\
477 0.288704786684234\\
478 0.28745618038824\\
479 0.286308007441019\\
480 0.285262728959473\\
481 0.284322752670447\\
482 0.283490428505868\\
483 0.28276804414936\\
484 0.282157820546716\\
485 0.281661907392821\\
486 0.281282378607906\\
487 0.28102122781663\\
488 0.28088036384434\\
489 0.280861606245913\\
490 0.280966680883282\\
491 0.281197215567594\\
492 0.281554735781009\\
493 0.282040660491814\\
494 0.282656298075188\\
495 0.283402842350881\\
496 0.284281368748523\\
497 0.285292830611305\\
498 0.286438055648955\\
499 0.287717742550677\\
500 0.289132457768281\\
501 0.290682632479998\\
502 0.292368559746629\\
503 0.294190391872657\\
504 0.296148137985117\\
505 0.298241661842331\\
506 0.300470679882925\\
507 0.302834759522722\\
508 0.305333317703927\\
509 0.307965619698623\\
510 0.310730778167401\\
511 0.313627752473403\\
512 0.316655348251724\\
513 0.319812217233758\\
514 0.323096857325715\\
515 0.326507612940078\\
516 0.330042675578388\\
517 0.333700084663275\\
518 0.337477728617265\\
519 0.341373346185476\\
520 0.345384527998925\\
521 0.349508718374771\\
522 0.353743217349466\\
523 0.358085182940407\\
524 0.362531633631354\\
525 0.367079451076531\\
526 0.371725383018007\\
527 0.376466046410665\\
528 0.381297930748753\\
529 0.386217401587764\\
530 0.391220704255089\\
531 0.396303967742687\\
532 0.401463208774743\\
533 0.406694336043109\\
534 0.411993154603093\\
535 0.417355370422088\\
536 0.422776595073786\\
537 0.428252350571505\\
538 0.433778074335443\\
539 0.439349124289462\\
540 0.44496078408298\\
541 0.450608268432618\\
542 0.456286728577374\\
543 0.461991257840822\\
544 0.467716897294237\\
545 0.473458641515325\\
546 0.479211444437516\\
547 0.484970225283858\\
548 0.490729874577924\\
549 0.49648526022244\\
550 0.502231233635503\\
551 0.507962635934553\\
552 0.513674304158889\\
553 0.519361077521621\\
554 0.525017803680694\\
555 0.530639345016474\\
556 0.536220584901511\\
557 0.541756433948188\\
558 0.547241836223173\\
559 0.55267177542238\\
560 0.558041281002699\\
561 0.563345434264793\\
562 0.568579374376647\\
563 0.573738304324767\\
564 0.578817496781284\\
565 0.583812299880273\\
566 0.588718142902937\\
567 0.593530541875468\\
568 0.598245105084006\\
569 0.602857538509303\\
570 0.607363651181442\\
571 0.611759360452241\\
572 0.616040697179583\\
573 0.620203810815155\\
574 0.624244974386608\\
575 0.628160589366101\\
576 0.631947190416331\\
577 0.635601450000895\\
578 0.639120182840417\\
579 0.642500350192959\\
580 0.645739063936947\\
581 0.648833590435979\\
582 0.651781354166928\\
583 0.654579941095752\\
584 0.65722710178855\\
585 0.659720754247889\\
586 0.662058986466258\\
587 0.664240058690021\\
588 0.666262405388633\\
589 0.668124636925236\\
590 0.669825540925978\\
591 0.671364083346578\\
592 0.672739409235774\\
593 0.673950843196292\\
594 0.674997889544977\\
595 0.675880232174584\\
596 0.676597734120599\\
597 0.677150436837218\\
598 0.677538559187353\\
599 0.677762496152176\\
600 0.677822817266347\\
601 0.677720264785621\\
602 0.677455751594029\\
603 0.677030358858237\\
604 0.676445333436889\\
605 0.675702085052759\\
606 0.674802183235335\\
607 0.673747354041124\\
608 0.672539476558655\\
609 0.671180579205208\\
610 0.669672835822973\\
611 0.668018561583158\\
612 0.666220208707354\\
613 0.664280362016057\\
614 0.662201734314686\\
615 0.659987161627747\\
616 0.657639598291634\\
617 0.655162111916016\\
618 0.652557878223022\\
619 0.649830175773156\\
620 0.646982380587252\\
621 0.644017960674476\\
622 0.640940470477063\\
623 0.637753545242882\\
624 0.634460895337322\\
625 0.631066300506026\\
626 0.627573604100202\\
627 0.623986707276739\\
628 0.620309563186149\\
629 0.616546171161642\\
630 0.612700570922372\\
631 0.608776836802966\\
632 0.60477907202042\\
633 0.600711402988692\\
634 0.596577973691114\\
635 0.592382940121085\\
636 0.588130464801949\\
637 0.58382471139702\\
638 0.579469839419881\\
639 0.575069999053358\\
640 0.570629326083547\\
641 0.566151936954048\\
642 0.561641923945254\\
643 0.557103350483785\\
644 0.552540246587714\\
645 0.547956604453722\\
646 0.54335637419175\\
647 0.538743459710686\\
648 0.53412171475636\\
649 0.529494939101932\\
650 0.524866874890666\\
651 0.520241203131232\\
652 0.515621540345717\\
653 0.511011435370898\\
654 0.506414366314176\\
655 0.501833737666104\\
656 0.497272877571124\\
657 0.492735035257398\\
658 0.488223378626047\\
659 0.483740991999808\\
660 0.479290874030867\\
661 0.474875935767387\\
662 0.470498998878064\\
663 0.466162794033793\\
664 0.461869959445335\\
665 0.457623039555677\\
666 0.453424483885559\\
667 0.449276646030491\\
668 0.445181782807355\\
669 0.441142053548575\\
670 0.437159519541607\\
671 0.433236143611399\\
672 0.429373789843282\\
673 0.425574223443628\\
674 0.42183911073547\\
675 0.418170019286146\\
676 0.414568418163912\\
677 0.411035678320375\\
678 0.407573073095486\\
679 0.404181778841877\\
680 0.400862875665382\\
681 0.397617348278735\\
682 0.394446086965625\\
683 0.391349888652748\\
684 0.388329458087945\\
685 0.385385409123043\\
686 0.382518266099924\\
687 0.379728465337902\\
688 0.37701635671945\\
689 0.374382205369993\\
690 0.371826193425807\\
691 0.369348421882278\\
692 0.366948912513506\\
693 0.364627609854615\\
694 0.362384383240621\\
695 0.360219028899617\\
696 0.358131272101156\\
697 0.356120769360564\\
698 0.354187110696545\\
699 0.352329821934894\\
700 0.35054836704832\\
701 0.348842150522616\\
702 0.347210519742935\\
703 0.345652767399236\\
704 0.344168133914688\\
705 0.342755809902869\\
706 0.341414938658613\\
707 0.340144618684668\\
708 0.338943906253957\\
709 0.337811818006188\\
710 0.336747333577152\\
711 0.335749398258457\\
712 0.334816925684561\\
713 0.333948800542985\\
714 0.333143881302627\\
715 0.33240100295438\\
716 0.331718979758302\\
717 0.331096607993038\\
718 0.330532668705776\\
719 0.330025930463316\\
720 0.329575152104728\\
721 0.329179085492678\\
722 0.328836478255155\\
723 0.328546076505107\\
724 0.328306627523669\\
725 0.328116882392577\\
726 0.327975598562203\\
727 0.327881542343436\\
728 0.327833491313837\\
729 0.327830236629879\\
730 0.32787058523739\\
731 0.327953361972133\\
732 0.328077411542756\\
733 0.328241600389101\\
734 0.328444818409855\\
735 0.328685980554549\\
736 0.328964028275865\\
737 0.329277930839107\\
738 0.329626686486541\\
739 0.330009323455132\\
740 0.330424900846914\\
741 0.330872509351954\\
742 0.33135127182453\\
743 0.331860343713717\\
744 0.332398913350199\\
745 0.332966202091569\\
746 0.333561464328931\\
747 0.334183987358011\\
748 0.334833091118405\\
749 0.335508127804976\\
750 0.3362084813557\\
751 0.336933566820577\\
752 0.337682829616314\\
753 0.338455744671309\\
754 0.33925181546501\\
755 0.340070572965238\\
756 0.340911574467147\\
757 0.341774402338201\\
758 0.342658662674313\\
759 0.34356398387257\\
760 0.344490015125713\\
761 0.345436424843498\\
762 0.346402899006608\\
763 0.347389139459494\\
764 0.348394862148756\\
765 0.349419795313303\\
766 0.350463677631954\\
767 0.351526256333866\\
768 0.352607285277706\\
769 0.353706523006175\\
770 0.354823730782851\\
771 0.355958670618045\\
772 0.357111103290132\\
773 0.35828078636921\\
774 0.35946747225081\\
775 0.360670906208221\\
776 0.361890824472359\\
777 0.363126952348046\\
778 0.364379002375243\\
779 0.365646672543314\\
780 0.366929644565896\\
781 0.368227582223461\\
782 0.369540129780194\\
783 0.370866910481267\\
784 0.372207525136018\\
785 0.37356155079168\\
786 0.374928539500959\\
787 0.376308017185014\\
788 0.377699482591895\\
789 0.379102406350125\\
790 0.380516230118011\\
791 0.381940365830637\\
792 0.383374195047463\\
793 0.384817068404294\\
794 0.386268305174076\\
795 0.387727192940766\\
796 0.389192987388579\\
797 0.390664912206304\\
798 0.392142159104694\\
799 0.393623887944987\\
800 0.395109226977779\\
801 0.39659727319266\\
802 0.398087092779377\\
803 0.399577721700981\\
804 0.401068166378886\\
805 0.402557404489204\\
806 0.404044385869187\\
807 0.405528033532153\\
808 0.407007244788767\\
809 0.40848089247219\\
810 0.409947826264154\\
811 0.411406874118711\\
812 0.412856843780044\\
813 0.414296524390448\\
814 0.415724688184308\\
815 0.417140092263643\\
816 0.418541480450603\\
817 0.41992758521206\\
818 0.421297129651308\\
819 0.422648829561711\\
820 0.423981395537019\\
821 0.425293535132977\\
822 0.426583955074757\\
823 0.427851363504659\\
824 0.429094472264405\\
825 0.430311999205954\\
826 0.431502670524184\\
827 0.432665223103904\\
828 0.433798406872707\\
829 0.434900987150426\\
830 0.435971746985678\\
831 0.437009489470216\\
832 0.43801304002259\\
833 0.438981248633702\\
834 0.439912992068106\\
835 0.440807176016033\\
836 0.441662737192085\\
837 0.442478645377389\\
838 0.443253905402783\\
839 0.44398755907133\\
840 0.444678687019051\\
841 0.445326410513227\\
842 0.445929893187593\\
843 0.446488342712967\\
844 0.447001012400295\\
845 0.447467202731209\\
846 0.447886262809586\\
847 0.448257591726378\\
848 0.448580639829266\\
849 0.448854909888626\\
850 0.449079958152238\\
851 0.449255395283227\\
852 0.449380887179107\\
853 0.449456155674156\\
854 0.44948097913133\\
855 0.449455192931962\\
856 0.449378689870834\\
857 0.449251420461503\\
858 0.449073393153257\\
859 0.448844674457796\\
860 0.448565388981658\\
861 0.4482357193601\\
862 0.447855906089626\\
863 0.447426247259294\\
864 0.446947098184418\\
865 0.446418870949671\\
866 0.445842033871237\\
867 0.44521711088877\\
868 0.444544680897218\\
869 0.443825377026239\\
870 0.443059885872136\\
871 0.44224894668493\\
872 0.441393350511802\\
873 0.440493939297912\\
874 0.43955160494652\\
875 0.438567288342297\\
876 0.437541978344\\
877 0.436476710754014\\
878 0.43537256727122\\
879 0.434230674429475\\
880 0.433052202517965\\
881 0.43183836447445\\
882 0.430590414739974\\
883 0.429309648063802\\
884 0.427997398248791\\
885 0.426655036829148\\
886 0.4252839716742\\
887 0.423885645513335\\
888 0.422461534378773\\
889 0.421013145964096\\
890 0.419542017897725\\
891 0.418049715931583\\
892 0.416537832046236\\
893 0.415007982474657\\
894 0.413461805647603\\
895 0.411900960064291\\
896 0.410327122092725\\
897 0.408741983704549\\
898 0.407147250149838\\
899 0.405544637577602\\
900 0.403935870608183\\
901 0.402322679863975\\
902 0.400706799465169\\
903 0.399089964497393\\
904 0.397473908458324\\
905 0.395860360690632\\
906 0.394251043809127\\
907 0.392647671130406\\
908 0.391051944113568\\
909 0.389465549820302\\
910 0.387890158402202\\
911 0.386327420622647\\
912 0.384778965419877\\
913 0.383246397517093\\
914 0.381731295084382\\
915 0.380235207456263\\
916 0.37875965290777\\
917 0.377306116491457\\
918 0.375876047937432\\
919 0.374470859618222\\
920 0.373091924579623\\
921 0.371740574637962\\
922 0.370418098543673\\
923 0.369125740210814\\
924 0.367864697011942\\
925 0.36663611813753\\
926 0.365441103019225\\
927 0.36428069981665\\
928 0.363155903967857\\
929 0.36206765680382\\
930 0.361016844227358\\
931 0.360004295457111\\
932 0.359030781837501\\
933 0.358097015715779\\
934 0.357203649387492\\
935 0.356351274112104\\
936 0.355540419201425\\
937 0.354771551184554\\
938 0.35404507305367\\
939 0.353361323595233\\
940 0.352720576810984\\
941 0.352123041432682\\
942 0.351568860533809\\
943 0.35105811124068\\
944 0.350590804544799\\
945 0.350166885218187\\
946 0.349786231833869\\
947 0.349448656894511\\
948 0.349153907072696\\
949 0.348901663565832\\
950 0.348691542567192\\
951 0.348523095852942\\
952 0.348395811484424\\
953 0.348309114625707\\
954 0.348262368477657\\
955 0.348254875330922\\
956 0.348285877741227\\
957 0.348354559830605\\
958 0.348460048716272\\
959 0.348601416065068\\
960 0.348777679767951\\
961 0.34898780572778\\
962 0.349230709754446\\
963 0.349505259562278\\
964 0.34981027686413\\
965 0.350144539555234\\
966 0.350506783978912\\
967 0.350895707265987\\
968 0.351309969739899\\
969 0.351748197379753\\
970 0.352208984333767\\
971 0.352690895475814\\
972 0.353192468998007\\
973 0.353712219032483\\
974 0.354248638295785\\
975 0.35480020074947\\
976 0.355365364270795\\
977 0.355942573327543\\
978 0.356530261651306\\
979 0.357126854903725\\
980 0.357730773330434\\
981 0.358340434397643\\
982 0.358954255406551\\
983 0.359570656080936\\
984 0.360188061123517\\
985 0.360804902736864\\
986 0.361419623104892\\
987 0.36203067683131\\
988 0.362636533331981\\
989 0.363235679178827\\
990 0.363826620393668\\
991 0.364407884691081\\
992 0.364978023670128\\
993 0.365535614955338\\
994 0.366079264287552\\
995 0.366607607565125\\
996 0.367119312836091\\
997 0.367613082242281\\
998 0.36808765391667\\
999 0.368541803834865\\
1000 0.368974347620351\\
1001 0.369384142301124\\
};
\addlegendentry{gefiltert};

\end{axis}
\end{tikzpicture}% 
\end{frame}



\begin{frame}{Unterscheidung}
    \centering 
    \setlength\figureheight{0.45\textwidth} 
    \setlength\figurewidth{0.8\textwidth} 
    % This file was created by matlab2tikz v0.4.3.
% Copyright (c) 2008--2013, Nico Schlömer <nico.schloemer@gmail.com>
% All rights reserved.
% 
% The latest updates can be retrieved from
%   http://www.mathworks.com/matlabcentral/fileexchange/22022-matlab2tikz
% where you can also make suggestions and rate matlab2tikz.
% 
%
% defining custom colors
\definecolor{mycolor1}{rgb}{0,0.75,0.75}%
%
\begin{tikzpicture}[font=\sffamily\sansmath]

\begin{axis}[%
width=\figurewidth,
height=\figureheight,
scale only axis,
xmin=0,
xmax=1000,
xtick={0,100,200,300,400,500,600,700,800,900,1000},
xticklabels={0,1,2,3,4,5,6,7,8,9,10},
xlabel={Zeit in s},
ymin=0,
ymax=0.9,
ylabel={Beschleunigungslevel},
legend style={draw=black,fill=white,legend cell align=left}
]
\addplot [
color=blue,
solid
]
table[row sep=crcr]{
1 0.00756082744702296\\
2 0.0074815024948087\\
3 0.00740511242821824\\
4 0.00733163732769468\\
5 0.0072610554690464\\
6 0.00719334337630905\\
7 0.00712847587497403\\
8 0.00706642614552894\\
9 0.00700716577725719\\
10 0.00695066482224533\\
11 0.00689689184954831\\
12 0.00684581399946441\\
13 0.00679739703787276\\
14 0.00675160541058712\\
15 0.00670840229767976\\
16 0.00666774966772488\\
17 0.00662960833190034\\
18 0.00659393799786519\\
19 0.00656069732331066\\
20 0.00652984396907896\\
21 0.00650133465176402\\
22 0.00647512519573442\\
23 0.00645117058453198\\
24 0.00642942501159781\\
25 0.00640984193027693\\
26 0.00639237410305334\\
27 0.00637697364996324\\
28 0.00636359209612866\\
29 0.00635218041836142\\
30 0.00634268909080763\\
31 0.00633506812961964\\
32 0.00632926713664663\\
33 0.0063252353421361\\
34 0.00632292164644466\\
35 0.00632227466077226\\
36 0.00632324274695969\\
37 0.00632577405641841\\
38 0.00632981656827835\\
39 0.00633531812683229\\
40 0.00634222647832758\\
41 0.00635048930712448\\
42 0.00636005427122107\\
43 0.00637086903714121\\
44 0.00638288131418374\\
45 0.00639603888803306\\
46 0.00641028965373326\\
47 0.00642558164803215\\
48 0.00644186308110435\\
49 0.00645908236766323\\
50 0.00647718815746809\\
51 0.00649612936522703\\
52 0.00651585519989582\\
53 0.00653631519338888\\
54 0.00655745922875553\\
55 0.00657923756791555\\
56 0.00660160087905725\\
57 0.00662450026375428\\
58 0.00664788728378235\\
59 0.00667171398756326\\
60 0.00669593293615724\\
61 0.0067204972287537\\
62 0.00674536052764703\\
63 0.00677047708270624\\
64 0.00679580175535173\\
65 0.00682129004204936\\
66 0.00684689809732768\\
67 0.00687258275632059\\
68 0.00689830155683402\\
69 0.00692401276093201\\
70 0.00694967537603509\\
71 0.00697524917552124\\
72 0.00700069471881733\\
73 0.00702597337096754\\
74 0.00705104732166306\\
75 0.00707587960371654\\
76 0.00710043411096327\\
77 0.00712467561557031\\
78 0.00714856978473415\\
79 0.00717208319674689\\
80 0.00719518335641069\\
81 0.00721783870977965\\
82 0.00724001865820562\\
83 0.00726169357165633\\
84 0.00728283480126519\\
85 0.00730341469107911\\
86 0.00732340658900545\\
87 0.00734278485700836\\
88 0.00736152488063205\\
89 0.00737960307791485\\
90 0.00739699690770946\\
91 0.00741368487736513\\
92 0.00742964654968086\\
93 0.00744486254902291\\
94 0.00745931456651262\\
95 0.00747298536421892\\
96 0.00748585877831953\\
97 0.00749791972121237\\
98 0.00750915418256169\\
99 0.00751954922925892\\
100 0.00752909300427696\\
101 0.0075377747243984\\
102 0.00754558467680208\\
103 0.00755251421450046\\
104 0.00755855575063208\\
105 0.00756370275162261\\
106 0.0075679497292309\\
107 0.00757129223150363\\
108 0.00757372683267849\\
109 0.00757525112208966\\
110 0.0075758636921277\\
111 0.0075755641252936\\
112 0.00757435298037752\\
113 0.00757223177780007\\
114 0.00756920298417429\\
115 0.00756526999616502\\
116 0.00756043712372122\\
117 0.00755470957273884\\
118 0.00754809342719026\\
119 0.00754059563074848\\
120 0.00753222396793157\\
121 0.00752298704478681\\
122 0.00751289426912609\\
123 0.00750195583032663\\
124 0.00749018267872148\\
125 0.00747758650461425\\
126 0.00746417971695828\\
127 0.00744997542174017\\
128 0.00743498740009585\\
129 0.00741923008617115\\
130 0.00740271854472856\\
131 0.0073854684484999\\
132 0.00736749605528709\\
133 0.00734881818481572\\
134 0.00732945219534879\\
135 0.00730941596007019\\
136 0.0072887278432497\\
137 0.00726740667620316\\
138 0.00724547173306332\\
139 0.00722294270637831\\
140 0.00719983968255625\\
141 0.00717618311717576\\
142 0.00715199381018329\\
143 0.00712729288099912\\
144 0.00710210174355478\\
145 0.00707644208128545\\
146 0.00705033582210137\\
147 0.00702380511336295\\
148 0.00699687229688476\\
149 0.00696955988399329\\
150 0.00694189053066384\\
151 0.00691388701275832\\
152 0.00688557220137911\\
153 0.00685696903834482\\
154 0.00682810051178803\\
155 0.00679898963187586\\
156 0.00676965940666405\\
157 0.00674013281811443\\
158 0.00671043279832383\\
159 0.0066805822060137\\
160 0.00665060380330985\\
161 0.00662052023281544\\
162 0.00659035399496305\\
163 0.00656012742562628\\
164 0.0065298626739747\\
165 0.0064995816805666\\
166 0.00646930615569355\\
167 0.00643905755800638\\
168 0.00640885707345193\\
169 0.00637872559453999\\
170 0.00634868369996016\\
171 0.00631875163459013\\
172 0.00628894928996721\\
173 0.00625929618531622\\
174 0.00622981144923429\\
175 0.00620051380212787\\
176 0.00617142153948031\\
177 0.00614255251601125\\
178 0.00611392413079127\\
179 0.00608555331339585\\
180 0.00605745651119987\\
181 0.00602964967790749\\
182 0.00600214826338209\\
183 0.00597496720480466\\
184 0.00594812091917577\\
185 0.0059216232972024\\
186 0.00589548769867027\\
187 0.00586972694946296\\
188 0.00584435334042396\\
189 0.00581937862825031\\
190 0.00579481403855943\\
191 0.00577067027119219\\
192 0.00574695750772602\\
193 0.00572368542110757\\
194 0.00570086318730922\\
195 0.00567849949897117\\
196 0.00565660258107392\\
197 0.00563518020874808\\
198 0.00561423972734671\\
199 0.00559378807488196\\
200 0.00557383180687092\\
201 0.00555437712356209\\
202 0.00553542989944857\\
203 0.00551699571493417\\
204 0.00549907988999801\\
205 0.00548168751969109\\
206 0.00546482351128906\\
207 0.0054484926229183\\
208 0.00543269950346588\\
209 0.00541744873358001\\
210 0.00540274486756407\\
211 0.00538859247596566\\
212 0.00537499618866139\\
213 0.00536196073823845\\
214 0.00534949100347553\\
215 0.00533759205272806\\
216 0.00532626918702573\\
217 0.00531552798269458\\
218 0.00530537433332032\\
219 0.00529581449087493\\
220 0.00528685510583245\\
221 0.00527850326609705\\
222 0.00527076653454902\\
223 0.00526365298498494\\
224 0.00525717123619632\\
225 0.00525133048390456\\
226 0.00524614053024652\\
227 0.00524161181048899\\
228 0.00523775541665973\\
229 0.00523458311782511\\
230 0.00523210737682422\\
231 0.00523034136338616\\
232 0.00522929896369099\\
233 0.00522899478654001\\
234 0.00522944416635316\\
235 0.00523066316323659\\
236 0.00523266856038745\\
237 0.00523547785909026\\
238 0.00523910927145564\\
239 0.00524358171088205\\
240 0.00524891478008801\\
241 0.00525512875653618\\
242 0.00526224457513934\\
243 0.00527028380823137\\
244 0.0052792686428318\\
245 0.00528922185520402\\
246 0.00530016678263506\\
247 0.00531212729231299\\
248 0.00532512774719033\\
249 0.00533919296879847\\
250 0.0053543481970872\\
251 0.00537061904746648\\
252 0.00538803146530405\\
253 0.0054066116781742\\
254 0.00542638614615109\\
255 0.00544738151038405\\
256 0.00546962454008355\\
257 0.00549314207792221\\
258 0.00551796098377256\\
259 0.00554410807668601\\
260 0.00557161007504624\\
261 0.00560049353488618\\
262 0.00563078478643999\\
263 0.00566250986908445\\
264 0.00569569446485633\\
265 0.00573036383071001\\
266 0.00576654272965071\\
267 0.00580425536087621\\
268 0.00584352528907264\\
269 0.0058843753730251\\
270 0.00592682769371618\\
271 0.0059709034820964\\
272 0.00601662304671933\\
273 0.00606400570144082\\
274 0.00611306969338749\\
275 0.00616383213140291\\
276 0.00621630891518263\\
277 0.00627051466530964\\
278 0.00632646265440188\\
279 0.00638416473958168\\
280 0.00644363129647435\\
281 0.00650487115493977\\
282 0.00656789153673606\\
283 0.00663269799530928\\
284 0.00669929435789733\\
285 0.0067676826701344\\
286 0.00683786314334842\\
287 0.00690983410476142\\
288 0.00698359195081726\\
289 0.00705913110385488\\
290 0.00713644397230934\\
291 0.00721552091457069\\
292 0.00729635020658944\\
293 0.0073789180132962\\
294 0.00746320836389223\\
295 0.00754920313106881\\
296 0.00763688201423705\\
297 0.00772622252688952\\
298 0.00781719998825775\\
299 0.00790978751946704\\
300 0.00800395604440537\\
301 0.00809967429550213\\
302 0.00819690882456315\\
303 0.00829562401874518\\
304 0.00839578212167281\\
305 0.00849734325960155\\
306 0.00860026547242876\\
307 0.00870450474928404\\
308 0.00881001506842371\\
309 0.0089167484412073\\
310 0.00902465496000923\\
311 0.00913368284997689\\
312 0.00924377852457609\\
313 0.0093548866448752\\
314 0.00946695018251022\\
315 0.00957991048623627\\
316 0.00969370735191387\\
317 0.00980827909571657\\
318 0.00992356263030167\\
319 0.0100394935436815\\
320 0.0101560061805989\\
321 0.0102730337263459\\
322 0.0103905082931187\\
323 0.0105083610090911\\
324 0.0106265221103651\\
325 0.0107449210358034\\
326 0.0108634865245493\\
327 0.0109821467158758\\
328 0.0111008292509387\\
329 0.011219461375992\\
330 0.0113379700466681\\
331 0.0114562820330138\\
332 0.0115743240250875\\
333 0.011692022739\\
334 0.0118093050233109\\
335 0.0119260979656852\\
336 0.0120423289997051\\
337 0.0121579260117154\\
338 0.0122728174475743\\
339 0.0123869324191665\\
340 0.012500200810533\\
341 0.0126125533834593\\
342 0.0127239218823624\\
343 0.0128342391383101\\
344 0.0129434391720028\\
345 0.0130514572955471\\
346 0.0131582302128472\\
347 0.013263696118441\\
348 0.0133677947946078\\
349 0.0134704677065749\\
350 0.013571658095655\\
351 0.0136713110701453\\
352 0.013769373693825\\
353 0.0138657950718918\\
354 0.0139605264341824\\
355 0.0140535212155324\\
356 0.0141447351331394\\
357 0.0142341262607757\\
358 0.0143216550996291\\
359 0.0144072846454421\\
360 0.0144909804515168\\
361 0.0145727106870863\\
362 0.0146524461905548\\
363 0.0147301605171651\\
364 0.0148058299807596\\
365 0.0148794336893937\\
366 0.0149509535746144\\
367 0.0150203744142287\\
368 0.0150876838483746\\
369 0.0151528723887111\\
370 0.0152159334206229\\
371 0.0152768631984953\\
372 0.0153356608343282\\
373 0.0153923282801206\\
374 0.0154468703045304\\
375 0.0154992944642969\\
376 0.0155496110708436\\
377 0.0155978331523792\\
378 0.0156439764117447\\
379 0.0156880591802824\\
380 0.0157301023680936\\
381 0.0157701294111507\\
382 0.0158081662157711\\
383 0.0158442411009494\\
384 0.0158783847389753\\
385 0.015910630094624\\
386 0.0159410123629954\\
387 0.0159695689058731\\
388 0.0159963391863429\\
389 0.0160213647013733\\
390 0.0160446889120818\\
391 0.0160663571714484\\
392 0.0160864166492784\\
393 0.0161049162542866\\
394 0.0161219065532711\\
395 0.0161374396874431\\
396 0.0161515692860391\\
397 0.0161643503773753\\
398 0.0161758392975202\\
399 0.0161860935967753\\
400 0.0161951719441688\\
401 0.0162031340301723\\
402 0.0162100404678646\\
403 0.0162159526927671\\
404 0.0162209328615864\\
405 0.0162250437500991\\
406 0.0162283486504169\\
407 0.0162309112678738\\
408 0.0162327956177733\\
409 0.0162340659222341\\
410 0.0162347865073717\\
411 0.0162350217010473\\
412 0.0162348357314135\\
413 0.0162342926264804\\
414 0.0162334561149212\\
415 0.0162323895283276\\
416 0.0162311557051207\\
417 0.0162298168963061\\
418 0.016228434673242\\
419 0.016227069837556\\
420 0.0162257823333063\\
421 0.0162246311614505\\
422 0.0162236742966714\\
423 0.0162229686066002\\
424 0.0162225697734609\\
425 0.016222532218121\\
426 0.0162229090265013\\
427 0.0162237518783139\\
428 0.0162251109781526\\
429 0.016227034989018\\
430 0.016229570968379\\
431 0.0162327643068484\\
432 0.0162366586694966\\
433 0.0162412959397665\\
434 0.0162467161659049\\
435 0.0162529575098043\\
436 0.0162600561981665\\
437 0.0162680464759195\\
438 0.0162769605618517\\
439 0.0162868286064443\\
440 0.0162976786519155\\
441 0.016309536594505\\
442 0.0163224261490499\\
443 0.016336368815901\\
444 0.0163513838502192\\
445 0.0163674882336637\\
446 0.016384696648462\\
447 0.0164030214538436\\
448 0.0164224726648304\\
449 0.0164430579333828\\
450 0.0164647825319178\\
451 0.0164876493392273\\
452 0.0165116588288396\\
453 0.0165368090598666\\
454 0.0165630956703748\\
455 0.0165905118733042\\
456 0.0166190484549556\\
457 0.0166486937760543\\
458 0.0166794337753971\\
459 0.0167112519760841\\
460 0.0167441294943446\\
461 0.0167780450509648\\
462 0.0168129749853318\\
463 0.0168488932721071\\
464 0.0168857715405443\\
465 0.0169235790964676\\
466 0.016962282946922\\
467 0.0170018478275106\\
468 0.0170422362324272\\
469 0.0170834084471926\\
470 0.017125322584099\\
471 0.0171679346203626\\
472 0.0172111984389833\\
473 0.017255065872301\\
474 0.0172994867482389\\
475 0.0173444089392164\\
476 0.0173897784137078\\
477 0.0174355392904232\\
478 0.0174816338950764\\
479 0.0175280028197057\\
480 0.0175745849845112\\
481 0.0176213177021804\\
482 0.0176681367446866\\
483 0.0177149764125653\\
484 0.0177617696066895\\
485 0.0178084479025661\\
486 0.0178549416271463\\
487 0.0179011799381041\\
488 0.0179470909055139\\
489 0.0179926015958578\\
490 0.0180376381583088\\
491 0.0180821259132523\\
492 0.0181259894430226\\
493 0.0181691526848331\\
494 0.018211539025876\\
495 0.0182530714005377\\
496 0.018293672389637\\
497 0.0183332643215601\\
498 0.0183717693751469\\
499 0.0184091096841823\\
500 0.0184452074433389\\
501 0.0184799850153968\\
502 0.0185133650395199\\
503 0.0185452705403323\\
504 0.0185756250375124\\
505 0.0186043526556168\\
506 0.0186313782338395\\
507 0.0186566274354117\\
508 0.0186800268563551\\
509 0.0187015041333171\\
510 0.0187209880502396\\
511 0.0187384086436284\\
512 0.018753697306204\\
513 0.0187667868887154\\
514 0.0187776117997075\\
515 0.0187861081030464\\
516 0.0187922136130115\\
517 0.0187958679867757\\
518 0.0187970128141012\\
519 0.0187955917041109\\
520 0.0187915503690234\\
521 0.0187848367047726\\
522 0.0187754008684388\\
523 0.0187631953524232\\
524 0.0187481750552999\\
525 0.0187302973492826\\
526 0.0187095221442564\\
527 0.0186858119483245\\
528 0.0186591319248123\\
529 0.0186294499456447\\
530 0.0185967366409903\\
531 0.0185609654450553\\
532 0.0185221126379168\\
533 0.0184801573832995\\
534 0.0184350817622142\\
535 0.0183868708023943\\
536 0.0183355125034782\\
537 0.0182809978579012\\
538 0.0182233208674723\\
539 0.0181624785556242\\
540 0.0180984709753375\\
541 0.0180313012127481\\
542 0.0179609753864608\\
543 0.0178875026426\\
544 0.017810895145638\\
545 0.0177311680650499\\
546 0.0176483395578535\\
547 0.0175624307470984\\
548 0.017473465696377\\
549 0.0173814713804352\\
550 0.0172864776519611\\
551 0.0171885172046352\\
552 0.0170876255325312\\
553 0.0169838408859732\\
554 0.016877204223978\\
555 0.016767759163436\\
556 0.0166555519251987\\
557 0.0165406312772541\\
558 0.0164230484751594\\
559 0.0163028571998774\\
560 0.0161801134931394\\
561 0.0160548756904405\\
562 0.0159272043517705\\
563 0.0157971621901932\\
564 0.0156648139983902\\
565 0.0155302265732976\\
566 0.0153934686389764\\
567 0.0152546107678694\\
568 0.0151137253006041\\
569 0.0149708862645019\\
570 0.0148261692909577\\
571 0.0146796515318411\\
572 0.014531411575048\\
573 0.0143815293593021\\
574 0.0142300860882832\\
575 0.0140771641441574\\
576 0.0139228470005927\\
577 0.013767219135349\\
578 0.0136103659425196\\
579 0.0134523736444944\\
580 0.013293329203716\\
581 0.0131333202343318\\
582 0.0129724349138681\\
583 0.0128107618950662\\
584 0.0126483902180159\\
585 0.012485409222706\\
586 0.0123219084620962\\
587 0.0121579776158242\\
588 0.0119937064046847\\
589 0.011829184506047\\
590 0.0116645014703917\\
591 0.0114997466391433\\
592 0.0113350090639615\\
593 0.0111703774276386\\
594 0.0110059399667322\\
595 0.010841784396049\\
596 0.0106779978350801\\
597 0.0105146667364762\\
598 0.0103518768166365\\
599 0.0101897129884764\\
600 0.0100282592964253\\
601 0.00986759885369843\\
602 0.00970781378187519\\
603 0.00954898515280876\\
604 0.00939119293288434\\
605 0.00923451592963469\\
606 0.00907903174071587\\
607 0.00892481670523935\\
608 0.00877194585745122\\
609 0.00862049288274363\\
610 0.0084705300759793\\
611 0.00832212830210481\\
612 0.0081753569590236\\
613 0.00803028394269053\\
614 0.0078869756143694\\
615 0.00774549676995792\\
616 0.00760591061124349\\
617 0.00746827871893056\\
618 0.00733266102728455\\
619 0.007199115800257\\
620 0.00706769960898004\\
621 0.00693846731054285\\
622 0.00681147202798891\\
623 0.00668676513148644\\
624 0.00656439622061779\\
625 0.00644441310771869\\
626 0.00632686180219547\\
627 0.00621178649575803\\
628 0.00609922954852403\\
629 0.005989231475987\\
630 0.00588183093689226\\
631 0.00577706472209605\\
632 0.00567496774447973\\
633 0.00557557302998992\\
634 0.00547891170989469\\
635 0.00538501301435731\\
636 0.00529390426740021\\
637 0.00520561088327091\\
638 0.00512015636416559\\
639 0.00503756229924564\\
640 0.0049578483649072\\
641 0.00488103232631545\\
642 0.00480713004026824\\
643 0.00473615545948178\\
644 0.00466812063838527\\
645 0.00460303574047744\\
646 0.00454090904725183\\
647 0.00448174696864989\\
648 0.00442555405495165\\
649 0.00437233300997145\\
650 0.00432208470539923\\
651 0.00427480819611884\\
652 0.00423050073634306\\
653 0.00418915779642744\\
654 0.00415077308025169\\
655 0.00411533854308157\\
656 0.00408284440984468\\
657 0.00405327919377062\\
658 0.00402662971536264\\
659 0.00400288112168134\\
660 0.00398201690593422\\
661 0.00396401892737521\\
662 0.00394886743152851\\
663 0.00393654107075886\\
664 0.00392701692521769\\
665 0.00392027052420021\\
666 0.00391627586795341\\
667 0.00391500544997876\\
668 0.00391643027987609\\
669 0.00392051990677724\\
670 0.0039272424434192\\
671 0.00393656459090704\\
672 0.00394845166421636\\
673 0.0039628676184802\\
674 0.00397977507609132\\
675 0.00399913535463059\\
676 0.00402090849562259\\
677 0.00404505329413298\\
678 0.00407152732924768\\
679 0.00410028699548245\\
680 0.00413128753514532\\
681 0.00416448307163316\\
682 0.00419982664361779\\
683 0.00423727024007143\\
684 0.00427676483607453\\
685 0.0043182604293372\\
686 0.00436170607736326\\
687 0.00440704993519208\\
688 0.00445423929366009\\
689 0.00450322061813036\\
690 0.00455393958764796\\
691 0.0046063411344827\\
692 0.00466036948402366\\
693 0.00471596819500023\\
694 0.0047730802000123\\
695 0.00483164784634012\\
696 0.00489161293698383\\
697 0.00495291677188136\\
698 0.00501550018927238\\
699 0.00507930360719925\\
700 0.00514426706515672\\
701 0.005210330265918\\
702 0.00527743261757442\\
703 0.00534551327583674\\
704 0.00541451118666358\\
705 0.00548436512928324\\
706 0.0055550137596387\\
707 0.0056263956542275\\
708 0.00569844935425856\\
709 0.00577111341002649\\
710 0.0058443264254063\\
711 0.00591802710238017\\
712 0.00599215428551059\\
713 0.00606664700627113\\
714 0.00614144452714616\\
715 0.00621648638541454\\
716 0.0062917124365385\\
717 0.00636706289708565\\
718 0.00644247838711829\\
719 0.00651789997199032\\
720 0.00659326920349731\\
721 0.00666852816033052\\
722 0.00674361948779046\\
723 0.00681848643671978\\
724 0.00689307290161965\\
725 0.00696732345791705\\
726 0.00704118339835446\\
727 0.00711459876847608\\
728 0.0071875164011879\\
729 0.00725988395037161\\
730 0.00733164992353451\\
731 0.00740276371348028\\
732 0.0074731756289872\\
733 0.00754283692448245\\
734 0.00761169982870271\\
735 0.00767971757233317\\
736 0.00774684441461789\\
737 0.00781303566893642\\
738 0.00787824772734182\\
739 0.00794243808405443\\
740 0.008005565357903\\
741 0.00806758931370418\\
742 0.00812847088257452\\
743 0.00818817218117257\\
744 0.00824665652986693\\
745 0.00830388846982319\\
746 0.00835983377900584\\
747 0.00841445948710849\\
748 0.00846773388944658\\
749 0.00851962655985075\\
750 0.0085701083625765\\
751 0.00861915146321605\\
752 0.00866672933859134\\
753 0.00871281678562728\\
754 0.00875738992923058\\
755 0.00880042622920613\\
756 0.00884190448623002\\
757 0.00888180484687637\\
758 0.00892010880766997\\
759 0.00895679921811852\\
760 0.00899186028267871\\
761 0.00902527756162685\\
762 0.0090570379708204\\
763 0.00908712978034577\\
764 0.0091155426120593\\
765 0.00914226743605028\\
766 0.00916729656608719\\
767 0.00919062365413353\\
768 0.00921224368402509\\
769 0.00923215296439274\\
770 0.00925034912090545\\
771 0.0092668310878909\\
772 0.00928159909935289\\
773 0.00929465467935694\\
774 0.00930600063172966\\
775 0.00931564102902868\\
776 0.00932358120077389\\
777 0.00932982772095372\\
778 0.00933438839481606\\
779 0.00933727224493466\\
780 0.00933848949652691\\
781 0.00933805156199535\\
782 0.00933597102466585\\
783 0.00933226162169816\\
784 0.00932693822614733\\
785 0.00932001682815703\\
786 0.00931151451526899\\
787 0.00930144945183541\\
788 0.00928984085752457\\
789 0.00927670898491251\\
790 0.00926207509615689\\
791 0.00924596143875185\\
792 0.00922839122036605\\
793 0.00920938858276855\\
794 0.00918897857485026\\
795 0.00916718712475174\\
796 0.00914404101111065\\
797 0.00911956783344504\\
798 0.00909379598169135\\
799 0.00906675460491855\\
800 0.00903847357924503\\
801 0.00900898347499591\\
802 0.00897831552315779\\
803 0.00894650158120728\\
804 0.00891357409839639\\
805 0.00887956608057344\\
806 0.00884451105461711\\
807 0.00880844303256741\\
808 0.00877139647554569\\
809 0.00873340625755609\\
810 0.00869450762925197\\
811 0.0086547361817357\\
812 0.00861412781044719\\
813 0.00857271867919286\\
814 0.0085305451843723\\
815 0.00848764391947451\\
816 0.00844405163994181\\
817 0.00839980522853081\\
818 0.00835494166131366\\
819 0.00830949797445302\\
820 0.00826351123186552\\
821 0.00821701849387079\\
822 0.00817005678690078\\
823 0.00812266307431921\\
824 0.00807487422839544\\
825 0.00802672700353783\\
826 0.00797825801103942\\
827 0.0079295036957643\\
828 0.00788050031528893\\
829 0.00783128392194632\\
830 0.00778189034803741\\
831 0.00773235519425025\\
832 0.00768271382113736\\
833 0.00763300134340917\\
834 0.00758325262683054\\
835 0.00753350228761648\\
836 0.00748378469433097\\
837 0.00743413397234262\\
838 0.00738458401087078\\
839 0.00733516847258959\\
840 0.0072859208056692\\
841 0.00723687425803781\\
842 0.00718806189357265\\
843 0.00713951660989346\\
844 0.0070912711574328\\
845 0.00704335815947189\\
846 0.0069958101328481\\
847 0.00694865950905627\\
848 0.00690193865548311\\
849 0.00685567989652962\\
850 0.00680991553439275\\
851 0.00676467786929305\\
852 0.00671999921895043\\
853 0.00667591193712531\\
854 0.00663244843105689\\
855 0.00658964117764497\\
856 0.00654752273823552\\
857 0.00650612577188393\\
858 0.00646548304698323\\
859 0.00642562745115737\\
860 0.00638659199933216\\
861 0.00634840983990857\\
862 0.00631111425897435\\
863 0.00627473868249979\\
864 0.00623931667647484\\
865 0.00620488194496742\\
866 0.00617146832611162\\
867 0.00613910978603275\\
868 0.00610784041065334\\
869 0.00607769439522953\\
870 0.00604870603141655\\
871 0.006020909691692\\
872 0.00599433981105955\\
873 0.00596903086607512\\
874 0.00594501735134775\\
875 0.00592233375372414\\
876 0.00590101452434218\\
877 0.00588109404866239\\
878 0.00586260661452235\\
879 0.00584558637823987\\
880 0.00583006732880205\\
881 0.00581608325019055\\
882 0.0058036676818869\\
883 0.00579285387757184\\
884 0.0057836747620103\\
885 0.00577616288612367\\
886 0.00577035038029868\\
887 0.00576626890605521\\
888 0.00576394960628272\\
889 0.00576342305431649\\
890 0.00576471920212279\\
891 0.00576786732779837\\
892 0.00577289598250279\\
893 0.00577983293686848\\
894 0.005788705126904\\
895 0.00579953859942832\\
896 0.00581235845712676\\
897 0.00582718880336273\\
898 0.00584405268690181\\
899 0.0058629720467101\\
900 0.00588396765698693\\
901 0.00590705907257702\\
902 0.00593226457487797\\
903 0.00595960111831765\\
904 0.00598908427743832\\
905 0.00602072819460363\\
906 0.00605454552834669\\
907 0.00609054740239404\\
908 0.00612874335542121\\
909 0.00616914129161454\\
910 0.00621174743213015\\
911 0.00625656626755556\\
912 0.00630360051149107\\
913 0.00635285105537793\\
914 0.00640431692470841\\
915 0.00645799523675875\\
916 0.00651388115999047\\
917 0.00657196787526864\\
918 0.00663224653904658\\
919 0.00669470624866725\\
920 0.0067593340099298\\
921 0.00682611470706781\\
922 0.00689503107528234\\
923 0.00696606367596873\\
924 0.00703919087477198\\
925 0.00711438882260967\\
926 0.00719163143981978\\
927 0.00727089040362065\\
928 0.00735213513908306\\
929 0.00743533281379026\\
930 0.00752044833630971\\
931 0.00760744435855537\\
932 0.00769628128211343\\
933 0.00778691726864229\\
934 0.0078793082545144\\
935 0.00797340796989505\\
936 0.00806916796241335\\
937 0.00816653762548567\\
938 0.00826546423125188\\
939 0.00836589296801702\\
940 0.00846776698205239\\
941 0.00857102742357826\\
942 0.00867561349670992\\
943 0.00878146251310121\\
944 0.00888850994898921\\
945 0.00899668950534963\\
946 0.00910593317092382\\
947 0.00921617128796486\\
948 0.0093273326206351\\
949 0.0094393444260333\\
950 0.00955213252782382\\
951 0.0096656213924068\\
952 0.00977973420752821\\
953 0.00989439296318395\\
954 0.0100095185346218\\
955 0.0101250307672012\\
956 0.0102408485628521\\
957 0.0103568899678713\\
958 0.0104730722618084\\
959 0.0105893120472109\\
960 0.0107055253400224\\
961 0.0108216276604453\\
962 0.010937534124088\\
963 0.0110531595332462\\
964 0.0111684184681848\\
965 0.011283225378304\\
966 0.0113974946730705\\
967 0.0115111408126041\\
968 0.0116240783978104\\
969 0.0117362222599573\\
970 0.011847487549599\\
971 0.011957789824752\\
972 0.0120670451382352\\
973 0.0121751701240877\\
974 0.0122820820829828\\
975 0.0123876990665588\\
976 0.0124919399605907\\
977 0.0125947245669313\\
978 0.0126959736841496\\
979 0.0127956091868002\\
980 0.0128935541032593\\
981 0.0129897326920625\\
982 0.0130840705166857\\
983 0.0131764945187105\\
984 0.0132669330893162\\
985 0.0133553161390444\\
986 0.0134415751657739\\
987 0.0135256433208322\\
988 0.0136074554731651\\
989 0.0136869482714949\\
990 0.0137640602044117\\
991 0.0138387316583462\\
992 0.0139109049733539\\
993 0.0139805244966251\\
994 0.0140475366336322\\
995 0.0141118898968303\\
996 0.0141735349518551\\
997 0.0142324246611826\\
998 0.014288514125215\\
999 0.0143417607207448\\
1000 0.014392124136738\\
1001 0.0144395664073863\\
};
\addlegendentry{In Ruhe};

\addplot [
color=green!50!black,
solid
]
table[row sep=crcr]{
1 0.350018661833295\\
2 0.350186853798652\\
3 0.350353138328617\\
4 0.350517877742661\\
5 0.35068152144753\\
6 0.350844606807404\\
7 0.351007759800017\\
8 0.35117169544851\\
9 0.351337218021128\\
10 0.351505220993894\\
11 0.35167668677369\\
12 0.351852686180712\\
13 0.352034377691886\\
14 0.352223006452171\\
15 0.352419903065399\\
16 0.352626482174603\\
17 0.352844240833177\\
18 0.353074756661216\\
19 0.353319685783473\\
20 0.353580760554926\\
21 0.353859787088762\\
22 0.354158642603413\\
23 0.354479272600399\\
24 0.354823687877081\\
25 0.355193961371621\\
26 0.355592224833357\\
27 0.356020665311344\\
28 0.356481521455852\\
29 0.356977079630223\\
30 0.357509669832336\\
31 0.358081661426556\\
32 0.358695458689097\\
33 0.359353496172084\\
34 0.360058233893125\\
35 0.360812152357633\\
36 0.361617747420896\\
37 0.362477524997226\\
38 0.36339399562448\\
39 0.364369668893665\\
40 0.365407047754753\\
41 0.366508622711396\\
42 0.367676865918898\\
43 0.368914225201573\\
44 0.370223118007125\\
45 0.371605925316583\\
46 0.373064985528268\\
47 0.374602588333602\\
48 0.37622096860178\\
49 0.377922300289409\\
50 0.37970869039027\\
51 0.381582172939133\\
52 0.383544703082852\\
53 0.385598151232063\\
54 0.387744297307871\\
55 0.389984825098748\\
56 0.392321316743031\\
57 0.394755247351875\\
58 0.397287979787133\\
59 0.399920759608738\\
60 0.402654710206277\\
61 0.405490828129224\\
62 0.408429978629595\\
63 0.411472891429905\\
64 0.414620156728399\\
65 0.417872221452583\\
66 0.421229385771215\\
67 0.424691799873961\\
68 0.428259461027042\\
69 0.431932210912301\\
70 0.435709733256227\\
71 0.43959155175459\\
72 0.443577028297508\\
73 0.447665361498885\\
74 0.451855585533363\\
75 0.456146569283084\\
76 0.460537015795805\\
77 0.465025462055085\\
78 0.469610279062582\\
79 0.474289672231808\\
80 0.47906168209236\\
81 0.483924185303305\\
82 0.488874895974039\\
83 0.493911367290045\\
84 0.499030993440344\\
85 0.50423101184377\\
86 0.509508505672886\\
87 0.514860406675851\\
88 0.520283498296284\\
89 0.525774419088852\\
90 0.531329666425727\\
91 0.536945600486907\\
92 0.542618448525237\\
93 0.54834430939472\\
94 0.554119158329481\\
95 0.55993885196092\\
96 0.565799133561012\\
97 0.57169563849948\\
98 0.577623899902085\\
99 0.583579354497583\\
100 0.5895573486418\\
101 0.595553144507595\\
102 0.60156192642867\\
103 0.607578807384071\\
104 0.613598835609942\\
105 0.619617001325814\\
106 0.62562824356329\\
107 0.631627457084832\\
108 0.637609499379207\\
109 0.643569197718746\\
110 0.649501356262361\\
111 0.655400763187739\\
112 0.661262197835849\\
113 0.667080437850777\\
114 0.672850266297962\\
115 0.678566478744928\\
116 0.684223890291298\\
117 0.689817342539786\\
118 0.695341710505871\\
119 0.70079190946875\\
120 0.70616290176733\\
121 0.711449703541717\\
122 0.716647391414031\\
123 0.721751109095457\\
124 0.726756073902138\\
125 0.731657583161782\\
126 0.736451020495287\\
127 0.741131861962144\\
128 0.745695682064091\\
129 0.750138159607358\\
130 0.754455083427757\\
131 0.758642357983257\\
132 0.762696008816096\\
133 0.766612187882199\\
134 0.77038717874109\\
135 0.774017401595684\\
136 0.777499418169174\\
137 0.780829936405844\\
138 0.784005814983137\\
139 0.787024067623156\\
140 0.789881867192693\\
141 0.792576549581809\\
142 0.795105617351839\\
143 0.797466743144606\\
144 0.799657772845461\\
145 0.801676728493585\\
146 0.803521810933831\\
147 0.805191402205134\\
148 0.806684067661313\\
149 0.807998557820832\\
150 0.809133809942794\\
151 0.810088949327173\\
152 0.810863290337931\\
153 0.811456337148363\\
154 0.81186778420863\\
155 0.812097516436079\\
156 0.812145609129696\\
157 0.812012327611352\\
158 0.811698126598417\\
159 0.811203649314533\\
160 0.810529726346113\\
161 0.809677374250407\\
162 0.808647793917096\\
163 0.807442368681405\\
164 0.806062662184659\\
165 0.804510415978527\\
166 0.802787546871636\\
167 0.800896144020867\\
168 0.798838465772699\\
169 0.796616936260985\\
170 0.794234141766849\\
171 0.791692826845667\\
172 0.788995890225849\\
173 0.786146380484023\\
174 0.78314749150092\\
175 0.780002557702406\\
176 0.776715049090921\\
177 0.773288566074073\\
178 0.769726834098755\\
179 0.766033698100291\\
180 0.762213116776774\\
181 0.758269156699295\\
182 0.754205986269831\\
183 0.750027869540239\\
184 0.745739159907623\\
185 0.74134429370285\\
186 0.736847783689672\\
187 0.732254212491988\\
188 0.727568225966212\\
189 0.722794526534897\\
190 0.717937866496669\\
191 0.713003041326334\\
192 0.707994882978067\\
193 0.702918253204093\\
194 0.697778036901329\\
195 0.692579135498542\\
196 0.68732646039633\\
197 0.682024926471633\\
198 0.676679445658241\\
199 0.671294920615041\\
200 0.665876238494231\\
201 0.660428264822044\\
202 0.654955837504951\\
203 0.649463760974649\\
204 0.643956800485056\\
205 0.638439676573654\\
206 0.632917059698123\\
207 0.627393565057681\\
208 0.621873747607109\\
209 0.616362097270068\\
210 0.610863034357065\\
211 0.605380905192255\\
212 0.599919977952202\\
213 0.59448443871872\\
214 0.589078387747015\\
215 0.583705835949492\\
216 0.578370701594867\\
217 0.573076807221481\\
218 0.567827876763131\\
219 0.562627532885127\\
220 0.557479294527778\\
221 0.552386574654065\\
222 0.547352678197828\\
223 0.542380800208432\\
224 0.537474024187582\\
225 0.532635320613642\\
226 0.527867545648565\\
227 0.52317344002206\\
228 0.518555628086686\\
229 0.514016617036383\\
230 0.509558796280074\\
231 0.505184436962075\\
232 0.50089569162179\\
233 0.496694593986167\\
234 0.492583058889623\\
235 0.488562882317142\\
236 0.484635741566884\\
237 0.480803195528643\\
238 0.47706668507452\\
239 0.473427533558965\\
240 0.469886947426865\\
241 0.466446016929809\\
242 0.463105716951353\\
243 0.459866907941833\\
244 0.456730336962653\\
245 0.453696638839533\\
246 0.450766337424109\\
247 0.447939846963534\\
248 0.445217473577556\\
249 0.442599416841443\\
250 0.440085771471209\\
251 0.437676529105701\\
252 0.435371580179126\\
253 0.433170715878107\\
254 0.431073630179858\\
255 0.429079921971499\\
256 0.427189097253166\\
257 0.425400571426913\\
258 0.423713671669336\\
259 0.422127639380165\\
260 0.420641632695643\\
261 0.419254729056455\\
262 0.417965927825372\\
263 0.416774152957332\\
264 0.415678255731147\\
265 0.414677017554634\\
266 0.413769152852957\\
267 0.412953312045497\\
268 0.412228084612263\\
269 0.411592002247594\\
270 0.411043542095816\\
271 0.410581130060483\\
272 0.410203144177656\\
273 0.40990791804498\\
274 0.409693744299403\\
275 0.409558878133735\\
276 0.409501540836292\\
277 0.409519923332404\\
278 0.409612189703799\\
279 0.409776480662027\\
280 0.410010916953809\\
281 0.410313602678849\\
282 0.410682628503282\\
283 0.411116074754514\\
284 0.411612014385569\\
285 0.412168515799194\\
286 0.412783645524048\\
287 0.41345547073708\\
288 0.414182061627945\\
289 0.414961493602816\\
290 0.415791849326391\\
291 0.416671220602147\\
292 0.417597710092076\\
293 0.418569432878154\\
294 0.419584517868735\\
295 0.420641109053887\\
296 0.421737366614408\\
297 0.422871467889906\\
298 0.424041608211873\\
299 0.42524600160815\\
300 0.42648288138547\\
301 0.427750500596738\\
302 0.429047132399388\\
303 0.430371070311113\\
304 0.431720628369765\\
305 0.433094141205398\\
306 0.43448996403372\\
307 0.43590647258108\\
308 0.437342062951084\\
309 0.438795151442118\\
310 0.440264174323666\\
311 0.441747587577827\\
312 0.443243866610941\\
313 0.444751505938863\\
314 0.446269018848456\\
315 0.447794937037376\\
316 0.449327810233967\\
317 0.450866205798672\\
318 0.452408708307863\\
319 0.453953919120645\\
320 0.455500455928951\\
321 0.457046952291123\\
322 0.458592057149526\\
323 0.460134434333312\\
324 0.461672762048197\\
325 0.46320573235549\\
326 0.464732050642983\\
327 0.466250435090772\\
328 0.467759616135774\\
329 0.469258335939345\\
330 0.470745347863023\\
331 0.47221941595778\\
332 0.473679314472433\\
333 0.475123827386971\\
334 0.47655174797643\\
335 0.477961878410677\\
336 0.479353029394971\\
337 0.480724019855721\\
338 0.48207367667553\\
339 0.483400834481638\\
340 0.484704335492131\\
341 0.485983029424412\\
342 0.487235773470082\\
343 0.488461432338997\\
344 0.489658878373006\\
345 0.490826991727704\\
346 0.4919646606196\\
347 0.493070781636584\\
348 0.49414426011075\\
349 0.495184010553931\\
350 0.496188957156912\\
351 0.497158034352496\\
352 0.498090187440396\\
353 0.498984373269755\\
354 0.499839560973865\\
355 0.500654732751469\\
356 0.501428884689204\\
357 0.502161027620001\\
358 0.502850188012486\\
359 0.50349540888668\\
360 0.504095750751485\\
361 0.504650292559698\\
362 0.505158132676441\\
363 0.50561838985716\\
364 0.506030204231466\\
365 0.506392738289315\\
366 0.506705177866194\\
367 0.506966733124115\\
368 0.507176639525425\\
369 0.507334158796553\\
370 0.507438579878992\\
371 0.507489219864946\\
372 0.507485424915207\\
373 0.507426571156971\\
374 0.507312065559426\\
375 0.507141346785063\\
376 0.506913886014794\\
377 0.506629187745105\\
378 0.506286790555677\\
379 0.505886267846355\\
380 0.505427228543192\\
381 0.504909317774471\\
382 0.504332217518815\\
383 0.503695647228079\\
384 0.502999364427728\\
385 0.502243165296898\\
386 0.501426885229732\\
387 0.500550399379042\\
388 0.499613623183036\\
389 0.498616512875984\\
390 0.49755906598401\\
391 0.496441321807437\\
392 0.495263361890881\\
393 0.494025310481966\\
394 0.492727334979179\\
395 0.491369646369464\\
396 0.489952499656454\\
397 0.488476194280795\\
398 0.486941074534275\\
399 0.485347529968836\\
400 0.48369599579989\\
401 0.481986953301143\\
402 0.480220930185906\\
403 0.478398500968115\\
404 0.476520287295235\\
405 0.474586958245187\\
406 0.472599230580287\\
407 0.470557868953003\\
408 0.46846368606103\\
409 0.466317542752601\\
410 0.464120348086366\\
411 0.461873059352769\\
412 0.45957668206528\\
413 0.457232269929875\\
414 0.45484092479966\\
415 0.452403796618155\\
416 0.449922083349866\\
417 0.447397030891479\\
418 0.444829932953502\\
419 0.442222130901693\\
420 0.439575013549958\\
421 0.436890016900576\\
422 0.43416862383247\\
423 0.4314123637428\\
424 0.42862281214992\\
425 0.425801590266186\\
426 0.422950364547447\\
427 0.420070846223233\\
428 0.417164790808209\\
429 0.414233997592346\\
430 0.411280309105772\\
431 0.408305610554372\\
432 0.405311829223197\\
433 0.402300933846192\\
434 0.399274933942799\\
435 0.396235879124274\\
436 0.39318585837376\\
437 0.390126999303978\\
438 0.3870614673955\\
439 0.383991465217463\\
440 0.38091923163161\\
441 0.37784704097968\\
442 0.374777202253434\\
443 0.371712058245929\\
444 0.368653984682128\\
445 0.365605389326461\\
446 0.362568711064592\\
447 0.359546418956325\\
448 0.356541011256393\\
449 0.353555014399677\\
450 0.350590981947344\\
451 0.347651493490321\\
452 0.344739153506542\\
453 0.341856590168482\\
454 0.339006454097541\\
455 0.336191417062024\\
456 0.333414170615302\\
457 0.330677424670133\\
458 0.327983906003231\\
459 0.32533635668105\\
460 0.322737532393703\\
461 0.320190200680631\\
462 0.317697139030897\\
463 0.315261132843135\\
464 0.312884973233871\\
465 0.310571454686877\\
466 0.308323372539958\\
467 0.306143520308971\\
468 0.304034686851503\\
469 0.301999653374332\\
470 0.300041190290007\\
471 0.298162053928632\\
472 0.29636498311125\\
473 0.294652695591282\\
474 0.293027884370897\\
475 0.291493213900059\\
476 0.290051316166471\\
477 0.288704786684234\\
478 0.28745618038824\\
479 0.286308007441019\\
480 0.285262728959473\\
481 0.284322752670447\\
482 0.283490428505868\\
483 0.28276804414936\\
484 0.282157820546716\\
485 0.281661907392821\\
486 0.281282378607906\\
487 0.28102122781663\\
488 0.28088036384434\\
489 0.280861606245913\\
490 0.280966680883282\\
491 0.281197215567594\\
492 0.281554735781009\\
493 0.282040660491814\\
494 0.282656298075188\\
495 0.283402842350881\\
496 0.284281368748523\\
497 0.285292830611305\\
498 0.286438055648955\\
499 0.287717742550677\\
500 0.289132457768281\\
501 0.290682632479998\\
502 0.292368559746629\\
503 0.294190391872657\\
504 0.296148137985117\\
505 0.298241661842331\\
506 0.300470679882925\\
507 0.302834759522722\\
508 0.305333317703927\\
509 0.307965619698623\\
510 0.310730778167401\\
511 0.313627752473403\\
512 0.316655348251724\\
513 0.319812217233758\\
514 0.323096857325715\\
515 0.326507612940078\\
516 0.330042675578388\\
517 0.333700084663275\\
518 0.337477728617265\\
519 0.341373346185476\\
520 0.345384527998925\\
521 0.349508718374771\\
522 0.353743217349466\\
523 0.358085182940407\\
524 0.362531633631354\\
525 0.367079451076531\\
526 0.371725383018007\\
527 0.376466046410665\\
528 0.381297930748753\\
529 0.386217401587764\\
530 0.391220704255089\\
531 0.396303967742687\\
532 0.401463208774743\\
533 0.406694336043109\\
534 0.411993154603093\\
535 0.417355370422088\\
536 0.422776595073786\\
537 0.428252350571505\\
538 0.433778074335443\\
539 0.439349124289462\\
540 0.44496078408298\\
541 0.450608268432618\\
542 0.456286728577374\\
543 0.461991257840822\\
544 0.467716897294237\\
545 0.473458641515325\\
546 0.479211444437516\\
547 0.484970225283858\\
548 0.490729874577924\\
549 0.49648526022244\\
550 0.502231233635503\\
551 0.507962635934553\\
552 0.513674304158889\\
553 0.519361077521621\\
554 0.525017803680694\\
555 0.530639345016474\\
556 0.536220584901511\\
557 0.541756433948188\\
558 0.547241836223173\\
559 0.55267177542238\\
560 0.558041281002699\\
561 0.563345434264793\\
562 0.568579374376647\\
563 0.573738304324767\\
564 0.578817496781284\\
565 0.583812299880273\\
566 0.588718142902937\\
567 0.593530541875468\\
568 0.598245105084006\\
569 0.602857538509303\\
570 0.607363651181442\\
571 0.611759360452241\\
572 0.616040697179583\\
573 0.620203810815155\\
574 0.624244974386608\\
575 0.628160589366101\\
576 0.631947190416331\\
577 0.635601450000895\\
578 0.639120182840417\\
579 0.642500350192959\\
580 0.645739063936947\\
581 0.648833590435979\\
582 0.651781354166928\\
583 0.654579941095752\\
584 0.65722710178855\\
585 0.659720754247889\\
586 0.662058986466258\\
587 0.664240058690021\\
588 0.666262405388633\\
589 0.668124636925236\\
590 0.669825540925978\\
591 0.671364083346578\\
592 0.672739409235774\\
593 0.673950843196292\\
594 0.674997889544977\\
595 0.675880232174584\\
596 0.676597734120599\\
597 0.677150436837218\\
598 0.677538559187353\\
599 0.677762496152176\\
600 0.677822817266347\\
601 0.677720264785621\\
602 0.677455751594029\\
603 0.677030358858237\\
604 0.676445333436889\\
605 0.675702085052759\\
606 0.674802183235335\\
607 0.673747354041124\\
608 0.672539476558655\\
609 0.671180579205208\\
610 0.669672835822973\\
611 0.668018561583158\\
612 0.666220208707354\\
613 0.664280362016057\\
614 0.662201734314686\\
615 0.659987161627747\\
616 0.657639598291634\\
617 0.655162111916016\\
618 0.652557878223022\\
619 0.649830175773156\\
620 0.646982380587252\\
621 0.644017960674476\\
622 0.640940470477063\\
623 0.637753545242882\\
624 0.634460895337322\\
625 0.631066300506026\\
626 0.627573604100202\\
627 0.623986707276739\\
628 0.620309563186149\\
629 0.616546171161642\\
630 0.612700570922372\\
631 0.608776836802966\\
632 0.60477907202042\\
633 0.600711402988692\\
634 0.596577973691114\\
635 0.592382940121085\\
636 0.588130464801949\\
637 0.58382471139702\\
638 0.579469839419881\\
639 0.575069999053358\\
640 0.570629326083547\\
641 0.566151936954048\\
642 0.561641923945254\\
643 0.557103350483785\\
644 0.552540246587714\\
645 0.547956604453722\\
646 0.54335637419175\\
647 0.538743459710686\\
648 0.53412171475636\\
649 0.529494939101932\\
650 0.524866874890666\\
651 0.520241203131232\\
652 0.515621540345717\\
653 0.511011435370898\\
654 0.506414366314176\\
655 0.501833737666104\\
656 0.497272877571124\\
657 0.492735035257398\\
658 0.488223378626047\\
659 0.483740991999808\\
660 0.479290874030867\\
661 0.474875935767387\\
662 0.470498998878064\\
663 0.466162794033793\\
664 0.461869959445335\\
665 0.457623039555677\\
666 0.453424483885559\\
667 0.449276646030491\\
668 0.445181782807355\\
669 0.441142053548575\\
670 0.437159519541607\\
671 0.433236143611399\\
672 0.429373789843282\\
673 0.425574223443628\\
674 0.42183911073547\\
675 0.418170019286146\\
676 0.414568418163912\\
677 0.411035678320375\\
678 0.407573073095486\\
679 0.404181778841877\\
680 0.400862875665382\\
681 0.397617348278735\\
682 0.394446086965625\\
683 0.391349888652748\\
684 0.388329458087945\\
685 0.385385409123043\\
686 0.382518266099924\\
687 0.379728465337902\\
688 0.37701635671945\\
689 0.374382205369993\\
690 0.371826193425807\\
691 0.369348421882278\\
692 0.366948912513506\\
693 0.364627609854615\\
694 0.362384383240621\\
695 0.360219028899617\\
696 0.358131272101156\\
697 0.356120769360564\\
698 0.354187110696545\\
699 0.352329821934894\\
700 0.35054836704832\\
701 0.348842150522616\\
702 0.347210519742935\\
703 0.345652767399236\\
704 0.344168133914688\\
705 0.342755809902869\\
706 0.341414938658613\\
707 0.340144618684668\\
708 0.338943906253957\\
709 0.337811818006188\\
710 0.336747333577152\\
711 0.335749398258457\\
712 0.334816925684561\\
713 0.333948800542985\\
714 0.333143881302627\\
715 0.33240100295438\\
716 0.331718979758302\\
717 0.331096607993038\\
718 0.330532668705776\\
719 0.330025930463316\\
720 0.329575152104728\\
721 0.329179085492678\\
722 0.328836478255155\\
723 0.328546076505107\\
724 0.328306627523669\\
725 0.328116882392577\\
726 0.327975598562203\\
727 0.327881542343436\\
728 0.327833491313837\\
729 0.327830236629879\\
730 0.32787058523739\\
731 0.327953361972133\\
732 0.328077411542756\\
733 0.328241600389101\\
734 0.328444818409855\\
735 0.328685980554549\\
736 0.328964028275865\\
737 0.329277930839107\\
738 0.329626686486541\\
739 0.330009323455132\\
740 0.330424900846914\\
741 0.330872509351954\\
742 0.33135127182453\\
743 0.331860343713717\\
744 0.332398913350199\\
745 0.332966202091569\\
746 0.333561464328931\\
747 0.334183987358011\\
748 0.334833091118405\\
749 0.335508127804976\\
750 0.3362084813557\\
751 0.336933566820577\\
752 0.337682829616314\\
753 0.338455744671309\\
754 0.33925181546501\\
755 0.340070572965238\\
756 0.340911574467147\\
757 0.341774402338201\\
758 0.342658662674313\\
759 0.34356398387257\\
760 0.344490015125713\\
761 0.345436424843498\\
762 0.346402899006608\\
763 0.347389139459494\\
764 0.348394862148756\\
765 0.349419795313303\\
766 0.350463677631954\\
767 0.351526256333866\\
768 0.352607285277706\\
769 0.353706523006175\\
770 0.354823730782851\\
771 0.355958670618045\\
772 0.357111103290132\\
773 0.35828078636921\\
774 0.35946747225081\\
775 0.360670906208221\\
776 0.361890824472359\\
777 0.363126952348046\\
778 0.364379002375243\\
779 0.365646672543314\\
780 0.366929644565896\\
781 0.368227582223461\\
782 0.369540129780194\\
783 0.370866910481267\\
784 0.372207525136018\\
785 0.37356155079168\\
786 0.374928539500959\\
787 0.376308017185014\\
788 0.377699482591895\\
789 0.379102406350125\\
790 0.380516230118011\\
791 0.381940365830637\\
792 0.383374195047463\\
793 0.384817068404294\\
794 0.386268305174076\\
795 0.387727192940766\\
796 0.389192987388579\\
797 0.390664912206304\\
798 0.392142159104694\\
799 0.393623887944987\\
800 0.395109226977779\\
801 0.39659727319266\\
802 0.398087092779377\\
803 0.399577721700981\\
804 0.401068166378886\\
805 0.402557404489204\\
806 0.404044385869187\\
807 0.405528033532153\\
808 0.407007244788767\\
809 0.40848089247219\\
810 0.409947826264154\\
811 0.411406874118711\\
812 0.412856843780044\\
813 0.414296524390448\\
814 0.415724688184308\\
815 0.417140092263643\\
816 0.418541480450603\\
817 0.41992758521206\\
818 0.421297129651308\\
819 0.422648829561711\\
820 0.423981395537019\\
821 0.425293535132977\\
822 0.426583955074757\\
823 0.427851363504659\\
824 0.429094472264405\\
825 0.430311999205954\\
826 0.431502670524184\\
827 0.432665223103904\\
828 0.433798406872707\\
829 0.434900987150426\\
830 0.435971746985678\\
831 0.437009489470216\\
832 0.43801304002259\\
833 0.438981248633702\\
834 0.439912992068106\\
835 0.440807176016033\\
836 0.441662737192085\\
837 0.442478645377389\\
838 0.443253905402783\\
839 0.44398755907133\\
840 0.444678687019051\\
841 0.445326410513227\\
842 0.445929893187593\\
843 0.446488342712967\\
844 0.447001012400295\\
845 0.447467202731209\\
846 0.447886262809586\\
847 0.448257591726378\\
848 0.448580639829266\\
849 0.448854909888626\\
850 0.449079958152238\\
851 0.449255395283227\\
852 0.449380887179107\\
853 0.449456155674156\\
854 0.44948097913133\\
855 0.449455192931962\\
856 0.449378689870834\\
857 0.449251420461503\\
858 0.449073393153257\\
859 0.448844674457796\\
860 0.448565388981658\\
861 0.4482357193601\\
862 0.447855906089626\\
863 0.447426247259294\\
864 0.446947098184418\\
865 0.446418870949671\\
866 0.445842033871237\\
867 0.44521711088877\\
868 0.444544680897218\\
869 0.443825377026239\\
870 0.443059885872136\\
871 0.44224894668493\\
872 0.441393350511802\\
873 0.440493939297912\\
874 0.43955160494652\\
875 0.438567288342297\\
876 0.437541978344\\
877 0.436476710754014\\
878 0.43537256727122\\
879 0.434230674429475\\
880 0.433052202517965\\
881 0.43183836447445\\
882 0.430590414739974\\
883 0.429309648063802\\
884 0.427997398248791\\
885 0.426655036829148\\
886 0.4252839716742\\
887 0.423885645513335\\
888 0.422461534378773\\
889 0.421013145964096\\
890 0.419542017897725\\
891 0.418049715931583\\
892 0.416537832046236\\
893 0.415007982474657\\
894 0.413461805647603\\
895 0.411900960064291\\
896 0.410327122092725\\
897 0.408741983704549\\
898 0.407147250149838\\
899 0.405544637577602\\
900 0.403935870608183\\
901 0.402322679863975\\
902 0.400706799465169\\
903 0.399089964497393\\
904 0.397473908458324\\
905 0.395860360690632\\
906 0.394251043809127\\
907 0.392647671130406\\
908 0.391051944113568\\
909 0.389465549820302\\
910 0.387890158402202\\
911 0.386327420622647\\
912 0.384778965419877\\
913 0.383246397517093\\
914 0.381731295084382\\
915 0.380235207456263\\
916 0.37875965290777\\
917 0.377306116491457\\
918 0.375876047937432\\
919 0.374470859618222\\
920 0.373091924579623\\
921 0.371740574637962\\
922 0.370418098543673\\
923 0.369125740210814\\
924 0.367864697011942\\
925 0.36663611813753\\
926 0.365441103019225\\
927 0.36428069981665\\
928 0.363155903967857\\
929 0.36206765680382\\
930 0.361016844227358\\
931 0.360004295457111\\
932 0.359030781837501\\
933 0.358097015715779\\
934 0.357203649387492\\
935 0.356351274112104\\
936 0.355540419201425\\
937 0.354771551184554\\
938 0.35404507305367\\
939 0.353361323595233\\
940 0.352720576810984\\
941 0.352123041432682\\
942 0.351568860533809\\
943 0.35105811124068\\
944 0.350590804544799\\
945 0.350166885218187\\
946 0.349786231833869\\
947 0.349448656894511\\
948 0.349153907072696\\
949 0.348901663565832\\
950 0.348691542567192\\
951 0.348523095852942\\
952 0.348395811484424\\
953 0.348309114625707\\
954 0.348262368477657\\
955 0.348254875330922\\
956 0.348285877741227\\
957 0.348354559830605\\
958 0.348460048716272\\
959 0.348601416065068\\
960 0.348777679767951\\
961 0.34898780572778\\
962 0.349230709754446\\
963 0.349505259562278\\
964 0.34981027686413\\
965 0.350144539555234\\
966 0.350506783978912\\
967 0.350895707265987\\
968 0.351309969739899\\
969 0.351748197379753\\
970 0.352208984333767\\
971 0.352690895475814\\
972 0.353192468998007\\
973 0.353712219032483\\
974 0.354248638295785\\
975 0.35480020074947\\
976 0.355365364270795\\
977 0.355942573327543\\
978 0.356530261651306\\
979 0.357126854903725\\
980 0.357730773330434\\
981 0.358340434397643\\
982 0.358954255406551\\
983 0.359570656080936\\
984 0.360188061123517\\
985 0.360804902736864\\
986 0.361419623104892\\
987 0.36203067683131\\
988 0.362636533331981\\
989 0.363235679178827\\
990 0.363826620393668\\
991 0.364407884691081\\
992 0.364978023670128\\
993 0.365535614955338\\
994 0.366079264287552\\
995 0.366607607565125\\
996 0.367119312836091\\
997 0.367613082242281\\
998 0.36808765391667\\
999 0.368541803834865\\
1000 0.368974347620351\\
1001 0.369384142301124\\
};
\addlegendentry{Bewegung};

\addplot [
color=red,
solid
]
table[row sep=crcr]{
1 0.099936580515523\\
2 0.10075248932234\\
3 0.101561430426685\\
4 0.102362874757831\\
5 0.103156303131262\\
6 0.103941206599376\\
7 0.104717086797335\\
8 0.105483456282679\\
9 0.106239838867466\\
10 0.106985769941774\\
11 0.107720796787524\\
12 0.108444478881656\\
13 0.109156388187802\\
14 0.10985610943567\\
15 0.11054324038744\\
16 0.111217392090542\\
17 0.111878189116281\\
18 0.112525269783811\\
19 0.113158286369053\\
20 0.113776905298205\\
21 0.114380807325536\\
22 0.114969687695243\\
23 0.115543256287165\\
24 0.116101237746238\\
25 0.116643371595575\\
26 0.117169412333141\\
27 0.117679129512025\\
28 0.118172307804393\\
29 0.118648747049432\\
30 0.119108262285911\\
31 0.119550683770423\\
32 0.119975856982666\\
33 0.120383642619245\\
34 0.120773916577594\\
35 0.121146569931886\\
36 0.121501508903071\\
37 0.12183865482508\\
38 0.12215794410868\\
39 0.122459328203406\\
40 0.122742773556932\\
41 0.123008261570277\\
42 0.123255788546733\\
43 0.12348536563238\\
44 0.123697018746268\\
45 0.123890788498601\\
46 0.124066730095258\\
47 0.124224913226944\\
48 0.12436542194135\\
49 0.124488354497119\\
50 0.124593823198774\\
51 0.124681954212037\\
52 0.124752887358904\\
53 0.124806775891754\\
54 0.124843786245783\\
55 0.124864097769312\\
56 0.124867902432023\\
57 0.124855404512003\\
58 0.12482682026339\\
59 0.124782377567264\\
60 0.124722315568821\\
61 0.124646884304064\\
62 0.12455634431914\\
63 0.124450966285402\\
64 0.124331030613014\\
65 0.12419682706575\\
66 0.124048654379318\\
67 0.12388681988533\\
68 0.12371163914285\\
69 0.123523435579216\\
70 0.12332254014147\\
71 0.123109290959071\\
72 0.122884033017793\\
73 0.122647117843871\\
74 0.122398903196835\\
75 0.12213975276933\\
76 0.121870035892453\\
77 0.121590127245626\\
78 0.121300406570387\\
79 0.121001258387429\\
80 0.120693071715884\\
81 0.120376239793258\\
82 0.120051159794024\\
83 0.119718232544608\\
84 0.119377862232614\\
85 0.119030456108342\\
86 0.118676424177011\\
87 0.118316178880351\\
88 0.117950134766577\\
89 0.117578708147966\\
90 0.117202316745546\\
91 0.116821379320601\\
92 0.116436315292928\\
93 0.116047544345948\\
94 0.115655486018957\\
95 0.115260559286954\\
96 0.114863182128632\\
97 0.114463771083228\\
98 0.114062740797046\\
99 0.113660503560575\\
100 0.113257468837192\\
101 0.112854042784531\\
102 0.112450627769651\\
103 0.112047621879189\\
104 0.111645418425734\\
105 0.11124440545165\\
106 0.110844965231584\\
107 0.110447473774863\\
108 0.110052300329004\\
109 0.109659806885627\\
110 0.109270347690248\\
111 0.108884268757618\\
112 0.108501907394509\\
113 0.108123591732127\\
114 0.107749640270407\\
115 0.107380361436482\\
116 0.107016053159351\\
117 0.106657002462556\\
118 0.106303485076794\\
119 0.105955765074821\\
120 0.105614094531446\\
121 0.105278713211255\\
122 0.104949848285977\\
123 0.104627714082512\\
124 0.104312511862036\\
125 0.104004429630381\\
126 0.103703641979882\\
127 0.103410309962862\\
128 0.103124580996916\\
129 0.102846588802\\
130 0.102576453369184\\
131 0.102314280960848\\
132 0.102060164142567\\
133 0.101814181847871\\
134 0.101576399478031\\
135 0.101346869039311\\
136 0.101125629319487\\
137 0.100912706104246\\
138 0.10070811243272\\
139 0.100511848890163\\
140 0.100323903934883\\
141 0.10014425425615\\
142 0.0999728651599175\\
143 0.0998096909796707\\
144 0.0996546755103103\\
145 0.0995077524634304\\
146 0.0993688459424286\\
147 0.0992378709356781\\
148 0.0991147338257326\\
149 0.0989993329123533\\
150 0.0988915589469185\\
151 0.0987912956753825\\
152 0.0986984203865222\\
153 0.0986128044619724\\
154 0.0985343139245886\\
155 0.0984628099819039\\
156 0.0983981495617382\\
157 0.0983401858373191\\
158 0.0982887687395574\\
159 0.0982437454543952\\
160 0.0982049609033923\\
161 0.0981722582059633\\
162 0.0981454791218965\\
163 0.0981244644730007\\
164 0.0981090545429216\\
165 0.0980990894543523\\
166 0.0980944095230394\\
167 0.0980948555881411\\
168 0.0981002693186476\\
169 0.0981104934957084\\
170 0.0981253722708414\\
171 0.0981447514001154\\
172 0.0981684784545071\\
173 0.0981964030067317\\
174 0.0982283767949385\\
175 0.0982642538637413\\
176 0.0983038906831345\\
177 0.0983471462459253\\
178 0.0983938821444436\\
179 0.0984439626275171\\
180 0.0984972546390288\\
181 0.0985536278397203\\
182 0.0986129546141505\\
183 0.0986751100647933\\
184 0.0987399719952279\\
185 0.0988074208844278\\
186 0.0988773398543741\\
187 0.0989496146334652\\
188 0.0990241335183699\\
189 0.0991007873371976\\
190 0.0991794694171838\\
191 0.0992600755601756\\
192 0.0993425040285521\\
193 0.0994266555427729\\
194 0.0995124332901066\\
195 0.099599742942999\\
196 0.0996884926851379\\
197 0.0997785932432007\\
198 0.0998699579223701\\
199 0.0999625026440439\\
200 0.100056145984559\\
201 0.100150809213839\\
202 0.100246416332696\\
203 0.10034289410748\\
204 0.100440172101075\\
205 0.100538182699441\\
206 0.100636861132762\\
207 0.100736145489766\\
208 0.100835976723575\\
209 0.10093629864758\\
210 0.101037057920299\\
211 0.10113820401867\\
212 0.101239689199554\\
213 0.101341468449629\\
214 0.101443499424152\\
215 0.101545742375418\\
216 0.101648160071821\\
217 0.101750717708223\\
218 0.101853382807981\\
219 0.101956125116686\\
220 0.102058916487719\\
221 0.102161730760073\\
222 0.102264543629374\\
223 0.102367332513311\\
224 0.102470076412895\\
225 0.10257275577091\\
226 0.102675352328826\\
227 0.102777848983348\\
228 0.102880229643621\\
229 0.102982479090056\\
230 0.103084582835609\\
231 0.103186526990293\\
232 0.103288298129583\\
233 0.103389883167339\\
234 0.103491269233759\\
235 0.103592443558848\\
236 0.103693393361814\\
237 0.103794105746723\\
238 0.103894567604735\\
239 0.10399476552314\\
240 0.104094685701412\\
241 0.104194313874417\\
242 0.104293635242885\\
243 0.104392634411222\\
244 0.104491295332686\\
245 0.104589601261917\\
246 0.104687534714807\\
247 0.104785077435624\\
248 0.104882210371298\\
249 0.104978913652772\\
250 0.105075166583294\\
251 0.105170947633654\\
252 0.105266234444513\\
253 0.105361003836126\\
254 0.105455231825835\\
255 0.105548893653724\\
256 0.105641963816914\\
257 0.105734416113174\\
258 0.105826223694839\\
259 0.105917359134403\\
260 0.106007794503411\\
261 0.106097501466443\\
262 0.106186451391838\\
263 0.10627461548048\\
264 0.106361964913376\\
265 0.106448471017984\\
266 0.106534105452309\\
267 0.106618840405161\\
268 0.106702648810648\\
269 0.106785504575194\\
270 0.106867382815871\\
271 0.106948260109718\\
272 0.107028114754802\\
273 0.107106927044893\\
274 0.107184679560394\\
275 0.107261357478085\\
276 0.107336948900889\\
277 0.107411445206943\\
278 0.107484841415276\\
279 0.107557136563886\\
280 0.107628334095164\\
281 0.107698442243716\\
282 0.107767474422495\\
283 0.107835449604066\\
284 0.107902392694233\\
285 0.107968334894878\\
286 0.108033314052379\\
287 0.108097374987551\\
288 0.108160569803189\\
289 0.108222958165677\\
290 0.108284607557986\\
291 0.108345593502324\\
292 0.108405999751662\\
293 0.108465918449994\\
294 0.108525450261565\\
295 0.108584704469397\\
296 0.108643799043361\\
297 0.108702860677908\\
298 0.108762024799185\\
299 0.108821435540832\\
300 0.108881245687387\\
301 0.108941616584122\\
302 0.109002718012343\\
303 0.109064728029432\\
304 0.109127832773247\\
305 0.109192226230784\\
306 0.109258109971238\\
307 0.109325692843886\\
308 0.109395190641415\\
309 0.109466825729514\\
310 0.109540826643764\\
311 0.109617427654988\\
312 0.109696868304423\\
313 0.109779392910145\\
314 0.109865250046369\\
315 0.109954691997295\\
316 0.110047974187292\\
317 0.11014535458927\\
318 0.110247093113176\\
319 0.110353450976579\\
320 0.110464690059342\\
321 0.110581072244271\\
322 0.110702858745475\\
323 0.110830309425841\\
324 0.110963682104876\\
325 0.111103231858188\\
326 0.11124921031018\\
327 0.111401864922049\\
328 0.11156143827752\\
329 0.11172816736897\\
330 0.111902282886576\\
331 0.112084008512952\\
332 0.112273560225462\\
333 0.112471145608045\\
334 0.112676963174118\\
335 0.112891201702081\\
336 0.113114039585083\\
337 0.113345644196749\\
338 0.11358617127447\\
339 0.113835764321469\\
340 0.114094554028435\\
341 0.114362657715348\\
342 0.11464017879455\\
343 0.114927206257199\\
344 0.115223814186566\\
345 0.115530061302383\\
346 0.115845990540598\\
347 0.116171628672557\\
348 0.116506985967417\\
349 0.116852055901373\\
350 0.117206814917236\\
351 0.117571222237753\\
352 0.117945219735782\\
353 0.118328731863871\\
354 0.118721665645064\\
355 0.119123910726018\\
356 0.11953533949276\\
357 0.1199558072487\\
358 0.120385152453685\\
359 0.120823197021959\\
360 0.121269746676074\\
361 0.121724591353318\\
362 0.122187505661224\\
363 0.12265824937914\\
364 0.123136568003806\\
365 0.12362219333816\\
366 0.124114844123776\\
367 0.124614226717828\\
368 0.125120035815411\\
369 0.125631955217621\\
370 0.126149658645322\\
371 0.126672810598089\\
372 0.127201067257405\\
373 0.127734077432805\\
374 0.12827148354934\\
375 0.128812922674408\\
376 0.129358027581739\\
377 0.129906427850068\\
378 0.130457750993825\\
379 0.13101162362296\\
380 0.131567672628915\\
381 0.132125526393687\\
382 0.13268481601924\\
383 0.133245176574758\\
384 0.133806248358934\\
385 0.134367678172772\\
386 0.134929120595713\\
387 0.135490239255804\\
388 0.136050708084352\\
389 0.136610212546757\\
390 0.137168450843001\\
391 0.13772513507251\\
392 0.138279992358849\\
393 0.138832765930143\\
394 0.139383216151463\\
395 0.139931121505727\\
396 0.140476279520001\\
397 0.141018507634381\\
398 0.141557644010924\\
399 0.142093548280383\\
400 0.142626102224633\\
401 0.143155210392698\\
402 0.143680800648217\\
403 0.144202824645994\\
404 0.144721258235037\\
405 0.145236101785228\\
406 0.145747380434509\\
407 0.146255144253478\\
408 0.146759468324578\\
409 0.147260452733602\\
410 0.147758222471962\\
411 0.148252927249153\\
412 0.148744741215925\\
413 0.149233862599714\\
414 0.149720513254338\\
415 0.150204938125916\\
416 0.150687404636455\\
417 0.151168201986284\\
418 0.151647640376674\\
419 0.152126050154543\\
420 0.152603780881937\\
421 0.153081200333966\\
422 0.15355869343029\\
423 0.154036661106671\\
424 0.154515519134118\\
425 0.154995696893611\\
426 0.155477636114288\\
427 0.155961789582158\\
428 0.156448619824712\\
429 0.15693859777489\\
430 0.157432201416656\\
431 0.157929914414646\\
432 0.158432224731243\\
433 0.158939623235165\\
434 0.15945260230579\\
435 0.15997165443724\\
436 0.160497270846137\\
437 0.161029940087264\\
438 0.161570146681902\\
439 0.162118369763976\\
440 0.162675081749015\\
441 0.163240747030474\\
442 0.163815820707167\\
443 0.164400747344545\\
444 0.164995959771613\\
445 0.165601877914912\\
446 0.166218907671331\\
447 0.166847439822094\\
448 0.16748784899055\\
449 0.168140492645996\\
450 0.168805710154597\\
451 0.169483821877039\\
452 0.170175128311359\\
453 0.170879909278772\\
454 0.171598423150315\\
455 0.172330906112329\\
456 0.173077571469138\\
457 0.173838608981538\\
458 0.174614184240022\\
459 0.175404438071871\\
460 0.176209485981462\\
461 0.177029417623361\\
462 0.177864296307915\\
463 0.178714158539219\\
464 0.179579013585461\\
465 0.180458843081776\\
466 0.181353600665815\\
467 0.182263211646339\\
468 0.183187572705188\\
469 0.184126551633068\\
470 0.185079987099593\\
471 0.186047688458075\\
472 0.187029435585569\\
473 0.188024978758674\\
474 0.189034038565607\\
475 0.190056305855088\\
476 0.191091441722729\\
477 0.192139077535927\\
478 0.193198814998535\\
479 0.194270226256622\\
480 0.195352854046168\\
481 0.1964462118829\\
482 0.197549784293996\\
483 0.198663027091257\\
484 0.199785367685362\\
485 0.200916205440993\\
486 0.20205491207276\\
487 0.203200832081989\\
488 0.204353283234454\\
489 0.205511557078994\\
490 0.206674919506682\\
491 0.207842611349971\\
492 0.209013849020902\\
493 0.210187825187287\\
494 0.211363709485726\\
495 0.212540649270289\\
496 0.213717770395729\\
497 0.214894178034023\\
498 0.216068957523032\\
499 0.217241175246047\\
500 0.218409879541087\\
501 0.219574101638942\\
502 0.220732856629112\\
503 0.221885144452989\\
504 0.223029950923881\\
505 0.224166248773623\\
506 0.22529299872568\\
507 0.226409150594623\\
508 0.227513644411856\\
509 0.228605411577334\\
510 0.229683376036846\\
511 0.230746455484147\\
512 0.231793562586787\\
513 0.232823606233917\\
514 0.233835492803705\\
515 0.234828127447547\\
516 0.23580041538805\\
517 0.236751263227975\\
518 0.237679580267515\\
519 0.238584279827558\\
520 0.239464280576638\\
521 0.240318507859283\\
522 0.241145895023382\\
523 0.241945384744113\\
524 0.242715930341911\\
525 0.243456497091968\\
526 0.244166063522733\\
527 0.244843622700921\\
528 0.245488183500595\\
529 0.24609877185397\\
530 0.246674431981757\\
531 0.247214227600941\\
532 0.24771724310802\\
533 0.248182584735856\\
534 0.248609381682376\\
535 0.248996787209487\\
536 0.249343979710654\\
537 0.249650163745711\\
538 0.249914571041556\\
539 0.25013646145748\\
540 0.250315123913965\\
541 0.250449877283902\\
542 0.250540071245213\\
543 0.250585087094008\\
544 0.250584338517436\\
545 0.250537272325517\\
546 0.250443369141264\\
547 0.250302144048536\\
548 0.250113147197081\\
549 0.249875964364302\\
550 0.24959021747331\\
551 0.249255565066985\\
552 0.248871702737898\\
553 0.248438363514226\\
554 0.247955318201943\\
555 0.24742237568376\\
556 0.24683938317534\\
557 0.246206226439365\\
558 0.245522829958003\\
559 0.244789157064203\\
560 0.244005210032073\\
561 0.243171030126451\\
562 0.242286697611582\\
563 0.241352331718745\\
564 0.240368090572693\\
565 0.239334171077\\
566 0.238250808758798\\
567 0.237118277573757\\
568 0.235936889672344\\
569 0.234706995128353\\
570 0.233428981630448\\
571 0.232103274137351\\
572 0.230730334497429\\
573 0.229310661033895\\
574 0.227844788097365\\
575 0.226333285587839\\
576 0.2247767584482\\
577 0.223175846131147\\
578 0.221531222041248\\
579 0.219843592953441\\
580 0.218113698408819\\
581 0.216342310087958\\
582 0.214530231161489\\
583 0.212678295617362\\
584 0.210787367564138\\
585 0.208858340509719\\
586 0.206892136615259\\
587 0.204889705924514\\
588 0.202852025569548\\
589 0.200780098953991\\
590 0.198674954915099\\
591 0.196537646865766\\
592 0.194369251918144\\
593 0.192170869991518\\
594 0.189943622908453\\
595 0.18768865348418\\
596 0.185407124614236\\
597 0.183100218364665\\
598 0.180769135067985\\
599 0.178415092427172\\
600 0.176039324629315\\
601 0.173643081470228\\
602 0.171227627491088\\
603 0.168794241127883\\
604 0.166344213874279\\
605 0.163878849458283\\
606 0.161399463032931\\
607 0.158907380381025\\
608 0.15640393713384\\
609 0.153890478003576\\
610 0.151368356029213\\
611 0.148838931835346\\
612 0.146303572903469\\
613 0.143763652855135\\
614 0.141220550746331\\
615 0.138675650372383\\
616 0.13613033958265\\
617 0.133586009604244\\
618 0.131044054373938\\
619 0.128505869877158\\
620 0.125972853492311\\
621 0.123446403337675\\
622 0.120927917617084\\
623 0.118418793960272\\
624 0.115920428754044\\
625 0.113434216461137\\
626 0.110961548924413\\
627 0.108503814654678\\
628 0.106062398101051\\
629 0.103638678903325\\
630 0.101234031126172\\
631 0.0988498224754018\\
632 0.0964874134966269\\
633 0.0941481567566913\\
634 0.0918333960080795\\
635 0.0895444653363473\\
636 0.0872826882905539\\
637 0.0850493769967891\\
638 0.082845831255075\\
639 0.0806733376200461\\
640 0.0785331684658894\\
641 0.0764265810361979\\
642 0.0743548164797664\\
643 0.0723190988738479\\
644 0.0703206342369082\\
645 0.0683606095334501\\
646 0.066440191674277\\
647 0.0645605265173988\\
648 0.0627227378783381\\
649 0.0609279265627249\\
650 0.0591771694358135\\
651 0.057471518541301\\
652 0.0558120002773851\\
653 0.0541996146342538\\
654 0.0526353344950974\\
655 0.051120105001278\\
656 0.0496548429805249\\
657 0.0482404364349838\\
658 0.0468777440843404\\
659 0.0455675949584939\\
660 0.0443107880342439\\
661 0.0431080919108933\\
662 0.0419602445203479\\
663 0.0408679528678614\\
664 0.0398318927995998\\
665 0.0388527087927532\\
666 0.0379310137636307\\
667 0.0370673888894796\\
668 0.0362623834404713\\
669 0.0355165146188972\\
670 0.03483026740283\\
671 0.0342040943914503\\
672 0.0336384156493081\\
673 0.0331336185470905\\
674 0.0326900575967057\\
675 0.032308054278521\\
676 0.0319878968586137\\
677 0.0317298401941485\\
678 0.0315341055254408\\
679 0.0314008802537691\\
680 0.0313303177044782\\
681 0.0313225368753513\\
682 0.0313776221706242\\
683 0.0314956231213764\\
684 0.0316765540933498\\
685 0.0319203939835351\\
686 0.0322270859071166\\
687 0.0325965368765901\\
688 0.0330286174750578\\
689 0.033523161525871\\
690 0.0340799657609259\\
691 0.034698789490033\\
692 0.0353793542738682\\
693 0.0361213436030826\\
694 0.036924402586194\\
695 0.0377881376489124\\
696 0.0387121162475736\\
697 0.0396958665994098\\
698 0.0407388774325507\\
699 0.0418405977589116\\
700 0.0430004366733412\\
701 0.0442177631823589\\
702 0.0454919060654845\\
703 0.0468221537716746\\
704 0.0482077543529248\\
705 0.0496479154368364\\
706 0.0511418042397939\\
707 0.0526885476221973\\
708 0.0542872321870001\\
709 0.0559369044228261\\
710 0.0576365708931763\\
711 0.0593851984736619\\
712 0.0611817146399235\\
713 0.0630250078097259\\
714 0.0649139277429085\\
715 0.0668472860020207\\
716 0.0688238564752247\\
717 0.0708423759622378\\
718 0.0729015448236906\\
719 0.075000027693739\\
720 0.0771364542549872\\
721 0.0793094200741919\\
722 0.0815174874971733\\
723 0.0837591866017315\\
724 0.0860330162077408\\
725 0.0883374449435418\\
726 0.0906709123671476\\
727 0.0930318301400466\\
728 0.0954185832510106\\
729 0.0978295312873617\\
730 0.100263009751216\\
731 0.102717331418569\\
732 0.105190787741157\\
733 0.107681650295738\\
734 0.110188172290398\\
735 0.112708590138432\\
736 0.115241125106477\\
737 0.117783985037867\\
738 0.120335366147499\\
739 0.122893454882186\\
740 0.125456429840055\\
741 0.128022463742941\\
742 0.130589725456322\\
743 0.133156382052547\\
744 0.135720600914882\\
745 0.138280551881289\\
746 0.140834409427239\\
747 0.143380354886439\\
748 0.145916578707585\\
749 0.148441282744631\\
750 0.150952682577388\\
751 0.15344900985878\\
752 0.155928514684577\\
753 0.158389467981011\\
754 0.160830163905354\\
755 0.163248922254221\\
756 0.165644090874114\\
757 0.16801404806855\\
758 0.17035720499594\\
759 0.172672008052316\\
760 0.174956941232892\\
761 0.177210528466468\\
762 0.179431335916653\\
763 0.181617974244029\\
764 0.18376910082407\\
765 0.185883421917815\\
766 0.187959694796315\\
767 0.189996729823835\\
768 0.191993392505406\\
769 0.193948605499968\\
770 0.195861350593248\\
771 0.197730670618442\\
772 0.199555671309657\\
773 0.201335523072371\\
774 0.203069462656059\\
775 0.204756794716389\\
776 0.206396893257216\\
777 0.207989202944297\\
778 0.209533240282151\\
779 0.211028594644504\\
780 0.212474929148716\\
781 0.213871981366095\\
782 0.215219563862161\\
783 0.216517564563064\\
784 0.217765946946114\\
785 0.218964750053567\\
786 0.220114088329494\\
787 0.221214151279916\\
788 0.222265202956663\\
789 0.223267581266008\\
790 0.224221697104061\\
791 0.225128033322242\\
792 0.225987143527155\\
793 0.226799650719718\\
794 0.227566245778316\\
795 0.228287685790697\\
796 0.228964792239465\\
797 0.229598449046478\\
798 0.230189600481857\\
799 0.23073924894366\\
800 0.23124845261451\\
801 0.231718323001811\\
802 0.232150022368832\\
803 0.232544761064894\\
804 0.232903794764023\\
805 0.233228421622185\\
806 0.233519979363307\\
807 0.233779842303856\\
808 0.234009418325108\\
809 0.234210145801804\\
810 0.234383490495717\\
811 0.234530942422786\\
812 0.234654012702681\\
813 0.234754230399774\\
814 0.234833139364427\\
815 0.234892295083208\\
816 0.234933261546276\\
817 0.234957608139702\\
818 0.234966906570098\\
819 0.234962727828461\\
820 0.234946639199745\\
821 0.234920201324259\\
822 0.234884965316583\\
823 0.234842469947338\\
824 0.23479423889273\\
825 0.234741778056455\\
826 0.234686572968187\\
827 0.234630086262517\\
828 0.234573755241902\\
829 0.234518989526818\\
830 0.234467168796057\\
831 0.234419640619732\\
832 0.234377718387333\\
833 0.234342679332829\\
834 0.234315762658589\\
835 0.234298167759581\\
836 0.23429105254904\\
837 0.234295531886487\\
838 0.234312676108627\\
839 0.234343509663253\\
840 0.23438900984578\\
841 0.234450105637382\\
842 0.234527676642954\\
843 0.2346225521262\\
844 0.234735510138469\\
845 0.234867276737735\\
846 0.235018525294582\\
847 0.235189875882999\\
848 0.23538189475454\\
849 0.2355950938949\\
850 0.235829930662191\\
851 0.236086807506561\\
852 0.236366071770741\\
853 0.236668015570368\\
854 0.236992875751827\\
855 0.237340833924651\\
856 0.237712016565669\\
857 0.238106495192853\\
858 0.238524286607888\\
859 0.238965353207286\\
860 0.239429603362338\\
861 0.239916891868253\\
862 0.240427020462854\\
863 0.240959738415247\\
864 0.241514743185045\\
865 0.242091681152841\\
866 0.242690148422579\\
867 0.243309691696289\\
868 0.243949809221195\\
869 0.244609951808663\\
870 0.245289523923767\\
871 0.245987884843396\\
872 0.246704349879791\\
873 0.247438191665413\\
874 0.24818864149433\\
875 0.248954890715146\\
876 0.249736092170693\\
877 0.250531361680139\\
878 0.251339779559707\\
879 0.252160392178558\\
880 0.252992213546582\\
881 0.253834226930833\\
882 0.254685386497443\\
883 0.255544618976115\\
884 0.256410825344558\\
885 0.257282882530381\\
886 0.25815964512803\\
887 0.259039947128354\\
888 0.259922603658406\\
889 0.260806412729105\\
890 0.261690156988391\\
891 0.262572605477549\\
892 0.263452515388394\\
893 0.26432863381903\\
894 0.265199699525948\\
895 0.266064444670252\\
896 0.26692159655584\\
897 0.267769879357412\\
898 0.268608015836206\\
899 0.269434729041426\\
900 0.270248743995352\\
901 0.271048789360195\\
902 0.271833599084775\\
903 0.272601914029197\\
904 0.273352483565705\\
905 0.27408406715397\\
906 0.27479543588904\\
907 0.275485374020134\\
908 0.276152680438334\\
909 0.276796170131166\\
910 0.277414675602075\\
911 0.278007048253138\\
912 0.278572159729772\\
913 0.279108903226599\\
914 0.279616194753914\\
915 0.280092974364238\\
916 0.280538207338414\\
917 0.28095088533053\\
918 0.281330027470809\\
919 0.281674681425504\\
920 0.281983924412804\\
921 0.282256864173912\\
922 0.282492639898838\\
923 0.282690423106946\\
924 0.282849418482654\\
925 0.282968864666673\\
926 0.283048035003131\\
927 0.283086238242785\\
928 0.283082819202422\\
929 0.283037159380366\\
930 0.282948677527976\\
931 0.282816830177105\\
932 0.282641112123836\\
933 0.282421056869254\\
934 0.282156237018387\\
935 0.281846264638479\\
936 0.281490791577359\\
937 0.281089509741997\\
938 0.280642151336729\\
939 0.280148489060355\\
940 0.279608336261209\\
941 0.279021547049365\\
942 0.278388016365154\\
943 0.27770768000329\\
944 0.27698051459206\\
945 0.276206537527286\\
946 0.275385806861015\\
947 0.274518421145118\\
948 0.273604519230148\\
949 0.272644280019906\\
950 0.271637922182203\\
951 0.270585703816383\\
952 0.269487922078238\\
953 0.268344912763032\\
954 0.267157049847431\\
955 0.265924744991213\\
956 0.264648446999706\\
957 0.263328641247947\\
958 0.261965849067633\\
959 0.260560627097975\\
960 0.259113566601622\\
961 0.257625292746854\\
962 0.256096463857297\\
963 0.254527770630432\\
964 0.252919935326213\\
965 0.251273710927121\\
966 0.249589880271016\\
967 0.247869255158163\\
968 0.246112675433829\\
969 0.244321008047857\\
970 0.242495146092625\\
971 0.24063600782084\\
972 0.238744535644606\\
973 0.236821695117307\\
974 0.234868473899856\\
975 0.232885880712754\\
976 0.230874944275228\\
977 0.228836712232513\\
978 0.226772250072331\\
979 0.224682640031634\\
980 0.222568979994697\\
981 0.220432382383619\\
982 0.218273973042195\\
983 0.216094890114171\\
984 0.213896282917102\\
985 0.211679310813325\\
986 0.209445142079921\\
987 0.207194952779712\\
988 0.204929925635231\\
989 0.202651248907429\\
990 0.200360115280846\\
991 0.198057720757244\\
992 0.19574526355986\\
993 0.193423943050306\\
994 0.191094958659843\\
995 0.188759508836527\\
996 0.186418790009805\\
997 0.184073995574434\\
998 0.18172631489572\\
999 0.179376932337861\\
1000 0.177027026316516\\
1001 0.174677768375934\\
};
\addlegendentry{Fahrrad};

\addplot [
color=mycolor1,
solid
]
table[row sep=crcr]{
1 0.035918217727969\\
2 0.0358656318090198\\
3 0.035817687587957\\
4 0.0357743813600413\\
5 0.0357357063007631\\
6 0.0357016524151982\\
7 0.0356722064902345\\
8 0.0356473520500968\\
9 0.0356270693154944\\
10 0.0356113351666016\\
11 0.0356001231101113\\
12 0.0355934032507143\\
13 0.0355911422674366\\
14 0.0355933033952574\\
15 0.0355998464123605\\
16 0.0356107276333067\\
17 0.0356258999083366\\
18 0.0356453126289553\\
19 0.0356689117398909\\
20 0.035696639757463\\
21 0.0357284357943535\\
22 0.0357642355907239\\
23 0.0358039715515871\\
24 0.0358475727903042\\
25 0.0358949651780479\\
26 0.0359460713990427\\
27 0.0360008110113717\\
28 0.0360591005131147\\
29 0.0361208534135677\\
30 0.0361859803092747\\
31 0.0362543889645928\\
32 0.0363259843965084\\
33 0.0364006689634591\\
34 0.0364783424579993\\
35 0.0365589022032078\\
36 0.0366422431526784\\
37 0.0367282579938741\\
38 0.0368168372547374\\
39 0.0369078694136626\\
40 0.0370012410129445\\
41 0.0370968367755554\\
42 0.0371945397248593\\
43 0.0372942313068216\\
44 0.0373957915142941\\
45 0.0374990990129316\\
46 0.0376040312683003\\
47 0.0377104646738393\\
48 0.0378182746794798\\
49 0.0379273359208136\\
50 0.0380375223486933\\
51 0.0381487073590371\\
52 0.0382607639224436\\
53 0.0383735647130849\\
54 0.0384869822363251\\
55 0.0386008889545965\\
56 0.0387151574112064\\
57 0.0388296603518802\\
58 0.0389442708439319\\
59 0.0390588623929673\\
60 0.0391733090570017\\
61 0.0392874855578799\\
62 0.0394012673899304\\
63 0.0395145309258113\\
64 0.0396271535194573\\
65 0.039739013606002\\
66 0.039849990798577\\
67 0.0399599659819624\\
68 0.0400688214031291\\
69 0.0401764407588073\\
70 0.0402827092803587\\
71 0.0403875138162764\\
72 0.0404907429125339\\
73 0.0405922868909167\\
74 0.0406920379254976\\
75 0.0407898901175\\
76 0.040885739568845\\
77 0.0409794844547044\\
78 0.0410710250954167\\
79 0.0411602640281143\\
80 0.0412471060783508\\
81 0.0413314584319166\\
82 0.0414132307069438\\
83 0.0414923350263176\\
84 0.0415686860903409\\
85 0.0416422012495317\\
86 0.0417128005773801\\
87 0.0417804069428444\\
88 0.0418449460823203\\
89 0.0419063466707901\\
90 0.0419645403918242\\
91 0.0420194620060885\\
92 0.0420710494179944\\
93 0.0421192437401131\\
94 0.0421639893549712\\
95 0.04220523397384\\
96 0.0422429286921305\\
97 0.042277028041012\\
98 0.0423074900348774\\
99 0.0423342762142888\\
100 0.0423573516840518\\
101 0.0423766851460861\\
102 0.0423922489267943\\
103 0.0424040189986718\\
104 0.0424119749959438\\
105 0.0424161002240732\\
106 0.042416381663098\\
107 0.042412809964919\\
108 0.0424053794448014\\
109 0.0423940880673845\\
110 0.0423789374274181\\
111 0.0423599327253447\\
112 0.0423370827378055\\
113 0.0423103997831504\\
114 0.0422798996820299\\
115 0.0422456017131109\\
116 0.0422075285639241\\
117 0.04216570627685\\
118 0.0421201641903076\\
119 0.0420709348753105\\
120 0.0420180540676499\\
121 0.0419615605960163\\
122 0.0419014963064066\\
123 0.0418379059832007\\
124 0.0417708372673179\\
125 0.041700340571887\\
126 0.0416264689958879\\
127 0.0415492782362545\\
128 0.041468826498912\\
129 0.0413851744091356\\
130 0.0412983849214738\\
131 0.0412085232293911\\
132 0.0411156566747556\\
133 0.0410198546572832\\
134 0.0409211885439966\\
135 0.0408197315786818\\
136 0.0407155587912951\\
137 0.0406087469073164\\
138 0.0404993742571431\\
139 0.0403875206856438\\
140 0.0402732674618884\\
141 0.0401566971889484\\
142 0.040037893713621\\
143 0.0399169420359524\\
144 0.0397939282184547\\
145 0.0396689392949375\\
146 0.0395420631789458\\
147 0.0394133885718702\\
148 0.0392830048708487\\
149 0.0391510020766057\\
150 0.0390174707014236\\
151 0.0388825016774447\\
152 0.0387461862654223\\
153 0.0386086159638998\\
154 0.0384698824187218\\
155 0.0383300773327772\\
156 0.0381892923759054\\
157 0.0380476190949359\\
158 0.0379051488238605\\
159 0.037761972594169\\
160 0.0376181810454017\\
161 0.0374738643359983\\
162 0.0373291120545382\\
163 0.0371840131314871\\
164 0.0370386557515797\\
165 0.0368931272669768\\
166 0.036747514111351\\
167 0.0366019017150561\\
168 0.0364563744215474\\
169 0.0363110154052211\\
170 0.0361659065908501\\
171 0.0360211285748199\\
172 0.0358767605484124\\
173 0.0357328802233764\\
174 0.0355895637599301\\
175 0.0354468856972148\\
176 0.0353049188861771\\
177 0.0351637344248999\\
178 0.0350234015964421\\
179 0.0348839878092887\\
180 0.034745558540623\\
181 0.0346081772828078\\
182 0.0344719054935751\\
183 0.0343368025503721\\
184 0.0342029257091158\\
185 0.0340703300673931\\
186 0.0339390685319954\\
187 0.0338091917906397\\
188 0.0336807482877808\\
189 0.0335537842045421\\
190 0.033428343442913\\
191 0.0333044676143876\\
192 0.0331821960331226\\
193 0.0330615657135472\\
194 0.0329426113722743\\
195 0.0328253654341991\\
196 0.0327098580427919\\
197 0.0325961170746682\\
198 0.0324841681584355\\
199 0.0323740346976388\\
200 0.0322657378975179\\
201 0.0321592967952919\\
202 0.0320547282936963\\
203 0.0319520471974367\\
204 0.0318512662521097\\
205 0.0317523961850711\\
206 0.0316554457477752\\
207 0.0315604217592208\\
208 0.0314673291502835\\
209 0.0313761710088976\\
210 0.0312869486262466\\
211 0.0311996615442608\\
212 0.0311143076047215\\
213 0.0310308830001825\\
214 0.0309493823269312\\
215 0.030869798640395\\
216 0.0307921235135696\\
217 0.0307163470989565\\
218 0.0306424581941401\\
219 0.0305704443107294\\
220 0.030500291746118\\
221 0.0304319856574067\\
222 0.0303655101368339\\
223 0.030300848288099\\
224 0.0302379823030168\\
225 0.0301768935379932\\
226 0.0301175625898573\\
227 0.030059969370633\\
228 0.0300040931808737\\
229 0.0299499127812279\\
230 0.0298974064619361\\
231 0.0298465521100011\\
232 0.0297973272738041\\
233 0.0297497092249715\\
234 0.0297036750173285\\
235 0.0296592015428003\\
236 0.0296162655841534\\
237 0.0295748438644861\\
238 0.0295349130934087\\
239 0.029496450009867\\
240 0.0294594314215936\\
241 0.0294238342412111\\
242 0.0293896355191007\\
243 0.0293568124732626\\
244 0.0293253425164889\\
245 0.0292952032811887\\
246 0.0292663726421387\\
247 0.0292388287373191\\
248 0.0292125499868815\\
249 0.0291875151102598\\
250 0.0291637031414959\\
251 0.0291410934429357\\
252 0.0291196657174625\\
253 0.0290994000193159\\
254 0.0290802767633722\\
255 0.0290622767326506\\
256 0.0290453810838238\\
257 0.0290295713506205\\
258 0.0290148294451004\\
259 0.0290011376568455\\
260 0.0289884786501467\\
261 0.0289768354592882\\
262 0.0289661914820038\\
263 0.0289565304711262\\
264 0.0289478365244064\\
265 0.0289400940724794\\
266 0.0289332878649791\\
267 0.0289274029548882\\
268 0.0289224246813308\\
269 0.0289183386511265\\
270 0.0289151307194612\\
271 0.0289127869699809\\
272 0.028911293694458\\
273 0.0289106373720388\\
274 0.0289108046480671\\
275 0.0289117823125551\\
276 0.0289135572784321\\
277 0.0289161165597316\\
278 0.0289194472499066\\
279 0.0289235365004995\\
280 0.02892837150041\\
281 0.0289339394559893\\
282 0.0289402275721304\\
283 0.0289472230344473\\
284 0.0289549129925703\\
285 0.0289632845445506\\
286 0.0289723247223617\\
287 0.028982020478503\\
288 0.0289923586737351\\
289 0.029003326066\\
290 0.029014909300595\\
291 0.0290270949016895\\
292 0.0290398692652808\\
293 0.0290532186536984\\
294 0.0290671291917699\\
295 0.0290815868647689\\
296 0.0290965775182624\\
297 0.0291120868599806\\
298 0.0291281004638255\\
299 0.0291446037761331\\
300 0.0291615821242995\\
301 0.0291790207278723\\
302 0.0291969047122037\\
303 0.0292152191247496\\
304 0.0292339489540747\\
305 0.0292530791515758\\
306 0.0292725946558792\\
307 0.0292924804198554\\
308 0.0293127214402659\\
309 0.0293333027901725\\
310 0.0293542096543495\\
311 0.0293754273679321\\
312 0.0293969414583792\\
313 0.0294187376905726\\
314 0.0294408021146554\\
315 0.0294631211161212\\
316 0.0294856814676365\\
317 0.0295084703820177\\
318 0.0295314755656712\\
319 0.0295546852717488\\
320 0.0295780883523357\\
321 0.0296016743091278\\
322 0.0296254333421662\\
323 0.0296493563962243\\
324 0.0296734352044633\\
325 0.0296976623290406\\
326 0.029722031198409\\
327 0.029746536141015\\
328 0.0297711724150588\\
329 0.0297959362340255\\
330 0.0298208247877991\\
331 0.0298458362592196\\
332 0.0298709698359753\\
333 0.0298962257177849\\
334 0.0299216051188491\\
335 0.02994711026551\\
336 0.0299727443890963\\
337 0.0299985117141567\\
338 0.0300244174425295\\
339 0.0300504677337655\\
340 0.030076669682279\\
341 0.0301030312914082\\
342 0.0301295614444458\\
343 0.030156269872674\\
344 0.0301831671204513\\
345 0.0302102645074376\\
346 0.0302375740881042\\
347 0.0302651086087348\\
348 0.030292881462152\\
349 0.0303209066404115\\
350 0.0303491986857092\\
351 0.0303777726397279\\
352 0.0304066439916079\\
353 0.0304358286246764\\
354 0.0304653427620414\\
355 0.030495202911143\\
356 0.0305254258073485\\
357 0.0305560283566691\\
358 0.0305870275776732\\
359 0.0306184405426637\\
360 0.030650284318187\\
361 0.0306825759049337\\
362 0.0307153321770943\\
363 0.0307485698212284\\
364 0.0307823052747053\\
365 0.030816554663777\\
366 0.0308513337413394\\
367 0.0308866578244439\\
368 0.0309225417316185\\
369 0.0309589997200596\\
370 0.0309960454227586\\
371 0.0310336917856278\\
372 0.0310719510046909\\
373 0.0311108344634176\\
374 0.0311503526703313\\
375 0.0311905151971486\\
376 0.0312313306178609\\
377 0.0312728064492931\\
378 0.0313149490937767\\
379 0.0313577637846619\\
380 0.0314012545353718\\
381 0.0314454240925475\\
382 0.0314902738936202\\
383 0.0315358040289625\\
384 0.0315820132086039\\
385 0.0316288987333418\\
386 0.0316764564699612\\
387 0.0317246808302456\\
388 0.0317735647535392\\
389 0.0318230996927731\\
390 0.0318732756040037\\
391 0.0319240809395538\\
392 0.0319755026448193\\
393 0.0320275261588398\\
394 0.0320801354188526\\
395 0.0321333128690903\\
396 0.0321870394739504\\
397 0.0322412947354403\\
398 0.0322960567146327\\
399 0.0323513020568046\\
400 0.0324070060199404\\
401 0.0324631425062927\\
402 0.0325196840967718\\
403 0.0325766020881067\\
404 0.0326338665328681\\
405 0.0326914462824036\\
406 0.0327493090325663\\
407 0.0328074213719951\\
408 0.0328657488326995\\
409 0.0329242559427749\\
410 0.0329829062811699\\
411 0.0330416625345012\\
412 0.033100486555884\\
413 0.0331593394256237\\
414 0.0332181815135198\\
415 0.0332769725425855\\
416 0.0333356716541866\\
417 0.0333942374748362\\
418 0.033452628184957\\
419 0.0335108015897925\\
420 0.0335687151924395\\
421 0.0336263262689231\\
422 0.0336835919453319\\
423 0.0337404692771053\\
424 0.0337969153305454\\
425 0.0338528872665418\\
426 0.0339083424264088\\
427 0.0339632384196487\\
428 0.0340175332133884\\
429 0.0340711852231705\\
430 0.0341241534047313\\
431 0.0341763973463528\\
432 0.0342278773613414\\
433 0.034278554580158\\
434 0.0343283910417028\\
435 0.0343773497832433\\
436 0.0344253949284666\\
437 0.0344724917731313\\
438 0.0345186068677986\\
439 0.0345637080971269\\
440 0.0346077647552478\\
441 0.0346507476168007\\
442 0.0346926290033093\\
443 0.0347333828446546\\
444 0.0347729847353971\\
445 0.0348114119856249\\
446 0.0348486436659687\\
447 0.0348846606465237\\
448 0.034919445629628\\
449 0.0349529831766445\\
450 0.0349852597290063\\
451 0.0350162636238088\\
452 0.035045985104216\\
453 0.0350744163249335\\
454 0.0351015513530151\\
455 0.035127386164393\\
456 0.0351519186367048\\
457 0.0351751485390536\\
458 0.0351970775191524\\
459 0.0352177090880146\\
460 0.0352370486021457\\
461 0.0352551032431368\\
462 0.0352718819945883\\
463 0.0352873956162679\\
464 0.0353016566152996\\
465 0.0353146792141142\\
466 0.0353264793149365\\
467 0.0353370744606747\\
468 0.0353464837921681\\
469 0.0353547280018926\\
470 0.0353618292844005\\
471 0.0353678112838755\\
472 0.0353726990392298\\
473 0.0353765189272296\\
474 0.035379298604227\\
475 0.0353810669470865\\
476 0.035381853993779\\
477 0.0353816908839305\\
478 0.0353806097995304\\
479 0.0353786439061347\\
480 0.0353758272950925\\
481 0.0353721949273295\\
482 0.0353677825789888\\
483 0.0353626267889726\\
484 0.035356764808296\\
485 0.0353502345511341\\
486 0.0353430745474472\\
487 0.0353353238970757\\
488 0.0353270222251969\\
489 0.0353182096390445\\
490 0.0353089266857914\\
491 0.0352992143115036\\
492 0.0352891138210774\\
493 0.0352786668390749\\
494 0.0352679152713787\\
495 0.0352569012675918\\
496 0.0352456671841128\\
497 0.0352342555478195\\
498 0.035222709020303\\
499 0.0352110703625929\\
500 0.0351993824003256\\
501 0.0351876879893055\\
502 0.0351760299814129\\
503 0.0351644511907944\\
504 0.0351529943602288\\
505 0.0351417021275144\\
506 0.0351306169917144\\
507 0.0351197812791395\\
508 0.0351092371090281\\
509 0.0350990263589166\\
510 0.0350891906296978\\
511 0.035079771210404\\
512 0.035070809042843\\
513 0.03506234468629\\
514 0.0350544182824536\\
515 0.0350470695209153\\
516 0.0350403376052345\\
517 0.0350342612198902\\
518 0.0350288784982001\\
519 0.035024226991306\\
520 0.0350203436382854\\
521 0.0350172647374234\\
522 0.0350150259186475\\
523 0.0350136621170815\\
524 0.0350132075476142\\
525 0.0350136956803399\\
526 0.0350151592167438\\
527 0.0350176300666029\\
528 0.0350211393256976\\
529 0.0350257172545162\\
530 0.0350313932581596\\
531 0.0350381958676976\\
532 0.0350461527232905\\
533 0.0350552905593962\\
534 0.0350656351923258\\
535 0.0350772115103693\\
536 0.0350900434667028\\
537 0.0351041540752924\\
538 0.0351195654099901\\
539 0.0351362986069898\\
540 0.035154373870701\\
541 0.0351738104829617\\
542 0.0351946268155007\\
543 0.0352168403456554\\
544 0.0352404676753513\\
545 0.035265524553169\\
546 0.0352920258991323\\
547 0.0353199858318102\\
548 0.0353494176974186\\
549 0.0353803341007011\\
550 0.0354127469373371\\
551 0.0354466674274873\\
552 0.0354821061499679\\
553 0.0355190730765168\\
554 0.0355575776056505\\
555 0.0355976285956627\\
556 0.0356392343963675\\
557 0.0356824028792393\\
558 0.0357271414656446\\
559 0.0357734571529077\\
560 0.0358213565379881\\
561 0.0358708458385889\\
562 0.0359219309115473\\
563 0.0359746172683931\\
564 0.0360289100879904\\
565 0.0360848142262068\\
566 0.036142334222579\\
567 0.0362014743039698\\
568 0.0362622383852336\\
569 0.0363246300669258\\
570 0.0363886526301119\\
571 0.0364543090283458\\
572 0.0365216018768836\\
573 0.0365905334391755\\
574 0.0366611056106393\\
575 0.0367333198997186\\
576 0.0368071774062724\\
577 0.0368826787973538\\
578 0.0369598242804086\\
579 0.0370386135739889\\
580 0.0371190458762285\\
581 0.037201119831451\\
582 0.0372848334953061\\
583 0.037370184298814\\
584 0.0374571690117282\\
585 0.0375457837056253\\
586 0.0376360237169827\\
587 0.0377278836102177\\
588 0.0378213571404543\\
589 0.0379164372157337\\
590 0.0380131158584383\\
591 0.0381113841657696\\
592 0.0382112322691999\\
593 0.038312649292941\\
594 0.0384156233116096\\
595 0.0385201413073488\\
596 0.0386261891265994\\
597 0.0387337514365992\\
598 0.0388428116816806\\
599 0.038953352039591\\
600 0.0390653533782748\\
601 0.0391787952136796\\
602 0.0392936556690874\\
603 0.0394099114362814\\
604 0.0395275377386839\\
605 0.0396465082965314\\
606 0.0397667952941468\\
607 0.0398883693493703\\
608 0.040011199485171\\
609 0.0401352531033843\\
610 0.04026049596043\\
611 0.0403868921448574\\
612 0.0405144040566873\\
613 0.0406429923887246\\
614 0.0407726161101843\\
615 0.0409032324529951\\
616 0.0410347969010381\\
617 0.0411672631824235\\
618 0.0413005832648273\\
619 0.0414347073539331\\
620 0.0415695838950502\\
621 0.0417051595779489\\
622 0.041841379344932\\
623 0.041978186402198\\
624 0.0421155222345511\\
625 0.0422533266234157\\
626 0.0423915376679838\\
627 0.0425300918093124\\
628 0.0426689238573508\\
629 0.0428079670211412\\
630 0.0429471529426164\\
631 0.0430864117344243\\
632 0.0432256720220993\\
633 0.0433648609907667\\
634 0.0435039044364645\\
635 0.0436427268220727\\
636 0.0437812513377602\\
637 0.0439193999657899\\
638 0.0440570935494595\\
639 0.0441942518659019\\
640 0.0443307937024251\\
641 0.0444666369360279\\
642 0.0446016986157021\\
643 0.0447358950470997\\
644 0.0448691418791278\\
645 0.0450013541920167\\
646 0.0451324465864022\\
647 0.0452623332729773\\
648 0.0453909281623546\\
649 0.0455181449549607\\
650 0.0456438972309931\\
651 0.0457680985405618\\
652 0.0458906624940935\\
653 0.0460115028530797\\
654 0.0461305336213846\\
655 0.0462476691374929\\
656 0.0463628241680659\\
657 0.0464759140029423\\
658 0.046586854551467\\
659 0.0466955624399199\\
660 0.0468019551098479\\
661 0.0469059509171752\\
662 0.0470074692320178\\
663 0.047106430539109\\
664 0.0472027565386883\\
665 0.0472963702476353\\
666 0.0473871961005711\\
667 0.0474751600505925\\
668 0.0475601896692448\\
669 0.0476422142452493\\
670 0.0477211648814379\\
671 0.0477969745894418\\
672 0.0478695783819241\\
673 0.0479389133623687\\
674 0.0480049188125068\\
675 0.0480675362774244\\
676 0.0481267096483573\\
677 0.0481823852431311\\
678 0.0482345118840809\\
679 0.0482830409731157\\
680 0.0483279265635455\\
681 0.0483691254284605\\
682 0.0484065971257508\\
683 0.0484403040600778\\
684 0.0484702115420798\\
685 0.0484962878448624\\
686 0.0485185042575438\\
687 0.0485368351354372\\
688 0.0485512579463522\\
689 0.0485617533124571\\
690 0.048568305047161\\
691 0.0485709001865495\\
692 0.0485695290149972\\
693 0.048564185084674\\
694 0.0485548652287409\\
695 0.0485415695681073\\
696 0.0485243015116907\\
697 0.048503067750178\\
698 0.0484778782433467\\
699 0.0484487462010486\\
700 0.0484156880580081\\
701 0.048378723442621\\
702 0.0483378751399767\\
703 0.0482931690493544\\
704 0.04824463413647\\
705 0.0481923023807708\\
706 0.0481362087180929\\
707 0.0480763909790122\\
708 0.0480128898232272\\
709 0.0479457486703236\\
710 0.0478750136272743\\
711 0.0478007334130354\\
712 0.047722959280633\\
713 0.0476417449372291\\
714 0.0475571464628137\\
715 0.0474692222282536\\
716 0.0473780328132891\\
717 0.047283640924786\\
718 0.0471861113153817\\
719 0.0470855107027262\\
720 0.0469819076897011\\
721 0.0468753726861939\\
722 0.0467659778331265\\
723 0.0466537969293997\\
724 0.0465389053622188\\
725 0.0464213800410347\\
726 0.0463012993351171\\
727 0.0461787430145447\\
728 0.0460537921941769\\
729 0.0459265292800239\\
730 0.0457970379174586\\
731 0.0456654029408689\\
732 0.045531710324537\\
733 0.0453960471346158\\
734 0.0452585014820381\\
735 0.0451191624761442\\
736 0.0449781201788447\\
737 0.0448354655592267\\
738 0.0446912904486738\\
739 0.0445456874967535\\
740 0.0443987501282581\\
741 0.0442505725017586\\
742 0.0441012494699038\\
743 0.043950876541588\\
744 0.0437995498460779\\
745 0.0436473660992676\\
746 0.0434944225723346\\
747 0.0433408170630342\\
748 0.0431866478696208\\
749 0.0430320137670571\\
750 0.0428770139849107\\
751 0.0427217481861428\\
752 0.0425663164459617\\
753 0.0424108192300825\\
754 0.0422553573719732\\
755 0.0421000320488031\\
756 0.0419449447558146\\
757 0.0417901972787832\\
758 0.0416358916641551\\
759 0.0414821301864352\\
760 0.0413290153124233\\
761 0.0411766496619647\\
762 0.0410251359649385\\
763 0.0408745770142697\\
764 0.0407250756148041\\
765 0.0405767345279366\\
766 0.0404296564119307\\
767 0.0402839437579107\\
768 0.0401396988215479\\
769 0.0399970235505001\\
770 0.0398560195076967\\
771 0.0397167877905931\\
772 0.0395794289465456\\
773 0.0394440428844851\\
774 0.0393107287830891\\
775 0.0391795849956714\\
776 0.0390507089520296\\
777 0.0389241970575078\\
778 0.0388001445895768\\
779 0.0386786455923355\\
780 0.0385597927694659\\
781 0.0384436773762285\\
782 0.0383303891109931\\
783 0.0382200160066607\\
784 0.0381126443223108\\
785 0.0380083584356265\\
786 0.0379072407370487\\
787 0.0378093715269425\\
788 0.0377148289170164\\
789 0.0376236887368551\\
790 0.037536024446136\\
791 0.0374519070531442\\
792 0.0373714050403389\\
793 0.0372945842976293\\
794 0.0372215080637285\\
795 0.0371522368756899\\
796 0.0370868285265241\\
797 0.0370253380306193\\
798 0.0369678175965017\\
799 0.0369143166062061\\
800 0.0368648816002574\\
801 0.0368195562671386\\
802 0.0367783814361803\\
803 0.0367413950729932\\
804 0.0367086322767976\\
805 0.0366801252792016\\
806 0.0366559034441171\\
807 0.0366359932686039\\
808 0.036620418384513\\
809 0.0366091995607929\\
810 0.0366023547062244\\
811 0.0365998988722268\\
812 0.036601844255321\\
813 0.0366082001988573\\
814 0.0366189731936782\\
815 0.0366341668774536\\
816 0.0366537820324829\\
817 0.0366778165818482\\
818 0.036706265583928\\
819 0.0367391212254052\\
820 0.0367763728129339\\
821 0.0368180067635892\\
822 0.036864006594167\\
823 0.0369143529093438\\
824 0.0369690233886746\\
825 0.0370279927724275\\
826 0.0370912328463247\\
827 0.0371587124253517\\
828 0.0372303973368899\\
829 0.0373062504035012\\
830 0.0373862314257632\\
831 0.0374702971656055\\
832 0.0375584013306441\\
833 0.0376504945600441\\
834 0.0377465244124672\\
835 0.0378464353566796\\
836 0.0379501687654023\\
837 0.0380576629129946\\
838 0.0381688529775532\\
839 0.0382836710480035\\
840 0.0384020461367453\\
841 0.0385239041983964\\
842 0.0386491681551543\\
843 0.038777757929264\\
844 0.0389095904830002\\
845 0.0390445798663722\\
846 0.0391826372724066\\
847 0.0393236710994774\\
848 0.0394675870199221\\
849 0.039614288054172\\
850 0.0397636746497352\\
851 0.0399156447644387\\
852 0.0400700939533588\\
853 0.040226915458924\\
854 0.0403860003038177\\
855 0.0405472373864612\\
856 0.0407105135789556\\
857 0.0408757138274181\\
858 0.0410427212546662\\
859 0.0412114172651633\\
860 0.0413816816520119\\
861 0.0415533927056395\\
862 0.0417264273237729\\
863 0.0419006611224226\\
864 0.0420759685478399\\
865 0.0422522229896002\\
866 0.0424292968949634\\
867 0.0426070618845426\\
868 0.0427853888692008\\
869 0.0429641481680896\\
870 0.0431432096278069\\
871 0.0433224427427225\\
872 0.0435017167765017\\
873 0.0436809008847146\\
874 0.0438598642382041\\
875 0.0440384761467218\\
876 0.0442166061823051\\
877 0.0443941243019414\\
878 0.0445709009691618\\
879 0.0447468072742846\\
880 0.0449217150530806\\
881 0.0450954970036785\\
882 0.0452680268015438\\
883 0.0454391792123247\\
884 0.0456088302022946\\
885 0.0457768570461475\\
886 0.0459431384320187\\
887 0.0461075545637199\\
888 0.0462699872602038\\
889 0.0464303200522363\\
890 0.0465884382761679\\
891 0.0467442291645413\\
892 0.0468975819330894\\
893 0.0470483878635436\\
894 0.0471965403816597\\
895 0.0473419351299381\\
896 0.0474844700346225\\
897 0.0476240453666574\\
898 0.0477605637963807\\
899 0.0478939304418076\\
900 0.0480240529104422\\
901 0.0481508413346151\\
902 0.0482742084004134\\
903 0.0483940693703168\\
904 0.048510342099708\\
905 0.0486229470474634\\
906 0.0487318072808706\\
907 0.0488368484751465\\
908 0.0489379989078617\\
909 0.0490351894485951\\
910 0.0491283535441604\\
911 0.0492174271997556\\
912 0.0493023489563937\\
913 0.0493830598650027\\
914 0.0494595034576777\\
915 0.0495316257167084\\
916 0.0495993750420834\\
917 0.0496627022180949\\
918 0.0497215603795245\\
919 0.0497759049778623\\
920 0.0498256937480925\\
921 0.0498708866766244\\
922 0.0499114459708832\\
923 0.04994733603096\\
924 0.049978523423624\\
925 0.0500049768589449\\
926 0.0500266671697699\\
927 0.0500435672943259\\
928 0.0500556522622299\\
929 0.0500628991841572\\
930 0.0500652872453382\\
931 0.0500627977029875\\
932 0.050055413887744\\
933 0.0500431212091832\\
934 0.0500259071654182\\
935 0.0500037613567093\\
936 0.049976675502898\\
937 0.0499446434644386\\
938 0.0499076612668063\\
939 0.0498657271280854\\
940 0.0498188414895116\\
941 0.0497670070486731\\
942 0.0497102287949528\\
943 0.0496485140466958\\
944 0.0495818724895905\\
945 0.0495103162159313\\
946 0.049433859764685\\
947 0.0493525201624598\\
948 0.0492663169655055\\
949 0.049175272302803\\
950 0.0490794109201942\\
951 0.0489787602253887\\
952 0.0488733503335722\\
953 0.0487632141132565\\
954 0.0486483872319369\\
955 0.048528908201067\\
956 0.0484048184198529\\
957 0.0482761622173902\\
958 0.0481429868927243\\
959 0.0480053427524685\\
960 0.0478632831456719\\
961 0.0477168644956669\\
962 0.0475661463286727\\
963 0.0474111912989779\\
964 0.0472520652105837\\
965 0.047088837035191\\
966 0.046921578926373\\
967 0.0467503662297988\\
968 0.0465752774894701\\
969 0.0463963944499969\\
970 0.0462138020548861\\
971 0.0460275884406601\\
972 0.0458378449264445\\
973 0.0456446659985099\\
974 0.0454481492891641\\
975 0.0452483955493723\\
976 0.0450455086145282\\
977 0.0448395953628804\\
978 0.0446307656662109\\
979 0.0444191323324507\\
980 0.0442048110399996\\
981 0.0439879202635956\\
982 0.0437685811916474\\
983 0.0435469176350128\\
984 0.0433230559272649\\
985 0.043097124816544\\
986 0.0428692553491455\\
987 0.0426395807450424\\
988 0.0424082362655805\\
989 0.0421753590736271\\
990 0.0419410880864872\\
991 0.0417055638219333\\
992 0.0414689282377221\\
993 0.0412313245649989\\
994 0.0409928971360083\\
995 0.0407537912065507\\
996 0.0405141527736507\\
997 0.0402741283889493\\
998 0.0400338649684424\\
999 0.0397935095993341\\
1000 0.039553209344906\\
1001 0.039313111048334\\
};
\addlegendentry{Auto};

\end{axis}
\end{tikzpicture}% 
\end{frame}



\subsection{Entscheidungsbaum}
\begin{frame}{Entscheidungsbaum für Bewegungsmuster}
    \centering 
    \setlength\figureheight{0.45\textwidth} 
    \setlength\figurewidth{0.8\textwidth} 
    % This file was created by matlab2tikz v0.4.4 running on MATLAB 7.14.
% Copyright (c) 2008--2013, Nico Schlömer <nico.schloemer@gmail.com>
% All rights reserved.
% 
% The latest updates can be retrieved from
%   http://www.mathworks.com/matlabcentral/fileexchange/22022-matlab2tikz
% where you can also make suggestions and rate matlab2tikz.
% 
\begin{tikzpicture}

\begin{axis}[%
width=\figurewidth,
height=\figureheight,
scale only axis,
xmin=0,
xmax=1,
ymin=0,
ymax=1,
hide axis
]
\addplot [
color=ohmred,
only marks,
mark=o,
mark options={solid},
forget plot
]
table[row sep=crcr]{
0.5 0.956521739130435\\
0.0265486725663717 0.217391304347826\\
0.52212389380531 0.91304347826087\\
0.0221238938053097 0.173913043478261\\
0.0442477876106195 0.0434782608695652\\
0.482300884955752 0.869565217391304\\
0.95575221238938 0.260869565217391\\
0.00884955752212389 0.0434782608695652\\
0.0265486725663717 0.130434782608696\\
0.47787610619469 0.826086956521739\\
0.911504424778761 0.0434782608695652\\
0.946902654867257 0.217391304347826\\
0.986725663716814 0.0869565217391304\\
0.0221238938053097 0.0869565217391304\\
0.0353982300884956 0.0434782608695652\\
0.265486725663717 0.782608695652174\\
0.694690265486726 0.739130434782609\\
0.924778761061947 0.0869565217391304\\
0.95575221238938 0.173913043478261\\
0.982300884955752 0.0434782608695652\\
0.991150442477876 0.0434782608695652\\
0.0176991150442478 0.0434782608695652\\
0.0265486725663717 0.0434782608695652\\
0.261061946902655 0.739130434782609\\
0.47787610619469 0.0434782608695652\\
0.486725663716814 0.0434782608695652\\
0.699115044247788 0.695652173913043\\
0.920353982300885 0.0434782608695652\\
0.929203539823009 0.0434782608695652\\
0.942477876106195 0.0869565217391304\\
0.964601769911504 0.130434782608696\\
0.185840707964602 0.695652173913043\\
0.398230088495575 0.347826086956522\\
0.561946902654867 0.434782608695652\\
0.769911504424779 0.652173913043478\\
0.938053097345133 0.0434782608695652\\
0.946902654867257 0.0434782608695652\\
0.960176991150442 0.0869565217391304\\
0.973451327433628 0.0434782608695652\\
0.115044247787611 0.652173913043478\\
0.252212389380531 0.391304347826087\\
0.358407079646018 0.304347826086957\\
0.433628318584071 0.304347826086957\\
0.548672566371681 0.391304347826087\\
0.619469026548673 0.130434782608696\\
0.63716814159292 0.0434782608695652\\
0.774336283185841 0.608695652173913\\
0.95575221238938 0.0434782608695652\\
0.964601769911504 0.0434782608695652\\
0.0530973451327434 0.0434782608695652\\
0.119469026548673 0.608695652173913\\
0.20353982300885 0.217391304347826\\
0.274336283185841 0.347826086956522\\
0.327433628318584 0.0434782608695652\\
0.36283185840708 0.260869565217391\\
0.429203539823009 0.260869565217391\\
0.469026548672566 0.0434782608695652\\
0.513274336283186 0.217391304347826\\
0.570796460176991 0.347826086956522\\
0.610619469026549 0.0434782608695652\\
0.623893805309734 0.0869565217391304\\
0.668141592920354 0.217391304347826\\
0.800884955752212 0.565217391304348\\
0.0619469026548673 0.0434782608695652\\
0.123893805309735 0.565217391304348\\
0.185840707964602 0.0434782608695652\\
0.207964601769911 0.173913043478261\\
0.243362831858407 0.173913043478261\\
0.292035398230089 0.304347826086957\\
0.358407079646018 0.217391304347826\\
0.389380530973451 0.0434782608695652\\
0.415929203539823 0.217391304347826\\
0.451327433628319 0.130434782608696\\
0.508849557522124 0.173913043478261\\
0.530973451327434 0.0434782608695652\\
0.566371681415929 0.304347826086957\\
0.601769911504425 0.0434782608695652\\
0.619469026548673 0.0434782608695652\\
0.628318584070796 0.0434782608695652\\
0.65929203539823 0.173913043478261\\
0.685840707964602 0.0869565217391304\\
0.787610619469027 0.521739130434783\\
0.893805309734513 0.130434782608696\\
0.119469026548673 0.521739130434783\\
0.176991150442478 0.0434782608695652\\
0.20353982300885 0.130434782608696\\
0.221238938053097 0.0434782608695652\\
0.230088495575221 0.0434782608695652\\
0.247787610619469 0.130434782608696\\
0.265486725663717 0.0434782608695652\\
0.29646017699115 0.260869565217391\\
0.336283185840708 0.0434782608695652\\
0.36283185840708 0.173913043478261\\
0.411504424778761 0.173913043478261\\
0.433628318584071 0.0434782608695652\\
0.446902654867257 0.0869565217391304\\
0.460176991150442 0.0434782608695652\\
0.495575221238938 0.0434782608695652\\
0.513274336283186 0.130434782608696\\
0.561946902654867 0.260869565217391\\
0.592920353982301 0.0434782608695652\\
0.654867256637168 0.130434782608696\\
0.672566371681416 0.0434782608695652\\
0.68141592920354 0.0434782608695652\\
0.690265486725664 0.0434782608695652\\
0.783185840707965 0.478260869565217\\
0.876106194690266 0.0434782608695652\\
0.884955752212389 0.0434782608695652\\
0.898230088495575 0.0869565217391304\\
0.115044247787611 0.478260869565217\\
0.168141592920354 0.0434782608695652\\
0.194690265486726 0.0434782608695652\\
0.207964601769911 0.0869565217391304\\
0.243362831858407 0.0869565217391304\\
0.256637168141593 0.0434782608695652\\
0.292035398230089 0.217391304347826\\
0.31858407079646 0.0434782608695652\\
0.353982300884956 0.130434782608696\\
0.376106194690265 0.0869565217391304\\
0.407079646017699 0.130434782608696\\
0.424778761061947 0.0434782608695652\\
0.442477876106195 0.0434782608695652\\
0.451327433628319 0.0434782608695652\\
0.508849557522124 0.0869565217391304\\
0.52212389380531 0.0434782608695652\\
0.539823008849557 0.0434782608695652\\
0.566371681415929 0.217391304347826\\
0.650442477876106 0.0869565217391304\\
0.663716814159292 0.0434782608695652\\
0.752212389380531 0.434782608695652\\
0.84070796460177 0.304347826086957\\
0.893805309734513 0.0434782608695652\\
0.902654867256637 0.0434782608695652\\
0.110619469026549 0.434782608695652\\
0.15929203539823 0.0434782608695652\\
0.20353982300885 0.0434782608695652\\
0.212389380530973 0.0434782608695652\\
0.238938053097345 0.0434782608695652\\
0.247787610619469 0.0434782608695652\\
0.287610619469027 0.173913043478261\\
0.309734513274336 0.0434782608695652\\
0.349557522123894 0.0869565217391304\\
0.36283185840708 0.0434782608695652\\
0.371681415929204 0.0434782608695652\\
0.380530973451327 0.0434782608695652\\
0.402654867256637 0.0869565217391304\\
0.415929203539823 0.0434782608695652\\
0.504424778761062 0.0434782608695652\\
0.513274336283186 0.0434782608695652\\
0.548672566371681 0.0434782608695652\\
0.570796460176991 0.173913043478261\\
0.646017699115044 0.0434782608695652\\
0.654867256637168 0.0434782608695652\\
0.738938053097345 0.391304347826087\\
0.79646017699115 0.130434782608696\\
0.814159292035398 0.0434782608695652\\
0.845132743362832 0.260869565217391\\
0.106194690265487 0.391304347826087\\
0.150442477876106 0.0434782608695652\\
0.283185840707965 0.130434782608696\\
0.300884955752212 0.0434782608695652\\
0.345132743362832 0.0434782608695652\\
0.353982300884956 0.0434782608695652\\
0.398230088495575 0.0434782608695652\\
0.407079646017699 0.0434782608695652\\
0.557522123893805 0.0434782608695652\\
0.575221238938053 0.130434782608696\\
0.734513274336283 0.347826086956522\\
0.778761061946903 0.0434782608695652\\
0.787610619469027 0.0434782608695652\\
0.800884955752212 0.0869565217391304\\
0.823008849557522 0.0434782608695652\\
0.849557522123894 0.217391304347826\\
0.101769911504425 0.347826086956522\\
0.141592920353982 0.0434782608695652\\
0.278761061946903 0.0869565217391304\\
0.292035398230089 0.0434782608695652\\
0.566371681415929 0.0434782608695652\\
0.579646017699115 0.0869565217391304\\
0.730088495575221 0.304347826086957\\
0.769911504424779 0.0434782608695652\\
0.79646017699115 0.0434782608695652\\
0.805309734513274 0.0434782608695652\\
0.831858407079646 0.0434782608695652\\
0.853982300884956 0.173913043478261\\
0.0973451327433628 0.304347826086957\\
0.132743362831858 0.0434782608695652\\
0.274336283185841 0.0434782608695652\\
0.283185840707965 0.0434782608695652\\
0.575221238938053 0.0434782608695652\\
0.584070796460177 0.0434782608695652\\
0.725663716814159 0.260869565217391\\
0.761061946902655 0.0434782608695652\\
0.849557522123894 0.130434782608696\\
0.867256637168142 0.0434782608695652\\
0.0929203539823009 0.260869565217391\\
0.123893805309735 0.0434782608695652\\
0.721238938053097 0.217391304347826\\
0.752212389380531 0.0434782608695652\\
0.845132743362832 0.0869565217391304\\
0.858407079646018 0.0434782608695652\\
0.0884955752212389 0.217391304347826\\
0.115044247787611 0.0434782608695652\\
0.712389380530973 0.173913043478261\\
0.738938053097345 0.0869565217391304\\
0.84070796460177 0.0434782608695652\\
0.849557522123894 0.0434782608695652\\
0.084070796460177 0.173913043478261\\
0.106194690265487 0.0434782608695652\\
0.707964601769911 0.130434782608696\\
0.725663716814159 0.0434782608695652\\
0.734513274336283 0.0434782608695652\\
0.743362831858407 0.0434782608695652\\
0.079646017699115 0.130434782608696\\
0.0973451327433628 0.0434782608695652\\
0.70353982300885 0.0869565217391304\\
0.716814159292035 0.0434782608695652\\
0.0752212389380531 0.0869565217391304\\
0.0884955752212389 0.0434782608695652\\
0.699115044247788 0.0434782608695652\\
0.707964601769911 0.0434782608695652\\
0.0707964601769911 0.0434782608695652\\
0.079646017699115 0.0434782608695652\\
};
\addplot [
color=ohmred,
solid,
forget plot
]
table[row sep=crcr]{
0.0265486725663717 0.217391304347826\\
0.5 0.956521739130435\\
};
\addplot [
color=ohmred,
solid,
forget plot
]
table[row sep=crcr]{
0.52212389380531 0.91304347826087\\
0.5 0.956521739130435\\
};
\addplot [
color=ohmred,
solid,
forget plot
]
table[row sep=crcr]{
0.0221238938053097 0.173913043478261\\
0.0265486725663717 0.217391304347826\\
};
\addplot [
color=ohmred,
solid,
forget plot
]
table[row sep=crcr]{
0.0442477876106195 0.0434782608695652\\
0.0265486725663717 0.217391304347826\\
};
\addplot [
color=ohmred,
solid,
forget plot
]
table[row sep=crcr]{
0.482300884955752 0.869565217391304\\
0.52212389380531 0.91304347826087\\
};
\addplot [
color=ohmred,
solid,
forget plot
]
table[row sep=crcr]{
0.95575221238938 0.260869565217391\\
0.52212389380531 0.91304347826087\\
};
\addplot [
color=ohmred,
solid,
forget plot
]
table[row sep=crcr]{
0.00884955752212389 0.0434782608695652\\
0.0221238938053097 0.173913043478261\\
};
\addplot [
color=ohmred,
solid,
forget plot
]
table[row sep=crcr]{
0.0265486725663717 0.130434782608696\\
0.0221238938053097 0.173913043478261\\
};
\addplot [
color=ohmred,
solid,
forget plot
]
table[row sep=crcr]{
0.47787610619469 0.826086956521739\\
0.482300884955752 0.869565217391304\\
};
\addplot [
color=ohmred,
solid,
forget plot
]
table[row sep=crcr]{
0.911504424778761 0.0434782608695652\\
0.482300884955752 0.869565217391304\\
};
\addplot [
color=ohmred,
solid,
forget plot
]
table[row sep=crcr]{
0.946902654867257 0.217391304347826\\
0.95575221238938 0.260869565217391\\
};
\addplot [
color=ohmred,
solid,
forget plot
]
table[row sep=crcr]{
0.986725663716814 0.0869565217391304\\
0.95575221238938 0.260869565217391\\
};
\addplot [
color=ohmred,
solid,
forget plot
]
table[row sep=crcr]{
0.0221238938053097 0.0869565217391304\\
0.0265486725663717 0.130434782608696\\
};
\addplot [
color=ohmred,
solid,
forget plot
]
table[row sep=crcr]{
0.0353982300884956 0.0434782608695652\\
0.0265486725663717 0.130434782608696\\
};
\addplot [
color=ohmred,
solid,
forget plot
]
table[row sep=crcr]{
0.265486725663717 0.782608695652174\\
0.47787610619469 0.826086956521739\\
};
\addplot [
color=ohmred,
solid,
forget plot
]
table[row sep=crcr]{
0.694690265486726 0.739130434782609\\
0.47787610619469 0.826086956521739\\
};
\addplot [
color=ohmred,
solid,
forget plot
]
table[row sep=crcr]{
0.924778761061947 0.0869565217391304\\
0.946902654867257 0.217391304347826\\
};
\addplot [
color=ohmred,
solid,
forget plot
]
table[row sep=crcr]{
0.95575221238938 0.173913043478261\\
0.946902654867257 0.217391304347826\\
};
\addplot [
color=ohmred,
solid,
forget plot
]
table[row sep=crcr]{
0.982300884955752 0.0434782608695652\\
0.986725663716814 0.0869565217391304\\
};
\addplot [
color=ohmred,
solid,
forget plot
]
table[row sep=crcr]{
0.991150442477876 0.0434782608695652\\
0.986725663716814 0.0869565217391304\\
};
\addplot [
color=ohmred,
solid,
forget plot
]
table[row sep=crcr]{
0.0176991150442478 0.0434782608695652\\
0.0221238938053097 0.0869565217391304\\
};
\addplot [
color=ohmred,
solid,
forget plot
]
table[row sep=crcr]{
0.0265486725663717 0.0434782608695652\\
0.0221238938053097 0.0869565217391304\\
};
\addplot [
color=ohmred,
solid,
forget plot
]
table[row sep=crcr]{
0.261061946902655 0.739130434782609\\
0.265486725663717 0.782608695652174\\
};
\addplot [
color=ohmred,
solid,
forget plot
]
table[row sep=crcr]{
0.47787610619469 0.0434782608695652\\
0.265486725663717 0.782608695652174\\
};
\addplot [
color=ohmred,
solid,
forget plot
]
table[row sep=crcr]{
0.486725663716814 0.0434782608695652\\
0.694690265486726 0.739130434782609\\
};
\addplot [
color=ohmred,
solid,
forget plot
]
table[row sep=crcr]{
0.699115044247788 0.695652173913043\\
0.694690265486726 0.739130434782609\\
};
\addplot [
color=ohmred,
solid,
forget plot
]
table[row sep=crcr]{
0.920353982300885 0.0434782608695652\\
0.924778761061947 0.0869565217391304\\
};
\addplot [
color=ohmred,
solid,
forget plot
]
table[row sep=crcr]{
0.929203539823009 0.0434782608695652\\
0.924778761061947 0.0869565217391304\\
};
\addplot [
color=ohmred,
solid,
forget plot
]
table[row sep=crcr]{
0.942477876106195 0.0869565217391304\\
0.95575221238938 0.173913043478261\\
};
\addplot [
color=ohmred,
solid,
forget plot
]
table[row sep=crcr]{
0.964601769911504 0.130434782608696\\
0.95575221238938 0.173913043478261\\
};
\addplot [
color=ohmred,
solid,
forget plot
]
table[row sep=crcr]{
0.185840707964602 0.695652173913043\\
0.261061946902655 0.739130434782609\\
};
\addplot [
color=ohmred,
solid,
forget plot
]
table[row sep=crcr]{
0.398230088495575 0.347826086956522\\
0.261061946902655 0.739130434782609\\
};
\addplot [
color=ohmred,
solid,
forget plot
]
table[row sep=crcr]{
0.561946902654867 0.434782608695652\\
0.699115044247788 0.695652173913043\\
};
\addplot [
color=ohmred,
solid,
forget plot
]
table[row sep=crcr]{
0.769911504424779 0.652173913043478\\
0.699115044247788 0.695652173913043\\
};
\addplot [
color=ohmred,
solid,
forget plot
]
table[row sep=crcr]{
0.938053097345133 0.0434782608695652\\
0.942477876106195 0.0869565217391304\\
};
\addplot [
color=ohmred,
solid,
forget plot
]
table[row sep=crcr]{
0.946902654867257 0.0434782608695652\\
0.942477876106195 0.0869565217391304\\
};
\addplot [
color=ohmred,
solid,
forget plot
]
table[row sep=crcr]{
0.960176991150442 0.0869565217391304\\
0.964601769911504 0.130434782608696\\
};
\addplot [
color=ohmred,
solid,
forget plot
]
table[row sep=crcr]{
0.973451327433628 0.0434782608695652\\
0.964601769911504 0.130434782608696\\
};
\addplot [
color=ohmred,
solid,
forget plot
]
table[row sep=crcr]{
0.115044247787611 0.652173913043478\\
0.185840707964602 0.695652173913043\\
};
\addplot [
color=ohmred,
solid,
forget plot
]
table[row sep=crcr]{
0.252212389380531 0.391304347826087\\
0.185840707964602 0.695652173913043\\
};
\addplot [
color=ohmred,
solid,
forget plot
]
table[row sep=crcr]{
0.358407079646018 0.304347826086957\\
0.398230088495575 0.347826086956522\\
};
\addplot [
color=ohmred,
solid,
forget plot
]
table[row sep=crcr]{
0.433628318584071 0.304347826086957\\
0.398230088495575 0.347826086956522\\
};
\addplot [
color=ohmred,
solid,
forget plot
]
table[row sep=crcr]{
0.548672566371681 0.391304347826087\\
0.561946902654867 0.434782608695652\\
};
\addplot [
color=ohmred,
solid,
forget plot
]
table[row sep=crcr]{
0.619469026548673 0.130434782608696\\
0.561946902654867 0.434782608695652\\
};
\addplot [
color=ohmred,
solid,
forget plot
]
table[row sep=crcr]{
0.63716814159292 0.0434782608695652\\
0.769911504424779 0.652173913043478\\
};
\addplot [
color=ohmred,
solid,
forget plot
]
table[row sep=crcr]{
0.774336283185841 0.608695652173913\\
0.769911504424779 0.652173913043478\\
};
\addplot [
color=ohmred,
solid,
forget plot
]
table[row sep=crcr]{
0.95575221238938 0.0434782608695652\\
0.960176991150442 0.0869565217391304\\
};
\addplot [
color=ohmred,
solid,
forget plot
]
table[row sep=crcr]{
0.964601769911504 0.0434782608695652\\
0.960176991150442 0.0869565217391304\\
};
\addplot [
color=ohmred,
solid,
forget plot
]
table[row sep=crcr]{
0.0530973451327434 0.0434782608695652\\
0.115044247787611 0.652173913043478\\
};
\addplot [
color=ohmred,
solid,
forget plot
]
table[row sep=crcr]{
0.119469026548673 0.608695652173913\\
0.115044247787611 0.652173913043478\\
};
\addplot [
color=ohmred,
solid,
forget plot
]
table[row sep=crcr]{
0.20353982300885 0.217391304347826\\
0.252212389380531 0.391304347826087\\
};
\addplot [
color=ohmred,
solid,
forget plot
]
table[row sep=crcr]{
0.274336283185841 0.347826086956522\\
0.252212389380531 0.391304347826087\\
};
\addplot [
color=ohmred,
solid,
forget plot
]
table[row sep=crcr]{
0.327433628318584 0.0434782608695652\\
0.358407079646018 0.304347826086957\\
};
\addplot [
color=ohmred,
solid,
forget plot
]
table[row sep=crcr]{
0.36283185840708 0.260869565217391\\
0.358407079646018 0.304347826086957\\
};
\addplot [
color=ohmred,
solid,
forget plot
]
table[row sep=crcr]{
0.429203539823009 0.260869565217391\\
0.433628318584071 0.304347826086957\\
};
\addplot [
color=ohmred,
solid,
forget plot
]
table[row sep=crcr]{
0.469026548672566 0.0434782608695652\\
0.433628318584071 0.304347826086957\\
};
\addplot [
color=ohmred,
solid,
forget plot
]
table[row sep=crcr]{
0.513274336283186 0.217391304347826\\
0.548672566371681 0.391304347826087\\
};
\addplot [
color=ohmred,
solid,
forget plot
]
table[row sep=crcr]{
0.570796460176991 0.347826086956522\\
0.548672566371681 0.391304347826087\\
};
\addplot [
color=ohmred,
solid,
forget plot
]
table[row sep=crcr]{
0.610619469026549 0.0434782608695652\\
0.619469026548673 0.130434782608696\\
};
\addplot [
color=ohmred,
solid,
forget plot
]
table[row sep=crcr]{
0.623893805309734 0.0869565217391304\\
0.619469026548673 0.130434782608696\\
};
\addplot [
color=ohmred,
solid,
forget plot
]
table[row sep=crcr]{
0.668141592920354 0.217391304347826\\
0.774336283185841 0.608695652173913\\
};
\addplot [
color=ohmred,
solid,
forget plot
]
table[row sep=crcr]{
0.800884955752212 0.565217391304348\\
0.774336283185841 0.608695652173913\\
};
\addplot [
color=ohmred,
solid,
forget plot
]
table[row sep=crcr]{
0.0619469026548673 0.0434782608695652\\
0.119469026548673 0.608695652173913\\
};
\addplot [
color=ohmred,
solid,
forget plot
]
table[row sep=crcr]{
0.123893805309735 0.565217391304348\\
0.119469026548673 0.608695652173913\\
};
\addplot [
color=ohmred,
solid,
forget plot
]
table[row sep=crcr]{
0.185840707964602 0.0434782608695652\\
0.20353982300885 0.217391304347826\\
};
\addplot [
color=ohmred,
solid,
forget plot
]
table[row sep=crcr]{
0.207964601769911 0.173913043478261\\
0.20353982300885 0.217391304347826\\
};
\addplot [
color=ohmred,
solid,
forget plot
]
table[row sep=crcr]{
0.243362831858407 0.173913043478261\\
0.274336283185841 0.347826086956522\\
};
\addplot [
color=ohmred,
solid,
forget plot
]
table[row sep=crcr]{
0.292035398230089 0.304347826086957\\
0.274336283185841 0.347826086956522\\
};
\addplot [
color=ohmred,
solid,
forget plot
]
table[row sep=crcr]{
0.358407079646018 0.217391304347826\\
0.36283185840708 0.260869565217391\\
};
\addplot [
color=ohmred,
solid,
forget plot
]
table[row sep=crcr]{
0.389380530973451 0.0434782608695652\\
0.36283185840708 0.260869565217391\\
};
\addplot [
color=ohmred,
solid,
forget plot
]
table[row sep=crcr]{
0.415929203539823 0.217391304347826\\
0.429203539823009 0.260869565217391\\
};
\addplot [
color=ohmred,
solid,
forget plot
]
table[row sep=crcr]{
0.451327433628319 0.130434782608696\\
0.429203539823009 0.260869565217391\\
};
\addplot [
color=ohmred,
solid,
forget plot
]
table[row sep=crcr]{
0.508849557522124 0.173913043478261\\
0.513274336283186 0.217391304347826\\
};
\addplot [
color=ohmred,
solid,
forget plot
]
table[row sep=crcr]{
0.530973451327434 0.0434782608695652\\
0.513274336283186 0.217391304347826\\
};
\addplot [
color=ohmred,
solid,
forget plot
]
table[row sep=crcr]{
0.566371681415929 0.304347826086957\\
0.570796460176991 0.347826086956522\\
};
\addplot [
color=ohmred,
solid,
forget plot
]
table[row sep=crcr]{
0.601769911504425 0.0434782608695652\\
0.570796460176991 0.347826086956522\\
};
\addplot [
color=ohmred,
solid,
forget plot
]
table[row sep=crcr]{
0.619469026548673 0.0434782608695652\\
0.623893805309734 0.0869565217391304\\
};
\addplot [
color=ohmred,
solid,
forget plot
]
table[row sep=crcr]{
0.628318584070796 0.0434782608695652\\
0.623893805309734 0.0869565217391304\\
};
\addplot [
color=ohmred,
solid,
forget plot
]
table[row sep=crcr]{
0.65929203539823 0.173913043478261\\
0.668141592920354 0.217391304347826\\
};
\addplot [
color=ohmred,
solid,
forget plot
]
table[row sep=crcr]{
0.685840707964602 0.0869565217391304\\
0.668141592920354 0.217391304347826\\
};
\addplot [
color=ohmred,
solid,
forget plot
]
table[row sep=crcr]{
0.787610619469027 0.521739130434783\\
0.800884955752212 0.565217391304348\\
};
\addplot [
color=ohmred,
solid,
forget plot
]
table[row sep=crcr]{
0.893805309734513 0.130434782608696\\
0.800884955752212 0.565217391304348\\
};
\addplot [
color=ohmred,
solid,
forget plot
]
table[row sep=crcr]{
0.119469026548673 0.521739130434783\\
0.123893805309735 0.565217391304348\\
};
\addplot [
color=ohmred,
solid,
forget plot
]
table[row sep=crcr]{
0.176991150442478 0.0434782608695652\\
0.123893805309735 0.565217391304348\\
};
\addplot [
color=ohmred,
solid,
forget plot
]
table[row sep=crcr]{
0.20353982300885 0.130434782608696\\
0.207964601769911 0.173913043478261\\
};
\addplot [
color=ohmred,
solid,
forget plot
]
table[row sep=crcr]{
0.221238938053097 0.0434782608695652\\
0.207964601769911 0.173913043478261\\
};
\addplot [
color=ohmred,
solid,
forget plot
]
table[row sep=crcr]{
0.230088495575221 0.0434782608695652\\
0.243362831858407 0.173913043478261\\
};
\addplot [
color=ohmred,
solid,
forget plot
]
table[row sep=crcr]{
0.247787610619469 0.130434782608696\\
0.243362831858407 0.173913043478261\\
};
\addplot [
color=ohmred,
solid,
forget plot
]
table[row sep=crcr]{
0.265486725663717 0.0434782608695652\\
0.292035398230089 0.304347826086957\\
};
\addplot [
color=ohmred,
solid,
forget plot
]
table[row sep=crcr]{
0.29646017699115 0.260869565217391\\
0.292035398230089 0.304347826086957\\
};
\addplot [
color=ohmred,
solid,
forget plot
]
table[row sep=crcr]{
0.336283185840708 0.0434782608695652\\
0.358407079646018 0.217391304347826\\
};
\addplot [
color=ohmred,
solid,
forget plot
]
table[row sep=crcr]{
0.36283185840708 0.173913043478261\\
0.358407079646018 0.217391304347826\\
};
\addplot [
color=ohmred,
solid,
forget plot
]
table[row sep=crcr]{
0.411504424778761 0.173913043478261\\
0.415929203539823 0.217391304347826\\
};
\addplot [
color=ohmred,
solid,
forget plot
]
table[row sep=crcr]{
0.433628318584071 0.0434782608695652\\
0.415929203539823 0.217391304347826\\
};
\addplot [
color=ohmred,
solid,
forget plot
]
table[row sep=crcr]{
0.446902654867257 0.0869565217391304\\
0.451327433628319 0.130434782608696\\
};
\addplot [
color=ohmred,
solid,
forget plot
]
table[row sep=crcr]{
0.460176991150442 0.0434782608695652\\
0.451327433628319 0.130434782608696\\
};
\addplot [
color=ohmred,
solid,
forget plot
]
table[row sep=crcr]{
0.495575221238938 0.0434782608695652\\
0.508849557522124 0.173913043478261\\
};
\addplot [
color=ohmred,
solid,
forget plot
]
table[row sep=crcr]{
0.513274336283186 0.130434782608696\\
0.508849557522124 0.173913043478261\\
};
\addplot [
color=ohmred,
solid,
forget plot
]
table[row sep=crcr]{
0.561946902654867 0.260869565217391\\
0.566371681415929 0.304347826086957\\
};
\addplot [
color=ohmred,
solid,
forget plot
]
table[row sep=crcr]{
0.592920353982301 0.0434782608695652\\
0.566371681415929 0.304347826086957\\
};
\addplot [
color=ohmred,
solid,
forget plot
]
table[row sep=crcr]{
0.654867256637168 0.130434782608696\\
0.65929203539823 0.173913043478261\\
};
\addplot [
color=ohmred,
solid,
forget plot
]
table[row sep=crcr]{
0.672566371681416 0.0434782608695652\\
0.65929203539823 0.173913043478261\\
};
\addplot [
color=ohmred,
solid,
forget plot
]
table[row sep=crcr]{
0.68141592920354 0.0434782608695652\\
0.685840707964602 0.0869565217391304\\
};
\addplot [
color=ohmred,
solid,
forget plot
]
table[row sep=crcr]{
0.690265486725664 0.0434782608695652\\
0.685840707964602 0.0869565217391304\\
};
\addplot [
color=ohmred,
solid,
forget plot
]
table[row sep=crcr]{
0.783185840707965 0.478260869565217\\
0.787610619469027 0.521739130434783\\
};
\addplot [
color=ohmred,
solid,
forget plot
]
table[row sep=crcr]{
0.876106194690266 0.0434782608695652\\
0.787610619469027 0.521739130434783\\
};
\addplot [
color=ohmred,
solid,
forget plot
]
table[row sep=crcr]{
0.884955752212389 0.0434782608695652\\
0.893805309734513 0.130434782608696\\
};
\addplot [
color=ohmred,
solid,
forget plot
]
table[row sep=crcr]{
0.898230088495575 0.0869565217391304\\
0.893805309734513 0.130434782608696\\
};
\addplot [
color=ohmred,
solid,
forget plot
]
table[row sep=crcr]{
0.115044247787611 0.478260869565217\\
0.119469026548673 0.521739130434783\\
};
\addplot [
color=ohmred,
solid,
forget plot
]
table[row sep=crcr]{
0.168141592920354 0.0434782608695652\\
0.119469026548673 0.521739130434783\\
};
\addplot [
color=ohmred,
solid,
forget plot
]
table[row sep=crcr]{
0.194690265486726 0.0434782608695652\\
0.20353982300885 0.130434782608696\\
};
\addplot [
color=ohmred,
solid,
forget plot
]
table[row sep=crcr]{
0.207964601769911 0.0869565217391304\\
0.20353982300885 0.130434782608696\\
};
\addplot [
color=ohmred,
solid,
forget plot
]
table[row sep=crcr]{
0.243362831858407 0.0869565217391304\\
0.247787610619469 0.130434782608696\\
};
\addplot [
color=ohmred,
solid,
forget plot
]
table[row sep=crcr]{
0.256637168141593 0.0434782608695652\\
0.247787610619469 0.130434782608696\\
};
\addplot [
color=ohmred,
solid,
forget plot
]
table[row sep=crcr]{
0.292035398230089 0.217391304347826\\
0.29646017699115 0.260869565217391\\
};
\addplot [
color=ohmred,
solid,
forget plot
]
table[row sep=crcr]{
0.31858407079646 0.0434782608695652\\
0.29646017699115 0.260869565217391\\
};
\addplot [
color=ohmred,
solid,
forget plot
]
table[row sep=crcr]{
0.353982300884956 0.130434782608696\\
0.36283185840708 0.173913043478261\\
};
\addplot [
color=ohmred,
solid,
forget plot
]
table[row sep=crcr]{
0.376106194690265 0.0869565217391304\\
0.36283185840708 0.173913043478261\\
};
\addplot [
color=ohmred,
solid,
forget plot
]
table[row sep=crcr]{
0.407079646017699 0.130434782608696\\
0.411504424778761 0.173913043478261\\
};
\addplot [
color=ohmred,
solid,
forget plot
]
table[row sep=crcr]{
0.424778761061947 0.0434782608695652\\
0.411504424778761 0.173913043478261\\
};
\addplot [
color=ohmred,
solid,
forget plot
]
table[row sep=crcr]{
0.442477876106195 0.0434782608695652\\
0.446902654867257 0.0869565217391304\\
};
\addplot [
color=ohmred,
solid,
forget plot
]
table[row sep=crcr]{
0.451327433628319 0.0434782608695652\\
0.446902654867257 0.0869565217391304\\
};
\addplot [
color=ohmred,
solid,
forget plot
]
table[row sep=crcr]{
0.508849557522124 0.0869565217391304\\
0.513274336283186 0.130434782608696\\
};
\addplot [
color=ohmred,
solid,
forget plot
]
table[row sep=crcr]{
0.52212389380531 0.0434782608695652\\
0.513274336283186 0.130434782608696\\
};
\addplot [
color=ohmred,
solid,
forget plot
]
table[row sep=crcr]{
0.539823008849557 0.0434782608695652\\
0.561946902654867 0.260869565217391\\
};
\addplot [
color=ohmred,
solid,
forget plot
]
table[row sep=crcr]{
0.566371681415929 0.217391304347826\\
0.561946902654867 0.260869565217391\\
};
\addplot [
color=ohmred,
solid,
forget plot
]
table[row sep=crcr]{
0.650442477876106 0.0869565217391304\\
0.654867256637168 0.130434782608696\\
};
\addplot [
color=ohmred,
solid,
forget plot
]
table[row sep=crcr]{
0.663716814159292 0.0434782608695652\\
0.654867256637168 0.130434782608696\\
};
\addplot [
color=ohmred,
solid,
forget plot
]
table[row sep=crcr]{
0.752212389380531 0.434782608695652\\
0.783185840707965 0.478260869565217\\
};
\addplot [
color=ohmred,
solid,
forget plot
]
table[row sep=crcr]{
0.84070796460177 0.304347826086957\\
0.783185840707965 0.478260869565217\\
};
\addplot [
color=ohmred,
solid,
forget plot
]
table[row sep=crcr]{
0.893805309734513 0.0434782608695652\\
0.898230088495575 0.0869565217391304\\
};
\addplot [
color=ohmred,
solid,
forget plot
]
table[row sep=crcr]{
0.902654867256637 0.0434782608695652\\
0.898230088495575 0.0869565217391304\\
};
\addplot [
color=ohmred,
solid,
forget plot
]
table[row sep=crcr]{
0.110619469026549 0.434782608695652\\
0.115044247787611 0.478260869565217\\
};
\addplot [
color=ohmred,
solid,
forget plot
]
table[row sep=crcr]{
0.15929203539823 0.0434782608695652\\
0.115044247787611 0.478260869565217\\
};
\addplot [
color=ohmred,
solid,
forget plot
]
table[row sep=crcr]{
0.20353982300885 0.0434782608695652\\
0.207964601769911 0.0869565217391304\\
};
\addplot [
color=ohmred,
solid,
forget plot
]
table[row sep=crcr]{
0.212389380530973 0.0434782608695652\\
0.207964601769911 0.0869565217391304\\
};
\addplot [
color=ohmred,
solid,
forget plot
]
table[row sep=crcr]{
0.238938053097345 0.0434782608695652\\
0.243362831858407 0.0869565217391304\\
};
\addplot [
color=ohmred,
solid,
forget plot
]
table[row sep=crcr]{
0.247787610619469 0.0434782608695652\\
0.243362831858407 0.0869565217391304\\
};
\addplot [
color=ohmred,
solid,
forget plot
]
table[row sep=crcr]{
0.287610619469027 0.173913043478261\\
0.292035398230089 0.217391304347826\\
};
\addplot [
color=ohmred,
solid,
forget plot
]
table[row sep=crcr]{
0.309734513274336 0.0434782608695652\\
0.292035398230089 0.217391304347826\\
};
\addplot [
color=ohmred,
solid,
forget plot
]
table[row sep=crcr]{
0.349557522123894 0.0869565217391304\\
0.353982300884956 0.130434782608696\\
};
\addplot [
color=ohmred,
solid,
forget plot
]
table[row sep=crcr]{
0.36283185840708 0.0434782608695652\\
0.353982300884956 0.130434782608696\\
};
\addplot [
color=ohmred,
solid,
forget plot
]
table[row sep=crcr]{
0.371681415929204 0.0434782608695652\\
0.376106194690265 0.0869565217391304\\
};
\addplot [
color=ohmred,
solid,
forget plot
]
table[row sep=crcr]{
0.380530973451327 0.0434782608695652\\
0.376106194690265 0.0869565217391304\\
};
\addplot [
color=ohmred,
solid,
forget plot
]
table[row sep=crcr]{
0.402654867256637 0.0869565217391304\\
0.407079646017699 0.130434782608696\\
};
\addplot [
color=ohmred,
solid,
forget plot
]
table[row sep=crcr]{
0.415929203539823 0.0434782608695652\\
0.407079646017699 0.130434782608696\\
};
\addplot [
color=ohmred,
solid,
forget plot
]
table[row sep=crcr]{
0.504424778761062 0.0434782608695652\\
0.508849557522124 0.0869565217391304\\
};
\addplot [
color=ohmred,
solid,
forget plot
]
table[row sep=crcr]{
0.513274336283186 0.0434782608695652\\
0.508849557522124 0.0869565217391304\\
};
\addplot [
color=ohmred,
solid,
forget plot
]
table[row sep=crcr]{
0.548672566371681 0.0434782608695652\\
0.566371681415929 0.217391304347826\\
};
\addplot [
color=ohmred,
solid,
forget plot
]
table[row sep=crcr]{
0.570796460176991 0.173913043478261\\
0.566371681415929 0.217391304347826\\
};
\addplot [
color=ohmred,
solid,
forget plot
]
table[row sep=crcr]{
0.646017699115044 0.0434782608695652\\
0.650442477876106 0.0869565217391304\\
};
\addplot [
color=ohmred,
solid,
forget plot
]
table[row sep=crcr]{
0.654867256637168 0.0434782608695652\\
0.650442477876106 0.0869565217391304\\
};
\addplot [
color=ohmred,
solid,
forget plot
]
table[row sep=crcr]{
0.738938053097345 0.391304347826087\\
0.752212389380531 0.434782608695652\\
};
\addplot [
color=ohmred,
solid,
forget plot
]
table[row sep=crcr]{
0.79646017699115 0.130434782608696\\
0.752212389380531 0.434782608695652\\
};
\addplot [
color=ohmred,
solid,
forget plot
]
table[row sep=crcr]{
0.814159292035398 0.0434782608695652\\
0.84070796460177 0.304347826086957\\
};
\addplot [
color=ohmred,
solid,
forget plot
]
table[row sep=crcr]{
0.845132743362832 0.260869565217391\\
0.84070796460177 0.304347826086957\\
};
\addplot [
color=ohmred,
solid,
forget plot
]
table[row sep=crcr]{
0.106194690265487 0.391304347826087\\
0.110619469026549 0.434782608695652\\
};
\addplot [
color=ohmred,
solid,
forget plot
]
table[row sep=crcr]{
0.150442477876106 0.0434782608695652\\
0.110619469026549 0.434782608695652\\
};
\addplot [
color=ohmred,
solid,
forget plot
]
table[row sep=crcr]{
0.283185840707965 0.130434782608696\\
0.287610619469027 0.173913043478261\\
};
\addplot [
color=ohmred,
solid,
forget plot
]
table[row sep=crcr]{
0.300884955752212 0.0434782608695652\\
0.287610619469027 0.173913043478261\\
};
\addplot [
color=ohmred,
solid,
forget plot
]
table[row sep=crcr]{
0.345132743362832 0.0434782608695652\\
0.349557522123894 0.0869565217391304\\
};
\addplot [
color=ohmred,
solid,
forget plot
]
table[row sep=crcr]{
0.353982300884956 0.0434782608695652\\
0.349557522123894 0.0869565217391304\\
};
\addplot [
color=ohmred,
solid,
forget plot
]
table[row sep=crcr]{
0.398230088495575 0.0434782608695652\\
0.402654867256637 0.0869565217391304\\
};
\addplot [
color=ohmred,
solid,
forget plot
]
table[row sep=crcr]{
0.407079646017699 0.0434782608695652\\
0.402654867256637 0.0869565217391304\\
};
\addplot [
color=ohmred,
solid,
forget plot
]
table[row sep=crcr]{
0.557522123893805 0.0434782608695652\\
0.570796460176991 0.173913043478261\\
};
\addplot [
color=ohmred,
solid,
forget plot
]
table[row sep=crcr]{
0.575221238938053 0.130434782608696\\
0.570796460176991 0.173913043478261\\
};
\addplot [
color=ohmred,
solid,
forget plot
]
table[row sep=crcr]{
0.734513274336283 0.347826086956522\\
0.738938053097345 0.391304347826087\\
};
\addplot [
color=ohmred,
solid,
forget plot
]
table[row sep=crcr]{
0.778761061946903 0.0434782608695652\\
0.738938053097345 0.391304347826087\\
};
\addplot [
color=ohmred,
solid,
forget plot
]
table[row sep=crcr]{
0.787610619469027 0.0434782608695652\\
0.79646017699115 0.130434782608696\\
};
\addplot [
color=ohmred,
solid,
forget plot
]
table[row sep=crcr]{
0.800884955752212 0.0869565217391304\\
0.79646017699115 0.130434782608696\\
};
\addplot [
color=ohmred,
solid,
forget plot
]
table[row sep=crcr]{
0.823008849557522 0.0434782608695652\\
0.845132743362832 0.260869565217391\\
};
\addplot [
color=ohmred,
solid,
forget plot
]
table[row sep=crcr]{
0.849557522123894 0.217391304347826\\
0.845132743362832 0.260869565217391\\
};
\addplot [
color=ohmred,
solid,
forget plot
]
table[row sep=crcr]{
0.101769911504425 0.347826086956522\\
0.106194690265487 0.391304347826087\\
};
\addplot [
color=ohmred,
solid,
forget plot
]
table[row sep=crcr]{
0.141592920353982 0.0434782608695652\\
0.106194690265487 0.391304347826087\\
};
\addplot [
color=ohmred,
solid,
forget plot
]
table[row sep=crcr]{
0.278761061946903 0.0869565217391304\\
0.283185840707965 0.130434782608696\\
};
\addplot [
color=ohmred,
solid,
forget plot
]
table[row sep=crcr]{
0.292035398230089 0.0434782608695652\\
0.283185840707965 0.130434782608696\\
};
\addplot [
color=ohmred,
solid,
forget plot
]
table[row sep=crcr]{
0.566371681415929 0.0434782608695652\\
0.575221238938053 0.130434782608696\\
};
\addplot [
color=ohmred,
solid,
forget plot
]
table[row sep=crcr]{
0.579646017699115 0.0869565217391304\\
0.575221238938053 0.130434782608696\\
};
\addplot [
color=ohmred,
solid,
forget plot
]
table[row sep=crcr]{
0.730088495575221 0.304347826086957\\
0.734513274336283 0.347826086956522\\
};
\addplot [
color=ohmred,
solid,
forget plot
]
table[row sep=crcr]{
0.769911504424779 0.0434782608695652\\
0.734513274336283 0.347826086956522\\
};
\addplot [
color=ohmred,
solid,
forget plot
]
table[row sep=crcr]{
0.79646017699115 0.0434782608695652\\
0.800884955752212 0.0869565217391304\\
};
\addplot [
color=ohmred,
solid,
forget plot
]
table[row sep=crcr]{
0.805309734513274 0.0434782608695652\\
0.800884955752212 0.0869565217391304\\
};
\addplot [
color=ohmred,
solid,
forget plot
]
table[row sep=crcr]{
0.831858407079646 0.0434782608695652\\
0.849557522123894 0.217391304347826\\
};
\addplot [
color=ohmred,
solid,
forget plot
]
table[row sep=crcr]{
0.853982300884956 0.173913043478261\\
0.849557522123894 0.217391304347826\\
};
\addplot [
color=ohmred,
solid,
forget plot
]
table[row sep=crcr]{
0.0973451327433628 0.304347826086957\\
0.101769911504425 0.347826086956522\\
};
\addplot [
color=ohmred,
solid,
forget plot
]
table[row sep=crcr]{
0.132743362831858 0.0434782608695652\\
0.101769911504425 0.347826086956522\\
};
\addplot [
color=ohmred,
solid,
forget plot
]
table[row sep=crcr]{
0.274336283185841 0.0434782608695652\\
0.278761061946903 0.0869565217391304\\
};
\addplot [
color=ohmred,
solid,
forget plot
]
table[row sep=crcr]{
0.283185840707965 0.0434782608695652\\
0.278761061946903 0.0869565217391304\\
};
\addplot [
color=ohmred,
solid,
forget plot
]
table[row sep=crcr]{
0.575221238938053 0.0434782608695652\\
0.579646017699115 0.0869565217391304\\
};
\addplot [
color=ohmred,
solid,
forget plot
]
table[row sep=crcr]{
0.584070796460177 0.0434782608695652\\
0.579646017699115 0.0869565217391304\\
};
\addplot [
color=ohmred,
solid,
forget plot
]
table[row sep=crcr]{
0.725663716814159 0.260869565217391\\
0.730088495575221 0.304347826086957\\
};
\addplot [
color=ohmred,
solid,
forget plot
]
table[row sep=crcr]{
0.761061946902655 0.0434782608695652\\
0.730088495575221 0.304347826086957\\
};
\addplot [
color=ohmred,
solid,
forget plot
]
table[row sep=crcr]{
0.849557522123894 0.130434782608696\\
0.853982300884956 0.173913043478261\\
};
\addplot [
color=ohmred,
solid,
forget plot
]
table[row sep=crcr]{
0.867256637168142 0.0434782608695652\\
0.853982300884956 0.173913043478261\\
};
\addplot [
color=ohmred,
solid,
forget plot
]
table[row sep=crcr]{
0.0929203539823009 0.260869565217391\\
0.0973451327433628 0.304347826086957\\
};
\addplot [
color=ohmred,
solid,
forget plot
]
table[row sep=crcr]{
0.123893805309735 0.0434782608695652\\
0.0973451327433628 0.304347826086957\\
};
\addplot [
color=ohmred,
solid,
forget plot
]
table[row sep=crcr]{
0.721238938053097 0.217391304347826\\
0.725663716814159 0.260869565217391\\
};
\addplot [
color=ohmred,
solid,
forget plot
]
table[row sep=crcr]{
0.752212389380531 0.0434782608695652\\
0.725663716814159 0.260869565217391\\
};
\addplot [
color=ohmred,
solid,
forget plot
]
table[row sep=crcr]{
0.845132743362832 0.0869565217391304\\
0.849557522123894 0.130434782608696\\
};
\addplot [
color=ohmred,
solid,
forget plot
]
table[row sep=crcr]{
0.858407079646018 0.0434782608695652\\
0.849557522123894 0.130434782608696\\
};
\addplot [
color=ohmred,
solid,
forget plot
]
table[row sep=crcr]{
0.0884955752212389 0.217391304347826\\
0.0929203539823009 0.260869565217391\\
};
\addplot [
color=ohmred,
solid,
forget plot
]
table[row sep=crcr]{
0.115044247787611 0.0434782608695652\\
0.0929203539823009 0.260869565217391\\
};
\addplot [
color=ohmred,
solid,
forget plot
]
table[row sep=crcr]{
0.712389380530973 0.173913043478261\\
0.721238938053097 0.217391304347826\\
};
\addplot [
color=ohmred,
solid,
forget plot
]
table[row sep=crcr]{
0.738938053097345 0.0869565217391304\\
0.721238938053097 0.217391304347826\\
};
\addplot [
color=ohmred,
solid,
forget plot
]
table[row sep=crcr]{
0.84070796460177 0.0434782608695652\\
0.845132743362832 0.0869565217391304\\
};
\addplot [
color=ohmred,
solid,
forget plot
]
table[row sep=crcr]{
0.849557522123894 0.0434782608695652\\
0.845132743362832 0.0869565217391304\\
};
\addplot [
color=ohmred,
solid,
forget plot
]
table[row sep=crcr]{
0.084070796460177 0.173913043478261\\
0.0884955752212389 0.217391304347826\\
};
\addplot [
color=ohmred,
solid,
forget plot
]
table[row sep=crcr]{
0.106194690265487 0.0434782608695652\\
0.0884955752212389 0.217391304347826\\
};
\addplot [
color=ohmred,
solid,
forget plot
]
table[row sep=crcr]{
0.707964601769911 0.130434782608696\\
0.712389380530973 0.173913043478261\\
};
\addplot [
color=ohmred,
solid,
forget plot
]
table[row sep=crcr]{
0.725663716814159 0.0434782608695652\\
0.712389380530973 0.173913043478261\\
};
\addplot [
color=ohmred,
solid,
forget plot
]
table[row sep=crcr]{
0.734513274336283 0.0434782608695652\\
0.738938053097345 0.0869565217391304\\
};
\addplot [
color=ohmred,
solid,
forget plot
]
table[row sep=crcr]{
0.743362831858407 0.0434782608695652\\
0.738938053097345 0.0869565217391304\\
};
\addplot [
color=ohmred,
solid,
forget plot
]
table[row sep=crcr]{
0.079646017699115 0.130434782608696\\
0.084070796460177 0.173913043478261\\
};
\addplot [
color=ohmred,
solid,
forget plot
]
table[row sep=crcr]{
0.0973451327433628 0.0434782608695652\\
0.084070796460177 0.173913043478261\\
};
\addplot [
color=ohmred,
solid,
forget plot
]
table[row sep=crcr]{
0.70353982300885 0.0869565217391304\\
0.707964601769911 0.130434782608696\\
};
\addplot [
color=ohmred,
solid,
forget plot
]
table[row sep=crcr]{
0.716814159292035 0.0434782608695652\\
0.707964601769911 0.130434782608696\\
};
\addplot [
color=ohmred,
solid,
forget plot
]
table[row sep=crcr]{
0.0752212389380531 0.0869565217391304\\
0.079646017699115 0.130434782608696\\
};
\addplot [
color=ohmred,
solid,
forget plot
]
table[row sep=crcr]{
0.0884955752212389 0.0434782608695652\\
0.079646017699115 0.130434782608696\\
};
\addplot [
color=ohmred,
solid,
forget plot
]
table[row sep=crcr]{
0.699115044247788 0.0434782608695652\\
0.70353982300885 0.0869565217391304\\
};
\addplot [
color=ohmred,
solid,
forget plot
]
table[row sep=crcr]{
0.707964601769911 0.0434782608695652\\
0.70353982300885 0.0869565217391304\\
};
\addplot [
color=ohmred,
solid,
forget plot
]
table[row sep=crcr]{
0.0707964601769911 0.0434782608695652\\
0.0752212389380531 0.0869565217391304\\
};
\addplot [
color=ohmred,
solid,
forget plot
]
table[row sep=crcr]{
0.079646017699115 0.0434782608695652\\
0.0752212389380531 0.0869565217391304\\
};
\end{axis}
\end{tikzpicture}% 
    
    Tiefe: 21
\end{frame}


\subsection{Under- \& Overfitting}

\begin{frame}{Under- \& Overfitting (1)}

\textbf{Was ist Overfitting? -- Eine anschauliche Erklärung}

Klassifikation einer Eiche anhand ihrer Silhouette

\pause

\begin{center}
    \includegraphics[height=0.9\textheight]{src/baumsilhouette}
\end{center}


\end{frame}



\begin{frame}{Blick zur Tafel}
\centering
\includegraphics[width=\linewidth]{src/blackboard}
\end{frame}



\begin{frame}{Under- \& Overfitting (1)}

\textbf{Was ist Overfitting? -- Eine anschauliche Erklärung}

Klassifikation einer Eiche anhand ihrer Silhouette

\begin{center}
\includegraphics[height=0.9\textheight]{src/baumsilhouette} \pause
\includegraphics[height=0.9\textheight]{src/baumsilhouettewinter}
\end{center}


\end{frame}


\begin{frame}{Under- \& Overfitting (1)}

\textbf{Was ist Overfitting? -- Eine anschauliche Erklärung}

Klassifikation einer Eiche anhand ihrer Silhouette

\begin{center}
\includegraphics[height=0.9\textheight]{src/lollischwarz} \pause
\includegraphics[height=0.9\textheight]{src/lollibunt}
\end{center}


\end{frame}





\begin{frame}{Under- \& Overfitting (2)}

\textbf{Was ist Overfitting? -- Eine mathematische Erklärung}\\[1em]
      
\setlength\figureheight{0.3\linewidth} 
\setlength\figurewidth{0.45\linewidth} 

\begin{columns}[t]
    \begin{column}{0.5\linewidth} \pause
        \centering 
        % This file was created by matlab2tikz v0.4.4 running on MATLAB 7.14.
% Copyright (c) 2008--2013, Nico Schlömer <nico.schloemer@gmail.com>
% All rights reserved.
% 
% The latest updates can be retrieved from
%   http://www.mathworks.com/matlabcentral/fileexchange/22022-matlab2tikz
% where you can also make suggestions and rate matlab2tikz.
% 
\begin{tikzpicture}

\begin{axis}[%
width=\figurewidth,
height=\figureheight,
scale only axis,
xmin=0,
xmax=2,
xtick={0,0.2,0.4,0.6,0.8,1,1.2,1.4,1.6,1.8,2},
xticklabels={\empty},
%xlabel={x},
ymin=-1,
ymax=2,
ytick={-1,-0.5,0,0.5,1,1.5,2},
yticklabels={\empty},
%ylabel={y}
]
\addplot [
color=ohmgreen,
dashed,
forget plot
]
table[row sep=crcr]{
0 0\\
0.105263157894737 0.283714827234291\\
0.210526315789474 0.50794576468873\\
0.315789473684211 0.679690917043301\\
0.421052631578947 0.805948388977985\\
0.526315789473684 0.893716285172766\\
0.631578947368421 0.949992710307625\\
0.736842105263158 0.981775769062546\\
0.842105263157895 0.99606356611751\\
0.947368421052632 0.9998542061525\\
1.05263157894737 1.0001457938475\\
1.15789473684211 1.00393643388249\\
1.26315789473684 1.01822423093745\\
1.36842105263158 1.05000728969237\\
1.47368421052632 1.10628371482723\\
1.57894736842105 1.19405161102201\\
1.68421052631579 1.3203090829567\\
1.78947368421053 1.49205423531127\\
1.89473684210526 1.71628517276571\\
2 2\\
};
\addplot [
color=blue,
only marks,
mark=+,
mark options={solid},
forget plot
]
table[row sep=crcr]{
0.210526315789474 0.740086643568247\\
0.421052631578947 1.15661215613222\\
0.526315789473684 1.26387315802164\\
0.631578947368421 1.54536820357361\\
0.736842105263158 1.57646482741354\\
1.26315789473684 1.84459521143237\\
1.36842105263158 1.83993016576459\\
1.47368421052632 1.25385192431208\\
1.68421052631579 1.67571298787751\\
};
\addplot [
color=blue,
only marks,
mark=o,
mark options={solid},
forget plot
]
table[row sep=crcr]{
0 -0.94973934981283\\
0.105263157894737 -0.0114064877954784\\
0.315789473684211 0.228558913453212\\
0.842105263157895 0.751642974658432\\
0.947368421052632 0.471435402852224\\
1.05263157894737 0.924340802677373\\
1.15789473684211 0.571509732448698\\
1.57894736842105 1.08048807903004\\
1.78947368421053 1.06391742312104\\
1.89473684210526 1.55377678479138\\
2 1.83476664984637\\
};
\end{axis}
\end{tikzpicture}% 
    \end{column}
    \begin{column}{0.5\linewidth} \pause
        \centering 
        \input{src/underfittingSample.tikz} 
    \end{column}
\end{columns}

\begin{columns}[t]
    \begin{column}{0.5\linewidth} \pause
        \centering 
    \input{src/goodfittingSample.tikz} 
    \end{column}
    \begin{column}{0.5\linewidth} \pause
        \centering 
    \input{src/overfittingSample.tikz} 
    \end{column}
\end{columns}


\end{frame}







\begin{frame}{Under- \& Overfitting (3)}

\textbf{Was ist Overfitting? -- Eine anwendungsorientierte Erklärung}\\[1em]

    \centering 
    \setlength\figureheight{0.45\textwidth} 
    \setlength\figurewidth{0.8\textwidth} 
    % This file was created by matlab2tikz v0.4.4 running on MATLAB 7.14.
% Copyright (c) 2008--2013, Nico Schlömer <nico.schloemer@gmail.com>
% All rights reserved.
% 
% The latest updates can be retrieved from
%   http://www.mathworks.com/matlabcentral/fileexchange/22022-matlab2tikz
% where you can also make suggestions and rate matlab2tikz.
% 
\begin{tikzpicture}[font=\sffamily\sansmath]

\begin{axis}[%
width=\figurewidth,
height=\figureheight,
scale only axis,
xmin=0,
xmax=150,
ymin=87,
ymax=95,
xlabel={Anzahl der Knoten},
ylabel={Genauigkeit in \%},
axis x line*=bottom,
axis y line*=left,
legend style={at={(0.03,0.97)},anchor=north west,draw=black,fill=white,legend cell align=left}
]
\addplot [
color=ohmred,
solid
]
table[row sep=crcr]{
9 88.89272658693\\
13 88.89272658693\\
15 88.89272658693\\
17 88.89272658693\\
19 88.89272658693\\
21 88.89272658693\\
23 89.1850636403039\\
25 89.1850636403039\\
27 89.1850636403039\\
31 89.1850636403039\\
33 89.2474492448748\\
37 90.3740148848439\\
39 90.5330118641541\\
41 90.6350676432999\\
43 91.2105477456059\\
45 91.413153670209\\
49 91.5355027729822\\
51 91.6170688414977\\
55 91.7643797116905\\
59 91.845945780206\\
61 91.9682948829792\\
63 92.6487838591103\\
67 93.0050460674491\\
73 93.1258785058728\\
75 93.1652485846129\\
77 93.3295538891465\\
79 93.4651057653867\\
81 93.5037175119521\\
87 93.6610155233515\\
93 93.7397556808318\\
97 93.8161296900072\\
101 93.8554997687474\\
111 94.0525813972391\\
115 94.1336877030243\\
123 94.2102440157609\\
125 94.2102440157609\\
131 94.2549224197245\\
133 94.2935341662898\\
141 94.3692409950711\\
};
\addlegendentry{Training};

\addplot [
color=ohmgreen,
solid
]
table[row sep=crcr]{
9 88.0593842469951\\
13 88.0593842469951\\
15 88.0593842469951\\
17 88.0593842469951\\
19 88.0593842469951\\
21 88.0593842469951\\
23 87.9363789823697\\
25 87.9363789823697\\
27 87.9363789823697\\
31 87.9363789823697\\
33 87.6823925984445\\
37 88.3392407115438\\
39 88.2404264822947\\
41 87.6351423632067\\
43 88.948854179426\\
45 89.4354109356141\\
49 89.5304679698347\\
51 89.7205820382757\\
55 90.572464430401\\
59 90.572464430401\\
61 90.6637642695931\\
63 91.5064480259885\\
67 92.3018820705657\\
73 92.3931819097577\\
75 92.3931819097577\\
77 92.4821557178367\\
79 92.4477142437416\\
81 92.3489000144926\\
87 92.5465284729906\\
93 92.6453427022396\\
97 92.5465284729906\\
101 92.5465284729906\\
111 91.9742915679692\\
115 92.0645442755504\\
123 92.566093431385\\
125 92.566093431385\\
131 92.4574387809659\\
133 92.1462354614638\\
141 92.0523414427998\\
};
\addlegendentry{Test};

\end{axis}
\end{tikzpicture}% 
\end{frame}



\subsection{Optimierungen}

\begin{frame}{Optimierungen}
\begin{itemize}[<+->]
    \item Pruning mit Schwellwert
    \item Reduced-Error-Pruning
    \item Cross Validation
    
        \pause 
        
        \textcolor{ohmgreen}{Trainingsdaten} und \textcolor{ohmred}{Testdaten}: \\[1em]
        
        \begin{tikzpicture}[rotate=0]
    \draw[fill=ohmred] (0,0) circle [radius=0.5];
    \clip (0,0) circle [radius=0.49];
    \draw[fill= ohmgreen] (0,0) -- (72:1) -- (144:1);
    \draw[fill= ohmgreen] (0,0) -- (144:1) -- (216:1);
    \draw[fill= ohmgreen] (0,0) -- (216:1) -- (288:1);
    \draw[fill= ohmgreen] (0,0) -- (288:1)-- (360:1);
    \draw[] (0,0) -- (360:1);
\end{tikzpicture}
% 
\begin{tikzpicture}[rotate=72]
    \draw[fill=ohmred] (0,0) circle [radius=0.5];
    \clip (0,0) circle [radius=0.49];
    \draw[fill= ohmgreen] (0,0) -- (72:1) -- (144:1);
    \draw[fill= ohmgreen] (0,0) -- (144:1) -- (216:1);
    \draw[fill= ohmgreen] (0,0) -- (216:1) -- (288:1);
    \draw[fill= ohmgreen] (0,0) -- (288:1)-- (360:1);
    \draw[] (0,0) -- (360:1);
\end{tikzpicture}
% 
\begin{tikzpicture}[rotate=144]
    \draw[fill=ohmred] (0,0) circle [radius=0.5];
    \clip (0,0) circle [radius=0.49];
    \draw[fill= ohmgreen] (0,0) -- (72:1) -- (144:1);
    \draw[fill= ohmgreen] (0,0) -- (144:1) -- (216:1);
    \draw[fill= ohmgreen] (0,0) -- (216:1) -- (288:1);
    \draw[fill= ohmgreen] (0,0) -- (288:1)-- (360:1);
    \draw[] (0,0) -- (360:1);
\end{tikzpicture}
% 
\begin{tikzpicture}[rotate=216]
    \draw[fill=ohmred] (0,0) circle [radius=0.5];
    \clip (0,0) circle [radius=0.49];
    \draw[fill= ohmgreen] (0,0) -- (72:1) -- (144:1);
    \draw[fill= ohmgreen] (0,0) -- (144:1) -- (216:1);
    \draw[fill= ohmgreen] (0,0) -- (216:1) -- (288:1);
    \draw[fill= ohmgreen] (0,0) -- (288:1)-- (360:1);
    \draw[] (0,0) -- (360:1);
\end{tikzpicture}
% 
\begin{tikzpicture}[rotate=288]
    \draw[fill=ohmred] (0,0) circle [radius=0.5];
    \clip (0,0) circle [radius=0.49];
    \draw[fill= ohmgreen] (0,0) -- (72:1) -- (144:1);
    \draw[fill= ohmgreen] (0,0) -- (144:1) -- (216:1);
    \draw[fill= ohmgreen] (0,0) -- (216:1) -- (288:1);
    \draw[fill= ohmgreen] (0,0) -- (288:1)-- (360:1);
    \draw[] (0,0) -- (360:1);
\end{tikzpicture}
% 
    \item Bagging (\underline{B}ootstrap \underline{agg}regat\underline{ing})
    \item Induktiver Bias
    \item Anderes Kriterium für Bestenauswahl (z.B. Gini-Index $\sum_i P(C_i)^2$)
\end{itemize}


\end{frame}



\subsection{Ergebnisse \& Vergleiche}
\begin{frame}{Ergebnis}
\setlength\figureheight{0.4\textwidth} 
\setlength\figurewidth{0.45\textwidth} 
\begin{columns}[b]
        \column[t]{0.45\textwidth} \centering \textbf{Kompletter Entscheidungsbaum} \\[1em] % This file was created by matlab2tikz v0.4.4 running on MATLAB 7.14.
% Copyright (c) 2008--2013, Nico Schlömer <nico.schloemer@gmail.com>
% All rights reserved.
% 
% The latest updates can be retrieved from
%   http://www.mathworks.com/matlabcentral/fileexchange/22022-matlab2tikz
% where you can also make suggestions and rate matlab2tikz.
% 
\begin{tikzpicture}

\begin{axis}[%
width=\figurewidth,
height=\figureheight,
scale only axis,
xmin=0,
xmax=1,
ymin=0,
ymax=1,
hide axis
]
\addplot [
color=ohmred,
only marks,
mark=o,
mark options={solid},
forget plot
]
table[row sep=crcr]{
0.5 0.956521739130435\\
0.0265486725663717 0.217391304347826\\
0.52212389380531 0.91304347826087\\
0.0221238938053097 0.173913043478261\\
0.0442477876106195 0.0434782608695652\\
0.482300884955752 0.869565217391304\\
0.95575221238938 0.260869565217391\\
0.00884955752212389 0.0434782608695652\\
0.0265486725663717 0.130434782608696\\
0.47787610619469 0.826086956521739\\
0.911504424778761 0.0434782608695652\\
0.946902654867257 0.217391304347826\\
0.986725663716814 0.0869565217391304\\
0.0221238938053097 0.0869565217391304\\
0.0353982300884956 0.0434782608695652\\
0.265486725663717 0.782608695652174\\
0.694690265486726 0.739130434782609\\
0.924778761061947 0.0869565217391304\\
0.95575221238938 0.173913043478261\\
0.982300884955752 0.0434782608695652\\
0.991150442477876 0.0434782608695652\\
0.0176991150442478 0.0434782608695652\\
0.0265486725663717 0.0434782608695652\\
0.261061946902655 0.739130434782609\\
0.47787610619469 0.0434782608695652\\
0.486725663716814 0.0434782608695652\\
0.699115044247788 0.695652173913043\\
0.920353982300885 0.0434782608695652\\
0.929203539823009 0.0434782608695652\\
0.942477876106195 0.0869565217391304\\
0.964601769911504 0.130434782608696\\
0.185840707964602 0.695652173913043\\
0.398230088495575 0.347826086956522\\
0.561946902654867 0.434782608695652\\
0.769911504424779 0.652173913043478\\
0.938053097345133 0.0434782608695652\\
0.946902654867257 0.0434782608695652\\
0.960176991150442 0.0869565217391304\\
0.973451327433628 0.0434782608695652\\
0.115044247787611 0.652173913043478\\
0.252212389380531 0.391304347826087\\
0.358407079646018 0.304347826086957\\
0.433628318584071 0.304347826086957\\
0.548672566371681 0.391304347826087\\
0.619469026548673 0.130434782608696\\
0.63716814159292 0.0434782608695652\\
0.774336283185841 0.608695652173913\\
0.95575221238938 0.0434782608695652\\
0.964601769911504 0.0434782608695652\\
0.0530973451327434 0.0434782608695652\\
0.119469026548673 0.608695652173913\\
0.20353982300885 0.217391304347826\\
0.274336283185841 0.347826086956522\\
0.327433628318584 0.0434782608695652\\
0.36283185840708 0.260869565217391\\
0.429203539823009 0.260869565217391\\
0.469026548672566 0.0434782608695652\\
0.513274336283186 0.217391304347826\\
0.570796460176991 0.347826086956522\\
0.610619469026549 0.0434782608695652\\
0.623893805309734 0.0869565217391304\\
0.668141592920354 0.217391304347826\\
0.800884955752212 0.565217391304348\\
0.0619469026548673 0.0434782608695652\\
0.123893805309735 0.565217391304348\\
0.185840707964602 0.0434782608695652\\
0.207964601769911 0.173913043478261\\
0.243362831858407 0.173913043478261\\
0.292035398230089 0.304347826086957\\
0.358407079646018 0.217391304347826\\
0.389380530973451 0.0434782608695652\\
0.415929203539823 0.217391304347826\\
0.451327433628319 0.130434782608696\\
0.508849557522124 0.173913043478261\\
0.530973451327434 0.0434782608695652\\
0.566371681415929 0.304347826086957\\
0.601769911504425 0.0434782608695652\\
0.619469026548673 0.0434782608695652\\
0.628318584070796 0.0434782608695652\\
0.65929203539823 0.173913043478261\\
0.685840707964602 0.0869565217391304\\
0.787610619469027 0.521739130434783\\
0.893805309734513 0.130434782608696\\
0.119469026548673 0.521739130434783\\
0.176991150442478 0.0434782608695652\\
0.20353982300885 0.130434782608696\\
0.221238938053097 0.0434782608695652\\
0.230088495575221 0.0434782608695652\\
0.247787610619469 0.130434782608696\\
0.265486725663717 0.0434782608695652\\
0.29646017699115 0.260869565217391\\
0.336283185840708 0.0434782608695652\\
0.36283185840708 0.173913043478261\\
0.411504424778761 0.173913043478261\\
0.433628318584071 0.0434782608695652\\
0.446902654867257 0.0869565217391304\\
0.460176991150442 0.0434782608695652\\
0.495575221238938 0.0434782608695652\\
0.513274336283186 0.130434782608696\\
0.561946902654867 0.260869565217391\\
0.592920353982301 0.0434782608695652\\
0.654867256637168 0.130434782608696\\
0.672566371681416 0.0434782608695652\\
0.68141592920354 0.0434782608695652\\
0.690265486725664 0.0434782608695652\\
0.783185840707965 0.478260869565217\\
0.876106194690266 0.0434782608695652\\
0.884955752212389 0.0434782608695652\\
0.898230088495575 0.0869565217391304\\
0.115044247787611 0.478260869565217\\
0.168141592920354 0.0434782608695652\\
0.194690265486726 0.0434782608695652\\
0.207964601769911 0.0869565217391304\\
0.243362831858407 0.0869565217391304\\
0.256637168141593 0.0434782608695652\\
0.292035398230089 0.217391304347826\\
0.31858407079646 0.0434782608695652\\
0.353982300884956 0.130434782608696\\
0.376106194690265 0.0869565217391304\\
0.407079646017699 0.130434782608696\\
0.424778761061947 0.0434782608695652\\
0.442477876106195 0.0434782608695652\\
0.451327433628319 0.0434782608695652\\
0.508849557522124 0.0869565217391304\\
0.52212389380531 0.0434782608695652\\
0.539823008849557 0.0434782608695652\\
0.566371681415929 0.217391304347826\\
0.650442477876106 0.0869565217391304\\
0.663716814159292 0.0434782608695652\\
0.752212389380531 0.434782608695652\\
0.84070796460177 0.304347826086957\\
0.893805309734513 0.0434782608695652\\
0.902654867256637 0.0434782608695652\\
0.110619469026549 0.434782608695652\\
0.15929203539823 0.0434782608695652\\
0.20353982300885 0.0434782608695652\\
0.212389380530973 0.0434782608695652\\
0.238938053097345 0.0434782608695652\\
0.247787610619469 0.0434782608695652\\
0.287610619469027 0.173913043478261\\
0.309734513274336 0.0434782608695652\\
0.349557522123894 0.0869565217391304\\
0.36283185840708 0.0434782608695652\\
0.371681415929204 0.0434782608695652\\
0.380530973451327 0.0434782608695652\\
0.402654867256637 0.0869565217391304\\
0.415929203539823 0.0434782608695652\\
0.504424778761062 0.0434782608695652\\
0.513274336283186 0.0434782608695652\\
0.548672566371681 0.0434782608695652\\
0.570796460176991 0.173913043478261\\
0.646017699115044 0.0434782608695652\\
0.654867256637168 0.0434782608695652\\
0.738938053097345 0.391304347826087\\
0.79646017699115 0.130434782608696\\
0.814159292035398 0.0434782608695652\\
0.845132743362832 0.260869565217391\\
0.106194690265487 0.391304347826087\\
0.150442477876106 0.0434782608695652\\
0.283185840707965 0.130434782608696\\
0.300884955752212 0.0434782608695652\\
0.345132743362832 0.0434782608695652\\
0.353982300884956 0.0434782608695652\\
0.398230088495575 0.0434782608695652\\
0.407079646017699 0.0434782608695652\\
0.557522123893805 0.0434782608695652\\
0.575221238938053 0.130434782608696\\
0.734513274336283 0.347826086956522\\
0.778761061946903 0.0434782608695652\\
0.787610619469027 0.0434782608695652\\
0.800884955752212 0.0869565217391304\\
0.823008849557522 0.0434782608695652\\
0.849557522123894 0.217391304347826\\
0.101769911504425 0.347826086956522\\
0.141592920353982 0.0434782608695652\\
0.278761061946903 0.0869565217391304\\
0.292035398230089 0.0434782608695652\\
0.566371681415929 0.0434782608695652\\
0.579646017699115 0.0869565217391304\\
0.730088495575221 0.304347826086957\\
0.769911504424779 0.0434782608695652\\
0.79646017699115 0.0434782608695652\\
0.805309734513274 0.0434782608695652\\
0.831858407079646 0.0434782608695652\\
0.853982300884956 0.173913043478261\\
0.0973451327433628 0.304347826086957\\
0.132743362831858 0.0434782608695652\\
0.274336283185841 0.0434782608695652\\
0.283185840707965 0.0434782608695652\\
0.575221238938053 0.0434782608695652\\
0.584070796460177 0.0434782608695652\\
0.725663716814159 0.260869565217391\\
0.761061946902655 0.0434782608695652\\
0.849557522123894 0.130434782608696\\
0.867256637168142 0.0434782608695652\\
0.0929203539823009 0.260869565217391\\
0.123893805309735 0.0434782608695652\\
0.721238938053097 0.217391304347826\\
0.752212389380531 0.0434782608695652\\
0.845132743362832 0.0869565217391304\\
0.858407079646018 0.0434782608695652\\
0.0884955752212389 0.217391304347826\\
0.115044247787611 0.0434782608695652\\
0.712389380530973 0.173913043478261\\
0.738938053097345 0.0869565217391304\\
0.84070796460177 0.0434782608695652\\
0.849557522123894 0.0434782608695652\\
0.084070796460177 0.173913043478261\\
0.106194690265487 0.0434782608695652\\
0.707964601769911 0.130434782608696\\
0.725663716814159 0.0434782608695652\\
0.734513274336283 0.0434782608695652\\
0.743362831858407 0.0434782608695652\\
0.079646017699115 0.130434782608696\\
0.0973451327433628 0.0434782608695652\\
0.70353982300885 0.0869565217391304\\
0.716814159292035 0.0434782608695652\\
0.0752212389380531 0.0869565217391304\\
0.0884955752212389 0.0434782608695652\\
0.699115044247788 0.0434782608695652\\
0.707964601769911 0.0434782608695652\\
0.0707964601769911 0.0434782608695652\\
0.079646017699115 0.0434782608695652\\
};
\addplot [
color=ohmred,
solid,
forget plot
]
table[row sep=crcr]{
0.0265486725663717 0.217391304347826\\
0.5 0.956521739130435\\
};
\addplot [
color=ohmred,
solid,
forget plot
]
table[row sep=crcr]{
0.52212389380531 0.91304347826087\\
0.5 0.956521739130435\\
};
\addplot [
color=ohmred,
solid,
forget plot
]
table[row sep=crcr]{
0.0221238938053097 0.173913043478261\\
0.0265486725663717 0.217391304347826\\
};
\addplot [
color=ohmred,
solid,
forget plot
]
table[row sep=crcr]{
0.0442477876106195 0.0434782608695652\\
0.0265486725663717 0.217391304347826\\
};
\addplot [
color=ohmred,
solid,
forget plot
]
table[row sep=crcr]{
0.482300884955752 0.869565217391304\\
0.52212389380531 0.91304347826087\\
};
\addplot [
color=ohmred,
solid,
forget plot
]
table[row sep=crcr]{
0.95575221238938 0.260869565217391\\
0.52212389380531 0.91304347826087\\
};
\addplot [
color=ohmred,
solid,
forget plot
]
table[row sep=crcr]{
0.00884955752212389 0.0434782608695652\\
0.0221238938053097 0.173913043478261\\
};
\addplot [
color=ohmred,
solid,
forget plot
]
table[row sep=crcr]{
0.0265486725663717 0.130434782608696\\
0.0221238938053097 0.173913043478261\\
};
\addplot [
color=ohmred,
solid,
forget plot
]
table[row sep=crcr]{
0.47787610619469 0.826086956521739\\
0.482300884955752 0.869565217391304\\
};
\addplot [
color=ohmred,
solid,
forget plot
]
table[row sep=crcr]{
0.911504424778761 0.0434782608695652\\
0.482300884955752 0.869565217391304\\
};
\addplot [
color=ohmred,
solid,
forget plot
]
table[row sep=crcr]{
0.946902654867257 0.217391304347826\\
0.95575221238938 0.260869565217391\\
};
\addplot [
color=ohmred,
solid,
forget plot
]
table[row sep=crcr]{
0.986725663716814 0.0869565217391304\\
0.95575221238938 0.260869565217391\\
};
\addplot [
color=ohmred,
solid,
forget plot
]
table[row sep=crcr]{
0.0221238938053097 0.0869565217391304\\
0.0265486725663717 0.130434782608696\\
};
\addplot [
color=ohmred,
solid,
forget plot
]
table[row sep=crcr]{
0.0353982300884956 0.0434782608695652\\
0.0265486725663717 0.130434782608696\\
};
\addplot [
color=ohmred,
solid,
forget plot
]
table[row sep=crcr]{
0.265486725663717 0.782608695652174\\
0.47787610619469 0.826086956521739\\
};
\addplot [
color=ohmred,
solid,
forget plot
]
table[row sep=crcr]{
0.694690265486726 0.739130434782609\\
0.47787610619469 0.826086956521739\\
};
\addplot [
color=ohmred,
solid,
forget plot
]
table[row sep=crcr]{
0.924778761061947 0.0869565217391304\\
0.946902654867257 0.217391304347826\\
};
\addplot [
color=ohmred,
solid,
forget plot
]
table[row sep=crcr]{
0.95575221238938 0.173913043478261\\
0.946902654867257 0.217391304347826\\
};
\addplot [
color=ohmred,
solid,
forget plot
]
table[row sep=crcr]{
0.982300884955752 0.0434782608695652\\
0.986725663716814 0.0869565217391304\\
};
\addplot [
color=ohmred,
solid,
forget plot
]
table[row sep=crcr]{
0.991150442477876 0.0434782608695652\\
0.986725663716814 0.0869565217391304\\
};
\addplot [
color=ohmred,
solid,
forget plot
]
table[row sep=crcr]{
0.0176991150442478 0.0434782608695652\\
0.0221238938053097 0.0869565217391304\\
};
\addplot [
color=ohmred,
solid,
forget plot
]
table[row sep=crcr]{
0.0265486725663717 0.0434782608695652\\
0.0221238938053097 0.0869565217391304\\
};
\addplot [
color=ohmred,
solid,
forget plot
]
table[row sep=crcr]{
0.261061946902655 0.739130434782609\\
0.265486725663717 0.782608695652174\\
};
\addplot [
color=ohmred,
solid,
forget plot
]
table[row sep=crcr]{
0.47787610619469 0.0434782608695652\\
0.265486725663717 0.782608695652174\\
};
\addplot [
color=ohmred,
solid,
forget plot
]
table[row sep=crcr]{
0.486725663716814 0.0434782608695652\\
0.694690265486726 0.739130434782609\\
};
\addplot [
color=ohmred,
solid,
forget plot
]
table[row sep=crcr]{
0.699115044247788 0.695652173913043\\
0.694690265486726 0.739130434782609\\
};
\addplot [
color=ohmred,
solid,
forget plot
]
table[row sep=crcr]{
0.920353982300885 0.0434782608695652\\
0.924778761061947 0.0869565217391304\\
};
\addplot [
color=ohmred,
solid,
forget plot
]
table[row sep=crcr]{
0.929203539823009 0.0434782608695652\\
0.924778761061947 0.0869565217391304\\
};
\addplot [
color=ohmred,
solid,
forget plot
]
table[row sep=crcr]{
0.942477876106195 0.0869565217391304\\
0.95575221238938 0.173913043478261\\
};
\addplot [
color=ohmred,
solid,
forget plot
]
table[row sep=crcr]{
0.964601769911504 0.130434782608696\\
0.95575221238938 0.173913043478261\\
};
\addplot [
color=ohmred,
solid,
forget plot
]
table[row sep=crcr]{
0.185840707964602 0.695652173913043\\
0.261061946902655 0.739130434782609\\
};
\addplot [
color=ohmred,
solid,
forget plot
]
table[row sep=crcr]{
0.398230088495575 0.347826086956522\\
0.261061946902655 0.739130434782609\\
};
\addplot [
color=ohmred,
solid,
forget plot
]
table[row sep=crcr]{
0.561946902654867 0.434782608695652\\
0.699115044247788 0.695652173913043\\
};
\addplot [
color=ohmred,
solid,
forget plot
]
table[row sep=crcr]{
0.769911504424779 0.652173913043478\\
0.699115044247788 0.695652173913043\\
};
\addplot [
color=ohmred,
solid,
forget plot
]
table[row sep=crcr]{
0.938053097345133 0.0434782608695652\\
0.942477876106195 0.0869565217391304\\
};
\addplot [
color=ohmred,
solid,
forget plot
]
table[row sep=crcr]{
0.946902654867257 0.0434782608695652\\
0.942477876106195 0.0869565217391304\\
};
\addplot [
color=ohmred,
solid,
forget plot
]
table[row sep=crcr]{
0.960176991150442 0.0869565217391304\\
0.964601769911504 0.130434782608696\\
};
\addplot [
color=ohmred,
solid,
forget plot
]
table[row sep=crcr]{
0.973451327433628 0.0434782608695652\\
0.964601769911504 0.130434782608696\\
};
\addplot [
color=ohmred,
solid,
forget plot
]
table[row sep=crcr]{
0.115044247787611 0.652173913043478\\
0.185840707964602 0.695652173913043\\
};
\addplot [
color=ohmred,
solid,
forget plot
]
table[row sep=crcr]{
0.252212389380531 0.391304347826087\\
0.185840707964602 0.695652173913043\\
};
\addplot [
color=ohmred,
solid,
forget plot
]
table[row sep=crcr]{
0.358407079646018 0.304347826086957\\
0.398230088495575 0.347826086956522\\
};
\addplot [
color=ohmred,
solid,
forget plot
]
table[row sep=crcr]{
0.433628318584071 0.304347826086957\\
0.398230088495575 0.347826086956522\\
};
\addplot [
color=ohmred,
solid,
forget plot
]
table[row sep=crcr]{
0.548672566371681 0.391304347826087\\
0.561946902654867 0.434782608695652\\
};
\addplot [
color=ohmred,
solid,
forget plot
]
table[row sep=crcr]{
0.619469026548673 0.130434782608696\\
0.561946902654867 0.434782608695652\\
};
\addplot [
color=ohmred,
solid,
forget plot
]
table[row sep=crcr]{
0.63716814159292 0.0434782608695652\\
0.769911504424779 0.652173913043478\\
};
\addplot [
color=ohmred,
solid,
forget plot
]
table[row sep=crcr]{
0.774336283185841 0.608695652173913\\
0.769911504424779 0.652173913043478\\
};
\addplot [
color=ohmred,
solid,
forget plot
]
table[row sep=crcr]{
0.95575221238938 0.0434782608695652\\
0.960176991150442 0.0869565217391304\\
};
\addplot [
color=ohmred,
solid,
forget plot
]
table[row sep=crcr]{
0.964601769911504 0.0434782608695652\\
0.960176991150442 0.0869565217391304\\
};
\addplot [
color=ohmred,
solid,
forget plot
]
table[row sep=crcr]{
0.0530973451327434 0.0434782608695652\\
0.115044247787611 0.652173913043478\\
};
\addplot [
color=ohmred,
solid,
forget plot
]
table[row sep=crcr]{
0.119469026548673 0.608695652173913\\
0.115044247787611 0.652173913043478\\
};
\addplot [
color=ohmred,
solid,
forget plot
]
table[row sep=crcr]{
0.20353982300885 0.217391304347826\\
0.252212389380531 0.391304347826087\\
};
\addplot [
color=ohmred,
solid,
forget plot
]
table[row sep=crcr]{
0.274336283185841 0.347826086956522\\
0.252212389380531 0.391304347826087\\
};
\addplot [
color=ohmred,
solid,
forget plot
]
table[row sep=crcr]{
0.327433628318584 0.0434782608695652\\
0.358407079646018 0.304347826086957\\
};
\addplot [
color=ohmred,
solid,
forget plot
]
table[row sep=crcr]{
0.36283185840708 0.260869565217391\\
0.358407079646018 0.304347826086957\\
};
\addplot [
color=ohmred,
solid,
forget plot
]
table[row sep=crcr]{
0.429203539823009 0.260869565217391\\
0.433628318584071 0.304347826086957\\
};
\addplot [
color=ohmred,
solid,
forget plot
]
table[row sep=crcr]{
0.469026548672566 0.0434782608695652\\
0.433628318584071 0.304347826086957\\
};
\addplot [
color=ohmred,
solid,
forget plot
]
table[row sep=crcr]{
0.513274336283186 0.217391304347826\\
0.548672566371681 0.391304347826087\\
};
\addplot [
color=ohmred,
solid,
forget plot
]
table[row sep=crcr]{
0.570796460176991 0.347826086956522\\
0.548672566371681 0.391304347826087\\
};
\addplot [
color=ohmred,
solid,
forget plot
]
table[row sep=crcr]{
0.610619469026549 0.0434782608695652\\
0.619469026548673 0.130434782608696\\
};
\addplot [
color=ohmred,
solid,
forget plot
]
table[row sep=crcr]{
0.623893805309734 0.0869565217391304\\
0.619469026548673 0.130434782608696\\
};
\addplot [
color=ohmred,
solid,
forget plot
]
table[row sep=crcr]{
0.668141592920354 0.217391304347826\\
0.774336283185841 0.608695652173913\\
};
\addplot [
color=ohmred,
solid,
forget plot
]
table[row sep=crcr]{
0.800884955752212 0.565217391304348\\
0.774336283185841 0.608695652173913\\
};
\addplot [
color=ohmred,
solid,
forget plot
]
table[row sep=crcr]{
0.0619469026548673 0.0434782608695652\\
0.119469026548673 0.608695652173913\\
};
\addplot [
color=ohmred,
solid,
forget plot
]
table[row sep=crcr]{
0.123893805309735 0.565217391304348\\
0.119469026548673 0.608695652173913\\
};
\addplot [
color=ohmred,
solid,
forget plot
]
table[row sep=crcr]{
0.185840707964602 0.0434782608695652\\
0.20353982300885 0.217391304347826\\
};
\addplot [
color=ohmred,
solid,
forget plot
]
table[row sep=crcr]{
0.207964601769911 0.173913043478261\\
0.20353982300885 0.217391304347826\\
};
\addplot [
color=ohmred,
solid,
forget plot
]
table[row sep=crcr]{
0.243362831858407 0.173913043478261\\
0.274336283185841 0.347826086956522\\
};
\addplot [
color=ohmred,
solid,
forget plot
]
table[row sep=crcr]{
0.292035398230089 0.304347826086957\\
0.274336283185841 0.347826086956522\\
};
\addplot [
color=ohmred,
solid,
forget plot
]
table[row sep=crcr]{
0.358407079646018 0.217391304347826\\
0.36283185840708 0.260869565217391\\
};
\addplot [
color=ohmred,
solid,
forget plot
]
table[row sep=crcr]{
0.389380530973451 0.0434782608695652\\
0.36283185840708 0.260869565217391\\
};
\addplot [
color=ohmred,
solid,
forget plot
]
table[row sep=crcr]{
0.415929203539823 0.217391304347826\\
0.429203539823009 0.260869565217391\\
};
\addplot [
color=ohmred,
solid,
forget plot
]
table[row sep=crcr]{
0.451327433628319 0.130434782608696\\
0.429203539823009 0.260869565217391\\
};
\addplot [
color=ohmred,
solid,
forget plot
]
table[row sep=crcr]{
0.508849557522124 0.173913043478261\\
0.513274336283186 0.217391304347826\\
};
\addplot [
color=ohmred,
solid,
forget plot
]
table[row sep=crcr]{
0.530973451327434 0.0434782608695652\\
0.513274336283186 0.217391304347826\\
};
\addplot [
color=ohmred,
solid,
forget plot
]
table[row sep=crcr]{
0.566371681415929 0.304347826086957\\
0.570796460176991 0.347826086956522\\
};
\addplot [
color=ohmred,
solid,
forget plot
]
table[row sep=crcr]{
0.601769911504425 0.0434782608695652\\
0.570796460176991 0.347826086956522\\
};
\addplot [
color=ohmred,
solid,
forget plot
]
table[row sep=crcr]{
0.619469026548673 0.0434782608695652\\
0.623893805309734 0.0869565217391304\\
};
\addplot [
color=ohmred,
solid,
forget plot
]
table[row sep=crcr]{
0.628318584070796 0.0434782608695652\\
0.623893805309734 0.0869565217391304\\
};
\addplot [
color=ohmred,
solid,
forget plot
]
table[row sep=crcr]{
0.65929203539823 0.173913043478261\\
0.668141592920354 0.217391304347826\\
};
\addplot [
color=ohmred,
solid,
forget plot
]
table[row sep=crcr]{
0.685840707964602 0.0869565217391304\\
0.668141592920354 0.217391304347826\\
};
\addplot [
color=ohmred,
solid,
forget plot
]
table[row sep=crcr]{
0.787610619469027 0.521739130434783\\
0.800884955752212 0.565217391304348\\
};
\addplot [
color=ohmred,
solid,
forget plot
]
table[row sep=crcr]{
0.893805309734513 0.130434782608696\\
0.800884955752212 0.565217391304348\\
};
\addplot [
color=ohmred,
solid,
forget plot
]
table[row sep=crcr]{
0.119469026548673 0.521739130434783\\
0.123893805309735 0.565217391304348\\
};
\addplot [
color=ohmred,
solid,
forget plot
]
table[row sep=crcr]{
0.176991150442478 0.0434782608695652\\
0.123893805309735 0.565217391304348\\
};
\addplot [
color=ohmred,
solid,
forget plot
]
table[row sep=crcr]{
0.20353982300885 0.130434782608696\\
0.207964601769911 0.173913043478261\\
};
\addplot [
color=ohmred,
solid,
forget plot
]
table[row sep=crcr]{
0.221238938053097 0.0434782608695652\\
0.207964601769911 0.173913043478261\\
};
\addplot [
color=ohmred,
solid,
forget plot
]
table[row sep=crcr]{
0.230088495575221 0.0434782608695652\\
0.243362831858407 0.173913043478261\\
};
\addplot [
color=ohmred,
solid,
forget plot
]
table[row sep=crcr]{
0.247787610619469 0.130434782608696\\
0.243362831858407 0.173913043478261\\
};
\addplot [
color=ohmred,
solid,
forget plot
]
table[row sep=crcr]{
0.265486725663717 0.0434782608695652\\
0.292035398230089 0.304347826086957\\
};
\addplot [
color=ohmred,
solid,
forget plot
]
table[row sep=crcr]{
0.29646017699115 0.260869565217391\\
0.292035398230089 0.304347826086957\\
};
\addplot [
color=ohmred,
solid,
forget plot
]
table[row sep=crcr]{
0.336283185840708 0.0434782608695652\\
0.358407079646018 0.217391304347826\\
};
\addplot [
color=ohmred,
solid,
forget plot
]
table[row sep=crcr]{
0.36283185840708 0.173913043478261\\
0.358407079646018 0.217391304347826\\
};
\addplot [
color=ohmred,
solid,
forget plot
]
table[row sep=crcr]{
0.411504424778761 0.173913043478261\\
0.415929203539823 0.217391304347826\\
};
\addplot [
color=ohmred,
solid,
forget plot
]
table[row sep=crcr]{
0.433628318584071 0.0434782608695652\\
0.415929203539823 0.217391304347826\\
};
\addplot [
color=ohmred,
solid,
forget plot
]
table[row sep=crcr]{
0.446902654867257 0.0869565217391304\\
0.451327433628319 0.130434782608696\\
};
\addplot [
color=ohmred,
solid,
forget plot
]
table[row sep=crcr]{
0.460176991150442 0.0434782608695652\\
0.451327433628319 0.130434782608696\\
};
\addplot [
color=ohmred,
solid,
forget plot
]
table[row sep=crcr]{
0.495575221238938 0.0434782608695652\\
0.508849557522124 0.173913043478261\\
};
\addplot [
color=ohmred,
solid,
forget plot
]
table[row sep=crcr]{
0.513274336283186 0.130434782608696\\
0.508849557522124 0.173913043478261\\
};
\addplot [
color=ohmred,
solid,
forget plot
]
table[row sep=crcr]{
0.561946902654867 0.260869565217391\\
0.566371681415929 0.304347826086957\\
};
\addplot [
color=ohmred,
solid,
forget plot
]
table[row sep=crcr]{
0.592920353982301 0.0434782608695652\\
0.566371681415929 0.304347826086957\\
};
\addplot [
color=ohmred,
solid,
forget plot
]
table[row sep=crcr]{
0.654867256637168 0.130434782608696\\
0.65929203539823 0.173913043478261\\
};
\addplot [
color=ohmred,
solid,
forget plot
]
table[row sep=crcr]{
0.672566371681416 0.0434782608695652\\
0.65929203539823 0.173913043478261\\
};
\addplot [
color=ohmred,
solid,
forget plot
]
table[row sep=crcr]{
0.68141592920354 0.0434782608695652\\
0.685840707964602 0.0869565217391304\\
};
\addplot [
color=ohmred,
solid,
forget plot
]
table[row sep=crcr]{
0.690265486725664 0.0434782608695652\\
0.685840707964602 0.0869565217391304\\
};
\addplot [
color=ohmred,
solid,
forget plot
]
table[row sep=crcr]{
0.783185840707965 0.478260869565217\\
0.787610619469027 0.521739130434783\\
};
\addplot [
color=ohmred,
solid,
forget plot
]
table[row sep=crcr]{
0.876106194690266 0.0434782608695652\\
0.787610619469027 0.521739130434783\\
};
\addplot [
color=ohmred,
solid,
forget plot
]
table[row sep=crcr]{
0.884955752212389 0.0434782608695652\\
0.893805309734513 0.130434782608696\\
};
\addplot [
color=ohmred,
solid,
forget plot
]
table[row sep=crcr]{
0.898230088495575 0.0869565217391304\\
0.893805309734513 0.130434782608696\\
};
\addplot [
color=ohmred,
solid,
forget plot
]
table[row sep=crcr]{
0.115044247787611 0.478260869565217\\
0.119469026548673 0.521739130434783\\
};
\addplot [
color=ohmred,
solid,
forget plot
]
table[row sep=crcr]{
0.168141592920354 0.0434782608695652\\
0.119469026548673 0.521739130434783\\
};
\addplot [
color=ohmred,
solid,
forget plot
]
table[row sep=crcr]{
0.194690265486726 0.0434782608695652\\
0.20353982300885 0.130434782608696\\
};
\addplot [
color=ohmred,
solid,
forget plot
]
table[row sep=crcr]{
0.207964601769911 0.0869565217391304\\
0.20353982300885 0.130434782608696\\
};
\addplot [
color=ohmred,
solid,
forget plot
]
table[row sep=crcr]{
0.243362831858407 0.0869565217391304\\
0.247787610619469 0.130434782608696\\
};
\addplot [
color=ohmred,
solid,
forget plot
]
table[row sep=crcr]{
0.256637168141593 0.0434782608695652\\
0.247787610619469 0.130434782608696\\
};
\addplot [
color=ohmred,
solid,
forget plot
]
table[row sep=crcr]{
0.292035398230089 0.217391304347826\\
0.29646017699115 0.260869565217391\\
};
\addplot [
color=ohmred,
solid,
forget plot
]
table[row sep=crcr]{
0.31858407079646 0.0434782608695652\\
0.29646017699115 0.260869565217391\\
};
\addplot [
color=ohmred,
solid,
forget plot
]
table[row sep=crcr]{
0.353982300884956 0.130434782608696\\
0.36283185840708 0.173913043478261\\
};
\addplot [
color=ohmred,
solid,
forget plot
]
table[row sep=crcr]{
0.376106194690265 0.0869565217391304\\
0.36283185840708 0.173913043478261\\
};
\addplot [
color=ohmred,
solid,
forget plot
]
table[row sep=crcr]{
0.407079646017699 0.130434782608696\\
0.411504424778761 0.173913043478261\\
};
\addplot [
color=ohmred,
solid,
forget plot
]
table[row sep=crcr]{
0.424778761061947 0.0434782608695652\\
0.411504424778761 0.173913043478261\\
};
\addplot [
color=ohmred,
solid,
forget plot
]
table[row sep=crcr]{
0.442477876106195 0.0434782608695652\\
0.446902654867257 0.0869565217391304\\
};
\addplot [
color=ohmred,
solid,
forget plot
]
table[row sep=crcr]{
0.451327433628319 0.0434782608695652\\
0.446902654867257 0.0869565217391304\\
};
\addplot [
color=ohmred,
solid,
forget plot
]
table[row sep=crcr]{
0.508849557522124 0.0869565217391304\\
0.513274336283186 0.130434782608696\\
};
\addplot [
color=ohmred,
solid,
forget plot
]
table[row sep=crcr]{
0.52212389380531 0.0434782608695652\\
0.513274336283186 0.130434782608696\\
};
\addplot [
color=ohmred,
solid,
forget plot
]
table[row sep=crcr]{
0.539823008849557 0.0434782608695652\\
0.561946902654867 0.260869565217391\\
};
\addplot [
color=ohmred,
solid,
forget plot
]
table[row sep=crcr]{
0.566371681415929 0.217391304347826\\
0.561946902654867 0.260869565217391\\
};
\addplot [
color=ohmred,
solid,
forget plot
]
table[row sep=crcr]{
0.650442477876106 0.0869565217391304\\
0.654867256637168 0.130434782608696\\
};
\addplot [
color=ohmred,
solid,
forget plot
]
table[row sep=crcr]{
0.663716814159292 0.0434782608695652\\
0.654867256637168 0.130434782608696\\
};
\addplot [
color=ohmred,
solid,
forget plot
]
table[row sep=crcr]{
0.752212389380531 0.434782608695652\\
0.783185840707965 0.478260869565217\\
};
\addplot [
color=ohmred,
solid,
forget plot
]
table[row sep=crcr]{
0.84070796460177 0.304347826086957\\
0.783185840707965 0.478260869565217\\
};
\addplot [
color=ohmred,
solid,
forget plot
]
table[row sep=crcr]{
0.893805309734513 0.0434782608695652\\
0.898230088495575 0.0869565217391304\\
};
\addplot [
color=ohmred,
solid,
forget plot
]
table[row sep=crcr]{
0.902654867256637 0.0434782608695652\\
0.898230088495575 0.0869565217391304\\
};
\addplot [
color=ohmred,
solid,
forget plot
]
table[row sep=crcr]{
0.110619469026549 0.434782608695652\\
0.115044247787611 0.478260869565217\\
};
\addplot [
color=ohmred,
solid,
forget plot
]
table[row sep=crcr]{
0.15929203539823 0.0434782608695652\\
0.115044247787611 0.478260869565217\\
};
\addplot [
color=ohmred,
solid,
forget plot
]
table[row sep=crcr]{
0.20353982300885 0.0434782608695652\\
0.207964601769911 0.0869565217391304\\
};
\addplot [
color=ohmred,
solid,
forget plot
]
table[row sep=crcr]{
0.212389380530973 0.0434782608695652\\
0.207964601769911 0.0869565217391304\\
};
\addplot [
color=ohmred,
solid,
forget plot
]
table[row sep=crcr]{
0.238938053097345 0.0434782608695652\\
0.243362831858407 0.0869565217391304\\
};
\addplot [
color=ohmred,
solid,
forget plot
]
table[row sep=crcr]{
0.247787610619469 0.0434782608695652\\
0.243362831858407 0.0869565217391304\\
};
\addplot [
color=ohmred,
solid,
forget plot
]
table[row sep=crcr]{
0.287610619469027 0.173913043478261\\
0.292035398230089 0.217391304347826\\
};
\addplot [
color=ohmred,
solid,
forget plot
]
table[row sep=crcr]{
0.309734513274336 0.0434782608695652\\
0.292035398230089 0.217391304347826\\
};
\addplot [
color=ohmred,
solid,
forget plot
]
table[row sep=crcr]{
0.349557522123894 0.0869565217391304\\
0.353982300884956 0.130434782608696\\
};
\addplot [
color=ohmred,
solid,
forget plot
]
table[row sep=crcr]{
0.36283185840708 0.0434782608695652\\
0.353982300884956 0.130434782608696\\
};
\addplot [
color=ohmred,
solid,
forget plot
]
table[row sep=crcr]{
0.371681415929204 0.0434782608695652\\
0.376106194690265 0.0869565217391304\\
};
\addplot [
color=ohmred,
solid,
forget plot
]
table[row sep=crcr]{
0.380530973451327 0.0434782608695652\\
0.376106194690265 0.0869565217391304\\
};
\addplot [
color=ohmred,
solid,
forget plot
]
table[row sep=crcr]{
0.402654867256637 0.0869565217391304\\
0.407079646017699 0.130434782608696\\
};
\addplot [
color=ohmred,
solid,
forget plot
]
table[row sep=crcr]{
0.415929203539823 0.0434782608695652\\
0.407079646017699 0.130434782608696\\
};
\addplot [
color=ohmred,
solid,
forget plot
]
table[row sep=crcr]{
0.504424778761062 0.0434782608695652\\
0.508849557522124 0.0869565217391304\\
};
\addplot [
color=ohmred,
solid,
forget plot
]
table[row sep=crcr]{
0.513274336283186 0.0434782608695652\\
0.508849557522124 0.0869565217391304\\
};
\addplot [
color=ohmred,
solid,
forget plot
]
table[row sep=crcr]{
0.548672566371681 0.0434782608695652\\
0.566371681415929 0.217391304347826\\
};
\addplot [
color=ohmred,
solid,
forget plot
]
table[row sep=crcr]{
0.570796460176991 0.173913043478261\\
0.566371681415929 0.217391304347826\\
};
\addplot [
color=ohmred,
solid,
forget plot
]
table[row sep=crcr]{
0.646017699115044 0.0434782608695652\\
0.650442477876106 0.0869565217391304\\
};
\addplot [
color=ohmred,
solid,
forget plot
]
table[row sep=crcr]{
0.654867256637168 0.0434782608695652\\
0.650442477876106 0.0869565217391304\\
};
\addplot [
color=ohmred,
solid,
forget plot
]
table[row sep=crcr]{
0.738938053097345 0.391304347826087\\
0.752212389380531 0.434782608695652\\
};
\addplot [
color=ohmred,
solid,
forget plot
]
table[row sep=crcr]{
0.79646017699115 0.130434782608696\\
0.752212389380531 0.434782608695652\\
};
\addplot [
color=ohmred,
solid,
forget plot
]
table[row sep=crcr]{
0.814159292035398 0.0434782608695652\\
0.84070796460177 0.304347826086957\\
};
\addplot [
color=ohmred,
solid,
forget plot
]
table[row sep=crcr]{
0.845132743362832 0.260869565217391\\
0.84070796460177 0.304347826086957\\
};
\addplot [
color=ohmred,
solid,
forget plot
]
table[row sep=crcr]{
0.106194690265487 0.391304347826087\\
0.110619469026549 0.434782608695652\\
};
\addplot [
color=ohmred,
solid,
forget plot
]
table[row sep=crcr]{
0.150442477876106 0.0434782608695652\\
0.110619469026549 0.434782608695652\\
};
\addplot [
color=ohmred,
solid,
forget plot
]
table[row sep=crcr]{
0.283185840707965 0.130434782608696\\
0.287610619469027 0.173913043478261\\
};
\addplot [
color=ohmred,
solid,
forget plot
]
table[row sep=crcr]{
0.300884955752212 0.0434782608695652\\
0.287610619469027 0.173913043478261\\
};
\addplot [
color=ohmred,
solid,
forget plot
]
table[row sep=crcr]{
0.345132743362832 0.0434782608695652\\
0.349557522123894 0.0869565217391304\\
};
\addplot [
color=ohmred,
solid,
forget plot
]
table[row sep=crcr]{
0.353982300884956 0.0434782608695652\\
0.349557522123894 0.0869565217391304\\
};
\addplot [
color=ohmred,
solid,
forget plot
]
table[row sep=crcr]{
0.398230088495575 0.0434782608695652\\
0.402654867256637 0.0869565217391304\\
};
\addplot [
color=ohmred,
solid,
forget plot
]
table[row sep=crcr]{
0.407079646017699 0.0434782608695652\\
0.402654867256637 0.0869565217391304\\
};
\addplot [
color=ohmred,
solid,
forget plot
]
table[row sep=crcr]{
0.557522123893805 0.0434782608695652\\
0.570796460176991 0.173913043478261\\
};
\addplot [
color=ohmred,
solid,
forget plot
]
table[row sep=crcr]{
0.575221238938053 0.130434782608696\\
0.570796460176991 0.173913043478261\\
};
\addplot [
color=ohmred,
solid,
forget plot
]
table[row sep=crcr]{
0.734513274336283 0.347826086956522\\
0.738938053097345 0.391304347826087\\
};
\addplot [
color=ohmred,
solid,
forget plot
]
table[row sep=crcr]{
0.778761061946903 0.0434782608695652\\
0.738938053097345 0.391304347826087\\
};
\addplot [
color=ohmred,
solid,
forget plot
]
table[row sep=crcr]{
0.787610619469027 0.0434782608695652\\
0.79646017699115 0.130434782608696\\
};
\addplot [
color=ohmred,
solid,
forget plot
]
table[row sep=crcr]{
0.800884955752212 0.0869565217391304\\
0.79646017699115 0.130434782608696\\
};
\addplot [
color=ohmred,
solid,
forget plot
]
table[row sep=crcr]{
0.823008849557522 0.0434782608695652\\
0.845132743362832 0.260869565217391\\
};
\addplot [
color=ohmred,
solid,
forget plot
]
table[row sep=crcr]{
0.849557522123894 0.217391304347826\\
0.845132743362832 0.260869565217391\\
};
\addplot [
color=ohmred,
solid,
forget plot
]
table[row sep=crcr]{
0.101769911504425 0.347826086956522\\
0.106194690265487 0.391304347826087\\
};
\addplot [
color=ohmred,
solid,
forget plot
]
table[row sep=crcr]{
0.141592920353982 0.0434782608695652\\
0.106194690265487 0.391304347826087\\
};
\addplot [
color=ohmred,
solid,
forget plot
]
table[row sep=crcr]{
0.278761061946903 0.0869565217391304\\
0.283185840707965 0.130434782608696\\
};
\addplot [
color=ohmred,
solid,
forget plot
]
table[row sep=crcr]{
0.292035398230089 0.0434782608695652\\
0.283185840707965 0.130434782608696\\
};
\addplot [
color=ohmred,
solid,
forget plot
]
table[row sep=crcr]{
0.566371681415929 0.0434782608695652\\
0.575221238938053 0.130434782608696\\
};
\addplot [
color=ohmred,
solid,
forget plot
]
table[row sep=crcr]{
0.579646017699115 0.0869565217391304\\
0.575221238938053 0.130434782608696\\
};
\addplot [
color=ohmred,
solid,
forget plot
]
table[row sep=crcr]{
0.730088495575221 0.304347826086957\\
0.734513274336283 0.347826086956522\\
};
\addplot [
color=ohmred,
solid,
forget plot
]
table[row sep=crcr]{
0.769911504424779 0.0434782608695652\\
0.734513274336283 0.347826086956522\\
};
\addplot [
color=ohmred,
solid,
forget plot
]
table[row sep=crcr]{
0.79646017699115 0.0434782608695652\\
0.800884955752212 0.0869565217391304\\
};
\addplot [
color=ohmred,
solid,
forget plot
]
table[row sep=crcr]{
0.805309734513274 0.0434782608695652\\
0.800884955752212 0.0869565217391304\\
};
\addplot [
color=ohmred,
solid,
forget plot
]
table[row sep=crcr]{
0.831858407079646 0.0434782608695652\\
0.849557522123894 0.217391304347826\\
};
\addplot [
color=ohmred,
solid,
forget plot
]
table[row sep=crcr]{
0.853982300884956 0.173913043478261\\
0.849557522123894 0.217391304347826\\
};
\addplot [
color=ohmred,
solid,
forget plot
]
table[row sep=crcr]{
0.0973451327433628 0.304347826086957\\
0.101769911504425 0.347826086956522\\
};
\addplot [
color=ohmred,
solid,
forget plot
]
table[row sep=crcr]{
0.132743362831858 0.0434782608695652\\
0.101769911504425 0.347826086956522\\
};
\addplot [
color=ohmred,
solid,
forget plot
]
table[row sep=crcr]{
0.274336283185841 0.0434782608695652\\
0.278761061946903 0.0869565217391304\\
};
\addplot [
color=ohmred,
solid,
forget plot
]
table[row sep=crcr]{
0.283185840707965 0.0434782608695652\\
0.278761061946903 0.0869565217391304\\
};
\addplot [
color=ohmred,
solid,
forget plot
]
table[row sep=crcr]{
0.575221238938053 0.0434782608695652\\
0.579646017699115 0.0869565217391304\\
};
\addplot [
color=ohmred,
solid,
forget plot
]
table[row sep=crcr]{
0.584070796460177 0.0434782608695652\\
0.579646017699115 0.0869565217391304\\
};
\addplot [
color=ohmred,
solid,
forget plot
]
table[row sep=crcr]{
0.725663716814159 0.260869565217391\\
0.730088495575221 0.304347826086957\\
};
\addplot [
color=ohmred,
solid,
forget plot
]
table[row sep=crcr]{
0.761061946902655 0.0434782608695652\\
0.730088495575221 0.304347826086957\\
};
\addplot [
color=ohmred,
solid,
forget plot
]
table[row sep=crcr]{
0.849557522123894 0.130434782608696\\
0.853982300884956 0.173913043478261\\
};
\addplot [
color=ohmred,
solid,
forget plot
]
table[row sep=crcr]{
0.867256637168142 0.0434782608695652\\
0.853982300884956 0.173913043478261\\
};
\addplot [
color=ohmred,
solid,
forget plot
]
table[row sep=crcr]{
0.0929203539823009 0.260869565217391\\
0.0973451327433628 0.304347826086957\\
};
\addplot [
color=ohmred,
solid,
forget plot
]
table[row sep=crcr]{
0.123893805309735 0.0434782608695652\\
0.0973451327433628 0.304347826086957\\
};
\addplot [
color=ohmred,
solid,
forget plot
]
table[row sep=crcr]{
0.721238938053097 0.217391304347826\\
0.725663716814159 0.260869565217391\\
};
\addplot [
color=ohmred,
solid,
forget plot
]
table[row sep=crcr]{
0.752212389380531 0.0434782608695652\\
0.725663716814159 0.260869565217391\\
};
\addplot [
color=ohmred,
solid,
forget plot
]
table[row sep=crcr]{
0.845132743362832 0.0869565217391304\\
0.849557522123894 0.130434782608696\\
};
\addplot [
color=ohmred,
solid,
forget plot
]
table[row sep=crcr]{
0.858407079646018 0.0434782608695652\\
0.849557522123894 0.130434782608696\\
};
\addplot [
color=ohmred,
solid,
forget plot
]
table[row sep=crcr]{
0.0884955752212389 0.217391304347826\\
0.0929203539823009 0.260869565217391\\
};
\addplot [
color=ohmred,
solid,
forget plot
]
table[row sep=crcr]{
0.115044247787611 0.0434782608695652\\
0.0929203539823009 0.260869565217391\\
};
\addplot [
color=ohmred,
solid,
forget plot
]
table[row sep=crcr]{
0.712389380530973 0.173913043478261\\
0.721238938053097 0.217391304347826\\
};
\addplot [
color=ohmred,
solid,
forget plot
]
table[row sep=crcr]{
0.738938053097345 0.0869565217391304\\
0.721238938053097 0.217391304347826\\
};
\addplot [
color=ohmred,
solid,
forget plot
]
table[row sep=crcr]{
0.84070796460177 0.0434782608695652\\
0.845132743362832 0.0869565217391304\\
};
\addplot [
color=ohmred,
solid,
forget plot
]
table[row sep=crcr]{
0.849557522123894 0.0434782608695652\\
0.845132743362832 0.0869565217391304\\
};
\addplot [
color=ohmred,
solid,
forget plot
]
table[row sep=crcr]{
0.084070796460177 0.173913043478261\\
0.0884955752212389 0.217391304347826\\
};
\addplot [
color=ohmred,
solid,
forget plot
]
table[row sep=crcr]{
0.106194690265487 0.0434782608695652\\
0.0884955752212389 0.217391304347826\\
};
\addplot [
color=ohmred,
solid,
forget plot
]
table[row sep=crcr]{
0.707964601769911 0.130434782608696\\
0.712389380530973 0.173913043478261\\
};
\addplot [
color=ohmred,
solid,
forget plot
]
table[row sep=crcr]{
0.725663716814159 0.0434782608695652\\
0.712389380530973 0.173913043478261\\
};
\addplot [
color=ohmred,
solid,
forget plot
]
table[row sep=crcr]{
0.734513274336283 0.0434782608695652\\
0.738938053097345 0.0869565217391304\\
};
\addplot [
color=ohmred,
solid,
forget plot
]
table[row sep=crcr]{
0.743362831858407 0.0434782608695652\\
0.738938053097345 0.0869565217391304\\
};
\addplot [
color=ohmred,
solid,
forget plot
]
table[row sep=crcr]{
0.079646017699115 0.130434782608696\\
0.084070796460177 0.173913043478261\\
};
\addplot [
color=ohmred,
solid,
forget plot
]
table[row sep=crcr]{
0.0973451327433628 0.0434782608695652\\
0.084070796460177 0.173913043478261\\
};
\addplot [
color=ohmred,
solid,
forget plot
]
table[row sep=crcr]{
0.70353982300885 0.0869565217391304\\
0.707964601769911 0.130434782608696\\
};
\addplot [
color=ohmred,
solid,
forget plot
]
table[row sep=crcr]{
0.716814159292035 0.0434782608695652\\
0.707964601769911 0.130434782608696\\
};
\addplot [
color=ohmred,
solid,
forget plot
]
table[row sep=crcr]{
0.0752212389380531 0.0869565217391304\\
0.079646017699115 0.130434782608696\\
};
\addplot [
color=ohmred,
solid,
forget plot
]
table[row sep=crcr]{
0.0884955752212389 0.0434782608695652\\
0.079646017699115 0.130434782608696\\
};
\addplot [
color=ohmred,
solid,
forget plot
]
table[row sep=crcr]{
0.699115044247788 0.0434782608695652\\
0.70353982300885 0.0869565217391304\\
};
\addplot [
color=ohmred,
solid,
forget plot
]
table[row sep=crcr]{
0.707964601769911 0.0434782608695652\\
0.70353982300885 0.0869565217391304\\
};
\addplot [
color=ohmred,
solid,
forget plot
]
table[row sep=crcr]{
0.0707964601769911 0.0434782608695652\\
0.0752212389380531 0.0869565217391304\\
};
\addplot [
color=ohmred,
solid,
forget plot
]
table[row sep=crcr]{
0.079646017699115 0.0434782608695652\\
0.0752212389380531 0.0869565217391304\\
};
\end{axis}
\end{tikzpicture}% \\ Tiefe: 21 \\ Genauigkeit: 91{,}1 \% \pause
        \column[t]{0.45\textwidth} \centering \textbf{Bester Entscheidungsbaum} \\[1em] % This file was created by matlab2tikz v0.4.4 running on MATLAB 7.14.
% Copyright (c) 2008--2013, Nico Schlömer <nico.schloemer@gmail.com>
% All rights reserved.
% 
% The latest updates can be retrieved from
%   http://www.mathworks.com/matlabcentral/fileexchange/22022-matlab2tikz
% where you can also make suggestions and rate matlab2tikz.
% 
\begin{tikzpicture}

\begin{axis}[%
width=\figurewidth,
height=\figureheight,
scale only axis,
xmin=0,
xmax=1,
ymin=0.1,
ymax=0.9,
hide axis
]
\addplot [
color=ohmred,
only marks,
mark=o,
mark options={solid},
forget plot
]
table[row sep=crcr]{
0.5 0.9\\
0.0833333333333333 0.1\\
0.541666666666667 0.8\\
0.458333333333333 0.7\\
0.875 0.2\\
0.416666666666667 0.6\\
0.75 0.1\\
0.833333333333333 0.1\\
0.916666666666667 0.1\\
0.333333333333333 0.5\\
0.625 0.2\\
0.291666666666667 0.4\\
0.5 0.1\\
0.583333333333333 0.1\\
0.666666666666667 0.1\\
0.166666666666667 0.1\\
0.333333333333333 0.3\\
0.291666666666667 0.2\\
0.416666666666667 0.1\\
0.25 0.1\\
0.333333333333333 0.1\\
};
\addplot [
color=ohmred,
solid,
forget plot
]
table[row sep=crcr]{
0.0833333333333333 0.1\\
0.5 0.9\\
};
\addplot [
color=ohmred,
solid,
forget plot
]
table[row sep=crcr]{
0.541666666666667 0.8\\
0.5 0.9\\
};
\addplot [
color=ohmred,
solid,
forget plot
]
table[row sep=crcr]{
0.458333333333333 0.7\\
0.541666666666667 0.8\\
};
\addplot [
color=ohmred,
solid,
forget plot
]
table[row sep=crcr]{
0.875 0.2\\
0.541666666666667 0.8\\
};
\addplot [
color=ohmred,
solid,
forget plot
]
table[row sep=crcr]{
0.416666666666667 0.6\\
0.458333333333333 0.7\\
};
\addplot [
color=ohmred,
solid,
forget plot
]
table[row sep=crcr]{
0.75 0.1\\
0.458333333333333 0.7\\
};
\addplot [
color=ohmred,
solid,
forget plot
]
table[row sep=crcr]{
0.833333333333333 0.1\\
0.875 0.2\\
};
\addplot [
color=ohmred,
solid,
forget plot
]
table[row sep=crcr]{
0.916666666666667 0.1\\
0.875 0.2\\
};
\addplot [
color=ohmred,
solid,
forget plot
]
table[row sep=crcr]{
0.333333333333333 0.5\\
0.416666666666667 0.6\\
};
\addplot [
color=ohmred,
solid,
forget plot
]
table[row sep=crcr]{
0.625 0.2\\
0.416666666666667 0.6\\
};
\addplot [
color=ohmred,
solid,
forget plot
]
table[row sep=crcr]{
0.291666666666667 0.4\\
0.333333333333333 0.5\\
};
\addplot [
color=ohmred,
solid,
forget plot
]
table[row sep=crcr]{
0.5 0.1\\
0.333333333333333 0.5\\
};
\addplot [
color=ohmred,
solid,
forget plot
]
table[row sep=crcr]{
0.583333333333333 0.1\\
0.625 0.2\\
};
\addplot [
color=ohmred,
solid,
forget plot
]
table[row sep=crcr]{
0.666666666666667 0.1\\
0.625 0.2\\
};
\addplot [
color=ohmred,
solid,
forget plot
]
table[row sep=crcr]{
0.166666666666667 0.1\\
0.291666666666667 0.4\\
};
\addplot [
color=ohmred,
solid,
forget plot
]
table[row sep=crcr]{
0.333333333333333 0.3\\
0.291666666666667 0.4\\
};
\addplot [
color=ohmred,
solid,
forget plot
]
table[row sep=crcr]{
0.291666666666667 0.2\\
0.333333333333333 0.3\\
};
\addplot [
color=ohmred,
solid,
forget plot
]
table[row sep=crcr]{
0.416666666666667 0.1\\
0.333333333333333 0.3\\
};
\addplot [
color=ohmred,
solid,
forget plot
]
table[row sep=crcr]{
0.25 0.1\\
0.291666666666667 0.2\\
};
\addplot [
color=ohmred,
solid,
forget plot
]
table[row sep=crcr]{
0.333333333333333 0.1\\
0.291666666666667 0.2\\
};
\end{axis}
\end{tikzpicture}% \\ Tiefe: 8 \\ Genauigkeit: 92{,}3 \% 
\end{columns}






\end{frame}


\section{Ausblick \& Zusammenfassung}

\subsection{Anwendungsgebiete}
\begin{frame}{Anwendungsgebiete}

\begin{itemize}[<+->]
    \item Medizin
    \item Klassifikation
    \item Regression
    \item Data-Mining
    \item SPAM-Filter
\end{itemize}


\end{frame}

\subsection{Ausblick}
\begin{frame}{Ausblick}
\begin{itemize}[<+->]
    \item CHAID %Chi-Quadrat-Unabhängigkeitstest verwendet
    \item CART
    \item C4.5 und C5.0
    \item Entscheidungswälder
    \item Kombination mit künstlichen neuronalen Netzen
\end{itemize}

\end{frame}


\subsection{Zusammenfassung}
\begin{frame}{Zusammenfassung}
\begin{itemize}[<+->]
    \item Was ist ein Entscheidungsbaum?
    \item Wie konstruiere ich einen Entscheidungsbaum?
    \item Warum haben Atomkraftwerke Kühltürme oder was ist Entropie?
    \item Was ist Informationsgehalt und -gewinn?
    \item Wie funktioniert der ID3-Algorithmus?
    \item Worin liegen die Stärken und Schwächen eines Entscheidungsbaumes?
    \item Was ist Overfitting?
    \item Wo gibt es Optimierungsmöglichkeiten?
    \item Was sind die Einsatzgebiete von Entscheidungsbäumen?
    \item Welche weiteren Algorithmen für Entscheidungsbäume gibt es neben des ID3-Algorithmus' noch?
    
\end{itemize}


\end{frame}


\begin{frame}{Ende}
\centering
Vielen Dank für die Aufmerksamkeit! \\[3em]
\textbf{Fragen?}
\end{frame}





%% All of the following is optional and typically not needed. 
%\appendix
%\section<presentation>*{\appendixname}
%\subsection<presentation>*{For Further Reading}
%
%\begin{frame}[allowframebreaks]
%  \frametitle<presentation>{For Further Reading}
%    
%  \begin{thebibliography}{10}
%    
%  \beamertemplatebookbibitems
%  % Start with overview books.
%
%  \bibitem{Author1990}
%    A.~Author.
%    \newblock {\em Handbook of Everything}.
%    \newblock Some Press, 1990.
% 
%    
%  \beamertemplatearticlebibitems
%  % Followed by interesting articles. Keep the list short. 
%
%  \bibitem{Someone2000}
%    S.~Someone.
%    \newblock On this and that.
%    \newblock {\em Journal of This and That}, 2(1):50--100,
%    2000.
%  \end{thebibliography}
%\end{frame}

\end{document}
